% Les activités agropastorales — Géographie 1re Tech — Cours signé OnBuch+
% Programme officiel MINESEC. Compiler avec : tectonic N04-activites-agropastorales.tex
\newcommand{\DOCMATIERE}{Géographie}
\newcommand{\DOCNIVEAU}{1re Tech}
\newcommand{\DOCMODULE}{Géographie – classe de Première technique}
\newcommand{\DOCLECON}{N04}
\newcommand{\DOCTITRE}{Les activités agropastorales}
% OnBuch+ — préambule « cours signés » ENRICHI (compilé avec tectonic / XeLaTeX)
% Système visuel « Pop » d'OnBuch+ : fond crème (jamais blanc), bordures encre
% épaisses, ombres « dures » décalées, titres Archivo Black en capitales.
% Version MATHÉMATIQUES : graphes (pgfplots/tikz), boîtes pédagogiques
% (définition, propriété, méthode, attention, le savais-tu, exercice résolu,
% exercice, TP, bilan), page de garde et signature OnBuch+ ancrée en bas.
%
% Macros attendues dans main.tex AVANT \input{preamble} :
%   \DOCMATIERE  \DOCNIVEAU  \DOCMODULE  \DOCLECON (numéro)  \DOCTITRE
\documentclass[11pt]{article}

\usepackage{fontspec}
\usepackage[french]{babel}
\usepackage[a4paper,margin=2cm,top=3.4cm,bottom=2.8cm,headheight=1.4cm,footskip=1.3cm]{geometry}
\usepackage{amsmath,amssymb,amsfonts}
\usepackage{mathtools} % \xmapsto, \coloneqq... souvent utilisés par les modèles
\usepackage{cancel} % simplifications barrées (bilans de forces)
% \abs{z} (module, valeur absolue) et \norm{u} (norme), utilisés par les modèles
\providecommand{\abs}{}\let\abs\relax\DeclarePairedDelimiter{\abs}{\lvert}{\rvert}
\providecommand{\norm}{}\let\norm\relax\DeclarePairedDelimiter{\norm}{\lVert}{\rVert}
\usepackage[table]{xcolor} % [table] : fournit aussi \rowcolors (alternance de lignes)
\usepackage{pagecolor}
\usepackage{tikz}
\usetikzlibrary{arrows.meta,positioning,calc,shapes.geometric,shapes.misc,shapes.arrows,decorations.pathmorphing,decorations.pathreplacing,decorations.markings,patterns,shadows,mindmap,trees,fit,backgrounds,matrix,intersections,angles,quotes}
\usepackage{pgfplots}
\usepackage[version=4]{mhchem} % équations nucléaires \ce{^{235}_{92}U}
\usepackage[european,siunitx]{circuitikz} % schémas électriques
\pgfplotsset{compat=1.18}
\usepgfplotslibrary{fillbetween,groupplots} % aires entre courbes (calcul intégral)
\usepackage{enumitem}
\usepackage{fancyhdr}
% [most] charge toutes les bibliothèques tcolorbox (listings, minted, raster,
% documentation...) et gonfle inutilement la mémoire TeX sur un document long
% (~300 boîtes) : on ne charge que ce qui est réellement utilisé.
\usepackage[skins,breakable]{tcolorbox}
\usepackage{graphicx}
\usepackage{colortbl}
\usepackage{booktabs}
\usepackage{tabularx}
\usepackage{multirow}
\usepackage{diagbox} % case d'en-tête diagonale des tableaux à double entrée
% Colonnes extensibles centrée / gauche / droite (C, L, R), souvent utilisées par les modèles
\newcolumntype{C}{>{\centering\arraybackslash}X}
\newcolumntype{L}{>{\raggedright\arraybackslash}X}
\newcolumntype{R}{>{\raggedleft\arraybackslash}X}
\usepackage{array}
\usepackage{multicol}
\usepackage{siunitx}
% parse-numbers=false : accepte aussi \SI{2,5 \times 10^{-3}}{...}
\sisetup{locale=FR,output-decimal-marker={,},group-minimum-digits=4,parse-numbers=false}
\DeclareSIUnit\ohm{\ensuremath{\symup{\Omega}}} % U+2126 absent de la police
\DeclareSIUnit\division{div} % réglages d'oscilloscope (ms/div, V/div)
\DeclareSIUnit\tr{tr} % tours (vitesse de rotation, ex. \SI{1500}{\tr\per\minute})
\DeclareSIUnit\molar{M} % concentration molaire (ex. \SI{3}{\molar}, \SI{1}{\micro\molar})
\DeclareSIUnit\an{an} % année (le modèle confond parfois avec \year, un compteur TeX) — ex. \SI{5}{\kilo\meter\per\an}
\DeclareSIUnit\FCFA{FCFA} % franc CFA (ex. \SI{500}{\FCFA\per\kilo\gram})
\DeclareSIUnit\degreeBrix{\textdegree Brix} % teneur en sucre d'un sirop/jus (réfractométrie)
\DeclareSIUnit\month{mois} % le modèle utilise parfois \month, un compteur TeX, comme unité de durée
\DeclareSIUnit\microliter{\micro\liter} % le modèle écrit parfois \microliter en un seul token
\DeclareSIUnit\million{millions} % dénombrement (ex. \SI{15}{\million\per\mL} spermatozoïdes)
\usepackage{qrcode}
% \degree : commande fréquemment utilisée par les modèles pour un angle ou une
% température en dehors de \SI{}{} (non fournie par les paquets chargés).
\providecommand{\degree}{^{\circ}}
% \qed (fin de démonstration), utilisé par les modèles sans amsthm
\providecommand{\qed}{\hfill$\square$}
% \barre : double barre des valeurs interdites dans un tableau de variations (tkz-tab)
\providecommand{\barre}{\ensuremath{\|}}
% environnement proof (amsthm non chargé) utilisé par les modèles
\ifdefined\proof\else\newenvironment{proof}[1][Démonstration]{\par\noindent\textit{#1.}\ }{\qed\par}\fi
\usepackage{lastpage}
\usepackage{hyperref}
\hypersetup{hidelinks}

% ============================================================
% Palette « Pop » OnBuch+ — mêmes hex que la classe P (plus_ui.dart)
% ============================================================
\definecolor{poporange}{HTML}{F59321}
\definecolor{popdark}{HTML}{E07A0C}
\definecolor{popcream}{HTML}{FFF6E8}
\definecolor{popink}{HTML}{1C1714}
\definecolor{poppurple}{HTML}{7A5AE0}
\definecolor{popgreen}{HTML}{2E9E6A}
\definecolor{poppink}{HTML}{E0457B}
\definecolor{popblue}{HTML}{2F7DE1}
\definecolor{popgold}{HTML}{C98A12}
\definecolor{popmuted}{HTML}{8A7F76}
\definecolor{popline}{HTML}{EADFD2}
\definecolor{popgray}{HTML}{8A7F76}
\colorlet{popgrey}{popgray}
% Alias tolérants (le modèle nomme parfois les couleurs différemment)
\colorlet{popwhite}{white}
\colorlet{popblack}{popink}
\colorlet{popred}{poppink}
\colorlet{popyellow}{popgold}
% Teintes claires (fonds de tableaux/encadrés) — le modèle nomme parfois
% les couleurs avec un suffixe L (ex. popblueL) sans les avoir définies.
\colorlet{poporangeL}{poporange!20}
\colorlet{popdarkL}{popdark!20}
\colorlet{poppurpleL}{poppurple!20}
\colorlet{popgreenL}{popgreen!20}
\colorlet{poppinkL}{poppink!20}
\colorlet{popblueL}{popblue!20}
\colorlet{popgoldL}{popgold!20}
\colorlet{popredL}{popred!20}
\colorlet{popgrayL}{popgray!20}
% Teintes claires pour fonds de tableaux / zones
\colorlet{popblueL}{popblue!10!white}
\colorlet{poporangeL}{poporange!14!white}
\colorlet{popgreenL}{popgreen!12!white}
\colorlet{poppurpleL}{poppurple!12!white}
\colorlet{poppinkL}{poppink!12!white}
\colorlet{popgoldL}{popgold!14!white}

\pagecolor{popcream}

% ============================================================
% Typographie
% ============================================================
\defaultfontfeatures{Ligatures=TeX}
\setmainfont{PlusJakartaSans-Medium.ttf}[Path=../../../fonts/,BoldFont=PlusJakartaSans-Bold.ttf]
% Maths en sans-serif (Fira Math) avec chiffres et lettres droites de Plus
% Jakarta, pour que les formules se fondent dans le texte.
\usepackage[math-style=ISO,bold-style=ISO]{unicode-math}
\setmathfont{FiraMath-Regular.otf}[Path=../../../fonts/]
\setmathfont{PlusJakartaSans-Medium.ttf}[Path=../../../fonts/,range={up/num,up/{Latin,latin},\mathup}]
\usepackage{icomma} % virgule décimale sans espace en mode math
% Caractères mathématiques Unicode absents des polices de texte (Plus Jakarta,
% Archivo Black) : rendus par leur équivalent mathématique, en texte comme en
% formule — sinon ils s'affichent en carré vide (ex. « Arithmétique dans ℤ »).
\usepackage{newunicodechar}
\newunicodechar{ℝ}{\ensuremath{\mathbb{R}}}
\newunicodechar{ℤ}{\ensuremath{\mathbb{Z}}}
\newunicodechar{ℂ}{\ensuremath{\mathbb{C}}}
\newunicodechar{ℕ}{\ensuremath{\mathbb{N}}}
\newunicodechar{ℚ}{\ensuremath{\mathbb{Q}}}
\newunicodechar{𝔻}{\ensuremath{\mathbb{D}}}
\newunicodechar{α}{\ensuremath{\alpha}}
\newunicodechar{β}{\ensuremath{\beta}}
\newunicodechar{γ}{\ensuremath{\gamma}}
\newunicodechar{δ}{\ensuremath{\delta}}
\newunicodechar{ε}{\ensuremath{\varepsilon}}
\newunicodechar{θ}{\ensuremath{\theta}}
\newunicodechar{λ}{\ensuremath{\lambda}}
\newunicodechar{μ}{\ensuremath{\mu}}
\newunicodechar{σ}{\ensuremath{\sigma}}
\newunicodechar{τ}{\ensuremath{\tau}}
\newunicodechar{φ}{\ensuremath{\varphi}}
\newunicodechar{ψ}{\ensuremath{\psi}}
\newunicodechar{ω}{\ensuremath{\omega}}
\newunicodechar{Σ}{\ensuremath{\Sigma}}
% majuscules produites par \MakeUppercase dans les titres (α-aminés -> Α)
\newunicodechar{Α}{\ensuremath{\alpha}}
\newunicodechar{Β}{\ensuremath{\beta}}
\newunicodechar{Γ}{\ensuremath{\Gamma}}
\newunicodechar{Δ}{\ensuremath{\Delta}}
\newunicodechar{Ε}{\ensuremath{\varepsilon}}
\newunicodechar{Θ}{\ensuremath{\Theta}}
\newunicodechar{Λ}{\ensuremath{\Lambda}}
\newunicodechar{Μ}{\ensuremath{\mu}}
\newunicodechar{Τ}{\ensuremath{\tau}}
\newunicodechar{Φ}{\ensuremath{\Phi}}
\newunicodechar{Ψ}{\ensuremath{\Psi}}
\newunicodechar{Ω}{\ensuremath{\Omega}}
\newunicodechar{↦}{\ensuremath{\mapsto}}
\newunicodechar{∈}{\ensuremath{\in}}
\newunicodechar{∉}{\ensuremath{\notin}}
\newunicodechar{∧}{\ensuremath{\wedge}}
\newunicodechar{⇔}{\ensuremath{\Leftrightarrow}}
\newunicodechar{⟺}{\ensuremath{\Longleftrightarrow}}
\newunicodechar{⇒}{\ensuremath{\Rightarrow}}
\newunicodechar{⟹}{\ensuremath{\Longrightarrow}}
\newunicodechar{∘}{\ensuremath{\circ}}
\newunicodechar{∪}{\ensuremath{\cup}}
\newunicodechar{∩}{\ensuremath{\cap}}
\newunicodechar{≡}{\ensuremath{\equiv}}
\newunicodechar{⊂}{\ensuremath{\subset}}
\newunicodechar{⊥}{\ensuremath{\perp}}
\newunicodechar{∃}{\ensuremath{\exists}}
\newunicodechar{∀}{\ensuremath{\forall}}
\newunicodechar{∛}{\ensuremath{\sqrt[3]{\,}}}
\newunicodechar{∜}{\ensuremath{\sqrt[4]{\,}}}
\newunicodechar{⃗}{\ensuremath{{}^{\to}}}
\newunicodechar{ⁿ}{\ifmmode^{n}\else\textsuperscript{\ensuremath{n}}\fi}
\newunicodechar{ˣ}{\ifmmode^{x}\else\textsuperscript{\ensuremath{x}}\fi}
\newunicodechar{ᵇ}{\ifmmode^{b}\else\textsuperscript{\ensuremath{b}}\fi}
\newunicodechar{ʳ}{\ifmmode^{r}\else\textsuperscript{\ensuremath{r}}\fi}
\newunicodechar{ᵘ}{\ifmmode^{u}\else\textsuperscript{\ensuremath{u}}\fi}
\newunicodechar{ʸ}{\ifmmode^{y}\else\textsuperscript{\ensuremath{y}}\fi}
\newunicodechar{ᵃ}{\ifmmode^{a}\else\textsuperscript{\ensuremath{a}}\fi}
\newunicodechar{ᵉ}{\ifmmode^{e}\else\textsuperscript{\ensuremath{e}}\fi}
\newunicodechar{ᵏ}{\ifmmode^{k}\else\textsuperscript{\ensuremath{k}}\fi}
\newunicodechar{ᵖ}{\ifmmode^{p}\else\textsuperscript{\ensuremath{p}}\fi}
\newunicodechar{ᴺ}{\ifmmode^{N}\else\textsuperscript{\ensuremath{N}}\fi}
\newunicodechar{ᶜ}{\ifmmode^{c}\else\textsuperscript{\ensuremath{c}}\fi}
\newunicodechar{ʰ}{\ifmmode^{h}\else\textsuperscript{\ensuremath{h}}\fi}
\newunicodechar{ᵠ}{\ifmmode^{q}\else\textsuperscript{\ensuremath{q}}\fi}
\newunicodechar{ₙ}{\ifmmode_{n}\else\textsubscript{\ensuremath{n}}\fi}
\newunicodechar{ₐ}{\ifmmode_{a}\else\textsubscript{\ensuremath{a}}\fi}
\newunicodechar{ᵢ}{\ifmmode_{i}\else\textsubscript{\ensuremath{i}}\fi}
\newunicodechar{ᵣ}{\ifmmode_{r}\else\textsubscript{\ensuremath{r}}\fi}
\newunicodechar{ₖ}{\ifmmode_{k}\else\textsubscript{\ensuremath{k}}\fi}
\newunicodechar{ᵦ}{\ifmmode_{\beta}\else\textsubscript{\ensuremath{\beta}}\fi}
\newunicodechar{ₓ}{\ifmmode_{x}\else\textsubscript{\ensuremath{x}}\fi}
\newfontfamily\popdisplay{ArchivoBlack-Regular.ttf}[Path=../../../fonts/]
\newfontfamily\poplogo{SpaceGrotesk-Bold.ttf}[Path=../../../fonts/]
\newfontfamily\popsemi{PlusJakartaSans-SemiBold.ttf}[Path=../../../fonts/]
\newfontfamily\popxbold{PlusJakartaSans-ExtraBold.ttf}[Path=../../../fonts/]
\newfontfamily\popxboldls{PlusJakartaSans-ExtraBold.ttf}[Path=../../../fonts/,LetterSpace=3.0]

\setlist[enumerate]{leftmargin=1.7em,itemsep=5pt,topsep=4pt}
\setlist[itemize]{leftmargin=1.7em,itemsep=4pt,topsep=4pt,label={\color{poporange}\textbullet}}
\setlength{\parindent}{0pt}
\setlength{\parskip}{5pt}
\renewcommand{\arraystretch}{1.25}

% pgfplots : style Pop par défaut
\pgfplotsset{
  every axis/.append style={axis line style={popink,line width=1pt},tick style={popink},
    grid style={popline,line width=0.5pt},label style={font=\small},tick label style={font=\footnotesize}},
  popcurve/.style={poporange,line width=1.8pt,no markers,smooth},
}
% Même style disponible dans un \draw TikZ simple (hors pgfplots)
\tikzset{popcurve/.style={poporange,line width=1.8pt}}
\tikzset{popgrid/.style={popline,very thin}}
\colorlet{popcurve}{poporange} % le nom du style est parfois utilisé comme couleur

% ============================================================
% Pill / squiggle
% ============================================================
\newcommand{\poppill}[3]{%
  \tikz[baseline=-0.55ex]\node[draw=popink,line width=1.2pt,rounded corners=2.6pt,%
    fill=#2,inner xsep=6.5pt,inner ysep=3pt]{%
    \popxboldls\fontsize{7}{7}\selectfont\color{#3}\MakeUppercase{#1}};%
}
\newcommand{\popsquiggle}[1]{%
  \tikz[baseline=-0.4ex]{
    \draw[#1,line width=2.6pt,line cap=round]
      plot [smooth] coordinates {(0,0.09) (0.34,0.01) (0.68,0.21) (1.02,0.01) (1.36,0.10)};
  }%
}
% Mot-clé surligné (marqueur orange sous le texte)
\newcommand{\cle}[1]{{\popxbold\color{popdark}#1}}

% ============================================================
% En-tête / pied
% ============================================================
\pagestyle{fancy}
\fancyhf{}
\renewcommand{\headrulewidth}{0pt}
\renewcommand{\footrulewidth}{0pt}
\lhead{\raisebox{-0.15ex}{\poplogo\fontsize{13}{13}\selectfont\color{popink}OnBuch\color{poporange}+}}
\rhead{\poppill{\DOCMATIERE\ \textbullet\ \DOCNIVEAU\ \textbullet\ Leçon \DOCLECON}{white}{popink}}
\lfoot{\popxbold\fontsize{7.6}{7.6}\selectfont\color{popmuted}\MakeUppercase{Cours signé OnBuch+}\ \textbullet\ {\popsemi\DOCTITRE}}
\rfoot{\tikz[baseline=-0.6ex]\node[rounded corners=8.5pt,fill=popink,minimum height=17pt,minimum width=17pt,inner xsep=5pt,inner sep=0]{\popxbold\fontsize{8}{8}\selectfont\color{popcream}\thepage\,/\,\pageref*{LastPage}};}

% ============================================================
% Titres
% ============================================================
\newcommand{\chaptitre}[2]{%
  \begin{tcolorbox}[enhanced,colback=poporange,colframe=popink,boxrule=2.6pt,arc=11pt,
      drop shadow=popink,left=15pt,right=15pt,top=12pt,bottom=11pt,before skip=3pt,after skip=18pt]
    \poppill{#2}{popink}{popcream}\par\vspace{9pt}
    {\raggedright\hyphenpenalty=10000\exhyphenpenalty=10000\popdisplay\fontsize{22}{24}\selectfont\color{popcream}\MakeUppercase{#1}\par}
  \end{tcolorbox}
}
% Section numérotée : pastille orange avec le numéro + titre Archivo
\newcounter{popsec}
\newcommand{\coursec}[1]{%
  \stepcounter{popsec}\setcounter{popsub}{0}%
  \par\vspace{16pt}\noindent
  \begin{minipage}{\linewidth}
  \tikz[baseline=-0.7ex]\node[draw=popink,line width=1.6pt,fill=poporange,rounded corners=4pt,minimum size=22pt,inner sep=0]{\popdisplay\fontsize{12}{12}\selectfont\color{popcream}\thepopsec};%
  \hspace{8pt}{\popdisplay\fontsize{15}{17}\selectfont\color{popink}\MakeUppercase{#1}}\par
  \vspace{4pt}\noindent{\color{poporange}\rule{36pt}{3.6pt}}
  \end{minipage}\par\nopagebreak\vspace{8pt}
}
\newcounter{popsub}
\newcommand{\courssub}[1]{%
  \stepcounter{popsub}%
  \par\vspace{9pt}\noindent{\popxbold\fontsize{11.5}{13}\selectfont\color{popink}{\color{poporange}\thepopsec.\thepopsub}\hspace{6pt}#1}\par\nopagebreak\vspace{4pt}
}

% ============================================================
% Boîtes pédagogiques — même recette : fond blanc, bordure épaisse colorée,
% ombre dure encre, badge plein. Titre optionnel [#1].
% ============================================================
\tcbset{popbase/.style={breakable,enhanced,colback=white,boxrule=2.3pt,arc=9pt,
  shadow={3pt}{-3pt}{0pt}{popink},left=14pt,right=14pt,top=11pt,bottom=11pt,before skip=11pt,after skip=15pt,
  fontupper=\popsemi\fontsize{10.6}{14.5}\selectfont\color{popink},
  fontlower=\popsemi\fontsize{10.6}{14.5}\selectfont\color{popink}}}
% Coche dessinée (le glyphe ✓ n'existe pas dans les polices du document)
\newcommand{\popcheck}[1]{\tikz[baseline=-0.5ex]\draw[#1,line width=1.6pt,line cap=round,line join=round](-0.13,0.0)--(-0.04,-0.09)--(0.13,0.1);}
\providecommand{\coche}{\popcheck{popgreen}}
\providecommand{\resultat}[1]{\par\smallskip\noindent\fcolorbox{popgreen}{popgreenL}{\parbox{\dimexpr\linewidth-2\fboxsep-2\fboxrule\relax}{\textbf{Résultat.} #1}}\par\smallskip}
\newcommand{\popboxhead}[3]{\poppill{#1}{#2}{white}\ifx\relax#3\relax\else\hspace{6pt}{\popxbold\fontsize{10.6}{12}\selectfont\color{popink}#3}\fi\par\vspace{7pt}}

\newtcolorbox{aretenir}[1][]{popbase,colframe=popgreen,before upper={\popboxhead{À retenir}{popgreen}{#1}}}
\newtcolorbox{exemplebox}[1][]{popbase,colframe=poporange,before upper={\popboxhead{Exemple}{poporange}{#1}}}
\newtcolorbox{definition}[1][]{popbase,colframe=popblue,before upper={\popboxhead{Définition}{popblue}{#1}}}
\newtcolorbox{propriete}[1][]{popbase,colframe=poppurple,before upper={\popboxhead{Propriété}{poppurple}{#1}}}
\newtcolorbox{methode}[1][]{popbase,colframe=popgold,before upper={\popboxhead{Méthode}{popgold}{#1}}}
\newtcolorbox{attention}[1][]{popbase,colframe=poppink,before upper={\popboxhead{Attention}{poppink}{#1}}}
\newtcolorbox{experience}[1][]{popbase,colframe=popblue,colback=popblueL,before upper={\popboxhead{Expérience}{popblue}{#1}}}
% « Le savais-tu ? » : carte violette pleine (contexte, culture, Cameroun)
\newtcolorbox{savaistu}[1][]{popbase,colback=poppurple,colframe=popink,
  fontupper=\popsemi\fontsize{10.4}{14.2}\selectfont\color{white},
  before upper={\poppill{Le savais-tu\,?}{white}{poppurple}\ifx\relax#1\relax\else\hspace{6pt}{\popxbold\color{white}#1}\fi\par\vspace{7pt}}}
% Exercice résolu : énoncé en haut, \tcblower puis la solution sur fond crème
\newcounter{popexo}\newcounter{popexr}
\newtcolorbox{exoresolu}[1][]{popbase,colframe=poporange,colbacklower=popcream,
  segmentation style={popink,line width=1.2pt,dashed},
  before upper={\stepcounter{popexr}\popboxhead{Exercice résolu \thepopexr}{poporange}{#1}},
  before lower={\poppill{Solution}{popink}{popcream}\par\vspace{6pt}}}
% Exercice (non résolu en place) — la correction est dans la partie Corrigés
\newtcolorbox{exercice}[1][]{popbase,colframe=popink,
  before upper={\stepcounter{popexo}\popboxhead{Exercice \thepopexo}{popink}{#1}}}
% Corrigé
\newtcolorbox{corrige}[1][]{popbase,colframe=popgreen,colback=popgreenL,
  before upper={\popboxhead{Corrigé}{popgreen}{#1}}}
% Objectifs / prérequis / bilan
\newtcolorbox{objectifs}{popbase,colframe=popink,colback=poporangeL,
  before upper={\popboxhead{Objectifs de la leçon}{popink}{}}}
\newtcolorbox{prerequis}{popbase,colframe=popink,colback=popblueL,
  before upper={\popboxhead{Prérequis}{popblue}{}}}
\newtcolorbox{bilan}{popbase,colframe=popink,colback=white,boxrule=2.8pt,
  before upper={\popboxhead{Fiche bilan}{poporange}{L'essentiel à connaître pour l'examen}}}
% Figure encadrée avec légende
\newtcolorbox{popfigure}[1][]{enhanced,breakable=false,colback=white,colframe=popline,boxrule=1.4pt,arc=8pt,
  left=10pt,right=10pt,top=10pt,bottom=8pt,before skip=10pt,after skip=12pt,halign=center,
  lower separated=false,halign lower=center,
  fontlower=\popsemi\fontsize{9}{11.5}\selectfont\color{popmuted},
  #1}
\newcommand{\legende}[1]{\par\vspace{6pt}{\popsemi\fontsize{9}{11.5}\selectfont\color{popmuted}#1}}

% ============================================================
% Signature OnBuch+
% ============================================================
\newcommand{\poplogoinline}[1]{{\poplogo\fontsize{#1}{#1}\selectfont\color{popink}OnBuch\color{poporange}+}}
\newcommand{\signatureonbuch}{%
  \begin{tcolorbox}[enhanced,colback=popink,colframe=popink,boxrule=2.2pt,arc=10pt,
      drop shadow=poporange,left=14pt,right=14pt,top=10pt,bottom=10pt,before skip=0pt,after skip=0pt]
    \tikz[baseline=-0.7ex]\node[circle,fill=popgreen,draw=popcream,line width=1.3pt,minimum size=22pt,inner sep=0]{\popcheck{white}};%
    \hspace{9pt}\begin{minipage}[c]{0.62\linewidth}
      {\popxbold\fontsize{12}{13}\selectfont\color{popcream}Signé\ \textbullet\ {\poplogo OnBuch\color{poporange}+}}\par\vspace{2pt}
      {\raggedright\popsemi\fontsize{8}{10.5}\selectfont\color{popline}\MakeUppercase{Équipe pédagogique OnBuch+}\\\MakeUppercase{Conforme au programme officiel MINESEC}\par}
    \end{minipage}\hfill
    {\poplogo\fontsize{22}{22}\selectfont\color{popcream}OnBuch\color{poporange}+}
  \end{tcolorbox}%
}

% ============================================================
% Page de garde
% #1 titre · #2 matière · #3 niveau · #4 module · #5 n° leçon · #6 durée ·
% #7 description / objectifs (texte ou itemize)
% ============================================================
\newcommand{\pagegarde}[7]{%
  \thispagestyle{empty}
  \newgeometry{a4paper,margin=1.7cm,top=1.6cm,bottom=1.4cm}%
  \begin{tikzpicture}[remember picture,overlay]
    \fill[poporange!18] ([xshift=-3.2cm,yshift=-3.6cm]current page.north east) circle (4.6cm);
    \fill[poppurple!14] ([xshift=2.4cm,yshift=7.6cm]current page.south west) circle (3.8cm);
  \end{tikzpicture}%
  {\poplogo\fontsize{30}{30}\selectfont\color{popink}OnBuch\color{poporange}+}\hfill
  \raisebox{0.6ex}{\poppill{Cours signé \textbullet\ Premium}{popink}{popcream}}\par
  \vspace{1pt}\popsquiggle{poppurple}\par
  \vspace{8pt}
  \begin{tcolorbox}[enhanced,colback=poporange,colframe=popink,boxrule=2.8pt,arc=13pt,
      drop shadow=popink,left=18pt,right=18pt,top=12pt,bottom=13pt,after skip=10pt]
    \poppill{#2 \textbullet\ #3}{popink}{popcream}\hspace{5pt}\poppill{#4}{white}{popink}\par\vspace{8pt}
    {\popdisplay\fontsize{14}{14}\selectfont\color{popink}LEÇON #5}\par\vspace{4pt}
    {\raggedright\hyphenpenalty=10000\exhyphenpenalty=10000\popdisplay\fontsize{25}{27}\selectfont\color{popcream}\MakeUppercase{#1}\par}
    \vspace{8pt}
    {\popxbold\fontsize{9.5}{9.5}\selectfont\color{popink}\MakeUppercase{Durée conseillée : #6}}
  \end{tcolorbox}
  \begin{tcolorbox}[enhanced,colback=white,colframe=popink,boxrule=2.2pt,arc=10pt,
      drop shadow=popink,left=16pt,right=16pt,top=10pt,bottom=10pt,after skip=8pt]
    \poppill{Dans cette leçon}{poppurple}{white}\par\vspace{5pt}
    \popsemi\fontsize{9.3}{12.6}\selectfont\color{popink}#7
  \end{tcolorbox}
  \noindent\begin{minipage}[t]{0.487\linewidth}
    \begin{tcolorbox}[enhanced,colback=popgreen,colframe=popink,boxrule=2.2pt,arc=10pt,
        drop shadow=popink,left=11pt,right=11pt,top=10pt,bottom=10pt,height=3.1cm,valign=center]
      \tcbox[colback=white,colframe=popink,boxrule=1.3pt,arc=5pt,boxsep=0pt,nobeforeafter,
        left=3pt,right=3pt,top=3pt,bottom=3pt,valign=center]{\qrcode[height=1.9cm]{https://play.google.com/store/apps/details?id=com.luvvix.onbuch&hl=fr}}%
      \hspace{8pt}\begin{minipage}[c]{0.5\linewidth}
        {\popdisplay\fontsize{11}{12.5}\selectfont\color{white}\MakeUppercase{Télécharge OnBuch}}\par\vspace{4pt}
        {\raggedright\popsemi\fontsize{8.4}{11}\selectfont\color{white}Gratuit \textbullet\ Google Play\\Tuteur IA, annales, concours, résultats\par}
      \end{minipage}
    \end{tcolorbox}
  \end{minipage}\hfill
  \begin{minipage}[t]{0.487\linewidth}
    \begin{tcolorbox}[enhanced,colback=poppurple,colframe=popink,boxrule=2.2pt,arc=10pt,
        drop shadow=popink,left=11pt,right=11pt,top=10pt,bottom=10pt,height=3.1cm,valign=center]
      \tikz[baseline=-0.7ex]\node[circle,fill=white,draw=popink,line width=1.3pt,minimum size=1.6cm,inner sep=0]{\popdisplay\fontsize{24}{24}\selectfont\color{poppurple}+};%
      \hspace{8pt}\begin{minipage}[c]{0.6\linewidth}
        {\popdisplay\fontsize{11}{12.5}\selectfont\color{white}\MakeUppercase{Passe à OnBuch+}}\par\vspace{4pt}
        {\raggedright\popsemi\fontsize{8.4}{11}\selectfont\color{white}Tous les cours signés, Tuteur IA illimité, fiches PDF — onglet OnBuch+ de l'app\par}
      \end{minipage}
    \end{tcolorbox}
  \end{minipage}
  \vspace{8pt}
  \popcontenu
  \vfill
  \signatureonbuch
  \restoregeometry
  \newpage
}

% Rangée « Ce cours contient » (page de garde)
\newcommand{\poptile}[3]{% #1 couleur #2 gros texte #3 légende
  \begin{tcolorbox}[enhanced,colback=#1,colframe=popink,boxrule=2pt,arc=9pt,shadow={2.5pt}{-2.5pt}{0pt}{popink},
      left=6pt,right=6pt,top=9pt,bottom=9pt,halign=center,valign=center,height=1.75cm,width=\linewidth,nobeforeafter]
    {\popdisplay\fontsize{12.5}{13}\selectfont\color{white}\MakeUppercase{#2}}\par\vspace{3pt}
    {\centering\popsemi\fontsize{7.8}{9.5}\selectfont\color{white}#3\par}
  \end{tcolorbox}}
\newcommand{\popcontenu}{%
  \par\noindent{\popxboldls\fontsize{7.5}{7.5}\selectfont\color{popmuted}CE COURS SIGNÉ CONTIENT}\par\vspace{7pt}
  \noindent\begin{minipage}[t]{0.235\linewidth}\poptile{popblue}{Cours}{complet, illustré, pas à pas}\end{minipage}\hfill
  \begin{minipage}[t]{0.235\linewidth}\poptile{poporange}{Méthodes}{et exercices résolus}\end{minipage}\hfill
  \begin{minipage}[t]{0.235\linewidth}\poptile{poppink}{Exercices}{type examen + corrigés}\end{minipage}\hfill
  \begin{minipage}[t]{0.235\linewidth}\poptile{popgreen}{Fiche bilan}{l'essentiel à retenir}\end{minipage}\par}

% Sommaire
\newtcolorbox{sommaire}{enhanced,colback=white,colframe=popink,boxrule=2.2pt,arc=10pt,
  drop shadow=popink,left=18pt,right=18pt,top=13pt,bottom=13pt,before skip=6pt,after skip=20pt,
  before upper={\poppill{Sommaire}{poppurple}{white}\par\vspace{9pt}\popsemi\fontsize{10.6}{14.5}\selectfont\color{popink}}}

% Signature de fin de cours — poussée tout en bas de la dernière page
\newcommand{\findecours}{%
  \par\vfill\nopagebreak
  \begin{center}\popsquiggle{poporange}\end{center}\vspace{4pt}
  \signatureonbuch
}
\AtBeginDocument{\def\llbracket{\mathopen{[\mkern-3.5mu[}}\def\rrbracket{\mathclose{]\mkern-3.5mu]}}} % intervalles d'entiers (glyphe ⟦ absent de la police)
% Glyphes absents de Fira Math : redéfinis à partir de glyphes disponibles
\AtBeginDocument{%
  \def\setminus{\mathbin{\backslash}}\let\smallsetminus\setminus
  \def\vdots{\mathord{\vbox{\baselineskip3.2pt\lineskiplimit0pt\kern4pt\hbox{.}\hbox{.}\hbox{.}}}}%
  \def\ddots{\mathinner{\mkern1mu\raise6.5pt\hbox{.}\mkern2mu\raise3.6pt\hbox{.}\mkern2mu\raise0.7pt\hbox{.}\mkern1mu}}%
  \def\triangle{\mathord{\scalebox{0.9}{$\symup{\Delta}$}}}%
  \def\nabla{\mathord{\raisebox{\depth}{\scalebox{1}[-1]{$\symup{\Delta}$}}}}%
  \def\checkmark{\popcheck{popdark}}%
  \def\star{\mathbin{\ast}}\def\bigstar{\mathord{\ast}}%
  \def\diamond{\mathbin{\scalebox{0.6}{$\Diamond$}}}%
  \def\bigcup{\mathop{\mathchoice{\raisebox{-0.3em}{\scalebox{1.7}{$\cup$}}}{\scalebox{1.25}{$\cup$}}{\cup}{\cup}}\displaylimits}%
  \def\bigcap{\mathop{\mathchoice{\raisebox{-0.3em}{\scalebox{1.7}{$\cap$}}}{\scalebox{1.25}{$\cap$}}{\cap}{\cap}}\displaylimits}%
  \def\lesseqqgtr{\mathrel{\lessgtr}}\def\bigoplus{\mathop{\oplus}\displaylimits}%
}
\providecommand{\ssi}{si et seulement si} % abréviation fréquente
\AtBeginDocument{\providecommand{\wideparen}{\overparen}} % arc de cercle

\begin{document}

\pagegarde{Les activités agropastorales}{Géographie}{1re Tech}{Géographie – classe de Première technique}{N04}{6 séances (fiche de progression harmonisée)}{Cette leçon étudie les activités agropastorales et piscicoles au Cameroun : les types d'agriculture, d'élevage et de pêche, leurs conditions de développement et leurs répartitions spatiales. Elle s'achève par un travail dirigé sur les problèmes du secteur de la production agropastorale et piscicole et les solutions envisageables.
\vspace{4pt}
\begin{multicols}{2}\small
\begin{itemize}
\item À la fin de cette leçon, je sais définir l'agriculture, l'élevage et la pêche et distinguer leurs principaux types au Cameroun
\item À la fin de cette leçon, je sais décrire les caractéristiques de l'agriculture vivrière et de l'agriculture de rente (plantations) à partir de documents
\item À la fin de cette leçon, je sais localiser sur une carte les grandes zones de production agricole, d'élevage et de pêche du Cameroun
\item À la fin de cette leçon, je sais expliquer les conditions naturelles et humaines du développement des activités agropastorales
\item À la fin de cette leçon, je sais construire un croquis, un diagramme ou une carte simple à partir de données chiffrées sur la production
\item À la fin de cette leçon, je sais identifier et hiérarchiser les problèmes du secteur agropastoral et piscicole camerounais
\end{itemize}\end{multicols}}

\begin{sommaire}
\begin{multicols}{2}
\begin{enumerate}
\item Généralités sur les activités agropastorales et piscicoles
\item L'agriculture au Cameroun
\item L'élevage au Cameroun
\item La pêche au Cameroun
\item TD : Problèmes et perspectives du secteur agropastoral et piscicole
\item Activité d'intégration
\item Méthodes et exercices résolus
\item Exercices
\item Corrigés des exercices
\item Fiche bilan
\end{enumerate}
\end{multicols}
\end{sommaire}

\begin{objectifs}
\begin{itemize}
\item À la fin de cette leçon, je sais définir l'agriculture, l'élevage et la pêche et distinguer leurs principaux types au Cameroun
\item À la fin de cette leçon, je sais décrire les caractéristiques de l'agriculture vivrière et de l'agriculture de rente (plantations) à partir de documents
\item À la fin de cette leçon, je sais localiser sur une carte les grandes zones de production agricole, d'élevage et de pêche du Cameroun
\item À la fin de cette leçon, je sais expliquer les conditions naturelles et humaines du développement des activités agropastorales
\item À la fin de cette leçon, je sais construire un croquis, un diagramme ou une carte simple à partir de données chiffrées sur la production
\item À la fin de cette leçon, je sais identifier et hiérarchiser les problèmes du secteur agropastoral et piscicole camerounais
\item À la fin de cette leçon, je sais proposer et justifier des solutions aux problèmes de la production agropastorale et piscicole
\item À la fin de cette leçon, je sais rédiger et justifier une démarche, une réponse ou une conclusion en mobilisant des documents variés
\end{itemize}
\end{objectifs}

\begin{prerequis}
\begin{itemize}
\item Les grandes zones climatiques et végétales du Cameroun (climat équatorial, tropical, soudanien)
\item Le relief et les sols du Cameroun (notion de 4e : milieux naturels)
\item La répartition de la population camerounaise et les activités rurales
\item La notion de secteur primaire et son rôle dans l'économie
\item La lecture simple d'une carte et d'un tableau statistique
\end{itemize}
\end{prerequis}

\newpage

\coursec{Généralités sur les activités agropastorales et piscicoles}

\textbf{Situation d'accroche.} Au marché Mokolo de Yaoundé, Amina constate que le prix du maïs a doublé en un an, tandis que son oncle éleveur de l'Adamaoua se plaint de la baisse de ses revenus malgré un troupeau en bonne santé. Pendant ce temps, à Kribi, les pêcheurs artisans rentrent avec des prises de plus en plus maigres. Pourtant, le Cameroun est souvent présenté comme « l'Afrique en miniature » aux potentialités agropastorales exceptionnelles. Comment expliquer ce paradoxe entre un potentiel remarquable et des difficultés persistantes ? Quelles sont les formes d'agriculture, d'élevage et de pêche pratiquées au Cameroun, et quels problèmes freinent ce secteur vital ?

\textbf{Problématique.} Avant de répondre à ces questions, il faut d'abord comprendre ce que recouvrent les activités agropastorales et piscicoles, quelle place elles occupent dans l'économie camerounaise, et sur quelles conditions naturelles et humaines elles reposent. C'est l'objet de cette première partie.

\courssub{Définitions et notions fondamentales}

\courssub{L'agriculture}

Dans la vie quotidienne, quand on dit qu'on « fait l'agriculture », on pense souvent à tout ce qui se passe à la campagne : les champs, les vaches, les poulets. En géographie, il faut être plus précis. L'intuition est simple : l'agriculture, c'est d'abord le travail de la \emph{terre} pour faire pousser des \emph{plantes}.

\begin{definition}[L'agriculture]
L'\cle{agriculture} est l'ensemble des activités de production végétale (les \cle{cultures}) destinées à l'alimentation humaine et animale ou à l'industrie. Elle comprend la préparation du sol, les semis, l'entretien des plantes, la récolte et parfois la première transformation des produits.
\end{definition}

On distingue deux grandes catégories de cultures au Cameroun :
\begin{itemize}
\item les \cle{cultures vivrières}, destinées à nourrir les hommes et les animaux : maïs, manioc, igname, plantain, mil, sorgho, riz, arachide, niébé ;
\item les \cle{cultures de rente} (ou cultures d'exportation), destinées à la vente, souvent à l'étranger, ou à l'industrie : cacao, café, coton, banane, hévéa (caoutchouc), palmier à huile, canne à sucre, tabac.
\end{itemize}

\courssub{L'élevage et la pêche}

\begin{definition}[L'élevage]
L'\cle{élevage} est l'ensemble des activités de production animale : il consiste à nourrir, protéger et faire se reproduire des animaux domestiques (bovins, ovins, caprins, porcins, volailles) pour obtenir de la viande, du lait, des œufs, des peaux ou de la force de travail.
\end{definition}

\begin{definition}[La pêche]
La \cle{pêche} est l'activité qui consiste à capturer des poissons et d'autres produits aquatiques (crevettes, crabes) dans les milieux naturels (mer, fleuves, lacs) ou dans des plans d'eau aménagés (pisciculture).
\end{definition}

\begin{definition}[Le secteur agropastoral et piscicole]
Le \cle{secteur agropastoral et piscicole} est l'ensemble regroupant l'\cle{agriculture} (production végétale), l'\cle{élevage} (production animale, d'où le terme \emph{pastoral}, du latin \emph{pastor}, le berger) et la \cle{pêche} (production halieutique, d'où le terme \emph{piscicole}, du latin \emph{piscis}, le poisson). Il constitue l'essentiel du \cle{secteur primaire} de l'économie camerounaise.
\end{definition}

\begin{attention}[Erreur fréquente : agriculture et élevage]
Beaucoup d'élèves confondent « agriculture » (les cultures, donc les \emph{végétaux}) et « élevage » (les \emph{animaux}). Ce sont deux activités distinctes, mais complémentaires : ensemble, elles forment le secteur \emph{agropastoral}. À l'examen, écrire « l'agriculture des bovins » est une erreur : on élève des bovins, on cultive du maïs.
\end{attention}

\courssub{La place du secteur primaire dans l'économie camerounaise}

\courssub{Un secteur qui nourrit et fait vivre la majorité des Camerounais}

Pour mesurer l'importance d'un secteur, le géographe utilise deux indicateurs : sa part dans le \cle{PIB} (Produit Intérieur Brut, c'est-à-dire la richesse totale produite par le pays en un an) et la part de la \cle{population active} qu'il emploie. Or ces deux indicateurs racontent deux histoires différentes au Cameroun.

\begin{exemplebox}[Un exemple chiffré]
Au Cameroun, l'agriculture occupe environ \cle{6 actifs sur 10} (soit environ 60 \% de la population active), alors que le secteur primaire ne fournit qu'environ 15 à 17 \% du PIB. Cela signifie qu'une majorité de Camerounais travaille beaucoup pour une part relativement faible de la richesse nationale : c'est le signe d'une agriculture encore peu mécanisée et à faible rendement. L'agriculture fournit aussi le cacao, le café et le coton, qui figurent parmi les principaux produits d'exportation du pays, aux côtés du pétrole et du bois.
\end{exemplebox}

\begin{popfigure}
\begin{center}
\begin{tikzpicture}
\fill[popblue] (0,0) -- (0:2.2) arc (0:61.2:2.2) -- cycle;
\fill[poporange] (0,0) -- (61.2:2.2) arc (61.2:158.4:2.2) -- cycle;
\fill[popgreen] (0,0) -- (158.4:2.2) arc (158.4:360:2.2) -- cycle;
\draw[white,thick] (0,0) circle (2.2);
\node[white] at (30:1.4) {17\,\%};
\node[white] at (110:1.4) {27\,\%};
\node[white] at (260:1.4) {56\,\%};
\node[anchor=west] at (2.6,1) {\textcolor{popblue}{$\blacksquare$} Secteur primaire};
\node[anchor=west] at (2.6,0.3) {\textcolor{poporange}{$\blacksquare$} Secteur secondaire};
\node[anchor=west] at (2.6,-0.4) {\textcolor{popgreen}{$\blacksquare$} Secteur tertiaire};
\end{tikzpicture}
\end{center}
\legende{Diagramme circulaire : part approximative des trois secteurs d'activité dans le PIB du Cameroun (ordres de grandeur récents). Le secteur primaire (agriculture, élevage, pêche, forêts) représente environ 17 \% du PIB, le secondaire (industries, BTP) environ 27 \% et le tertiaire (services, commerce) environ 56 \%. À comparer avec les 60 \% d'actifs occupés par l'agriculture : le contraste est frappant.}
\end{popfigure}

\textbf{Comment lire ce paradoxe ?} La démarche du géographe consiste à croiser les deux données : 60 \% des actifs pour 17 \% du PIB. Le calcul du rapport est éloquent : un actif agricole produit en moyenne une richesse par tête bien inférieure à celle d'un actif du tertiaire. Ce déséquilibre explique à la fois la pauvreté rurale et l'exode des jeunes vers les villes — et il est au cœur de la situation d'accroche : l'oncle d'Amina travaille dur, mais ses revenus baissent.

\begin{aretenir}[L'essentiel sur la place du secteur]
\begin{itemize}
\item Le secteur primaire fournit environ 15 à 17 \% du PIB camerounais.
\item L'agriculture emploie environ 60 \% de la population active.
\item Cacao, café et coton comptent parmi les principaux produits d'exportation.
\item Le contraste « beaucoup d'actifs, peu de richesse produite » traduit une faible productivité agricole.
\end{itemize}
\end{aretenir}

\courssub{Les conditions naturelles favorables}

Pourquoi le Cameroun peut-il presque tout produire ? Observons une carte climatique : du sud au nord, le pays s'étire sur plus de \SI{1200}{km}, de l'équateur presque jusqu'au Sahel. Cette observation cartographique simple permet d'établir le raisonnement suivant : à chaque climat correspond un type de végétation, donc un type de production possible.

\begin{itemize}
\item \textbf{La diversité des climats.} Le Sud connaît un \cle{climat équatorial} (chaud et humide toute l'année, quatre saisons), favorable aux cultures pérennes (cacao, palmier à huile, hévéa, banane). Le Nord connaît un \cle{climat soudanien} puis sahélien (une longue saison sèche et une courte saison des pluies), favorable aux cultures annuelles (mil, sorgho, coton, arachide) et à l'élevage. L'Ouest et le Nord-Ouest connaissent un climat tropical de montagne, plus frais, favorable au café arabica, aux maraîchages et aux pommes de terre.
\item \textbf{La diversité des sols.} Sols ferrallitiques du Sud (adaptés aux cultures pérennes malgré leur pauvreté), sols volcaniques très fertiles de l'Ouest (mont Cameroun, hautes terres), sols tropicaux du Nord adaptés au coton.
\item \textbf{L'hydrographie.} De nombreux fleuves et rivières (Sanaga, Wouri, Nyong, Bénoué, Logone), des lacs (lac Tchad, lacs de Barombi) et des précipitations abondantes au Sud assurent l'eau nécessaire aux cultures, au bétail et à la pêche continentale.
\item \textbf{La façade maritime.} Environ \SI{400}{km} de côte sur l'océan Atlantique (golfe de Guinée), de Kribi à Limbe en passant par Douala, offrent un domaine de pêche maritime et des ports pour exporter les produits agricoles.
\end{itemize}

\begin{savaistu}[Le Cameroun, « l'Afrique en miniature »]
Le Cameroun est surnommé \emph{l'Afrique en miniature} car il réunit presque tous les climats et paysages du continent : forêt dense équatoriale au Sud, savanes au Nord, hautes montagnes à l'Ouest (le mont Cameroun culmine à \SI{4095}{m}), steppes sahéliennes à l'Extrême-Nord, et une ouverture sur la mer. Cette diversité exceptionnelle explique la variété de ses productions : on y cultive aussi bien la banane et le cacao que le mil et le coton, on y élève des zébus comme on y pêche des capitaines.
\end{savaistu}

\courssub{Les conditions humaines}

Les potentialités naturelles ne suffisent pas : il faut des hommes, des savoir-faire et des décisions politiques pour les transformer en productions.

\begin{itemize}
\item \textbf{Une main-d'œuvre rurale abondante.} Une grande partie de la population camerounaise vit encore en milieu rural et travaille dans les champs. Cette main-d'œuvre nombreuse compense la faible mécanisation : la houe, la machette et la daba restent les outils dominants.
\item \textbf{Des traditions culturales anciennes.} Chaque région possède des savoir-faire hérités : cultures sur brûlis et rotation des parcelles dans le Sud forestier, élevage transhumant des Peuls (Foulbé) sur les pâturages de l'Adamaoua et du Nord, pêche artisanale chez les communautés côtières et lacustres, cultures de montagne en terrasses dans les Grassfields de l'Ouest.
\item \textbf{Les politiques agricoles de l'État.} Depuis l'indépendance, l'État a créé des organismes d'encadrement et de développement : la SODECAO pour le cacao et le café, la SODECOTON pour le coton dans le Nord, la CDC et la SOCAPALM pour les plantations industrielles (banane, palmier à huile, hévéa), le MINEPIA (ministère de l'Élevage, des Pêches et des Industries animales). Des plans et stratégies de développement rural visent à moderniser le secteur et à assurer la sécurité alimentaire.
\end{itemize}

\begin{methode}[Analyser les facteurs d'une activité agricole (démarche exigible)]
Pour expliquer pourquoi une production est localisée dans une région, procéder en trois étapes :
\begin{enumerate}
\item Identifier les \textbf{facteurs naturels} : climat (pluies, températures), sols, relief, eau disponible.
\item Identifier les \textbf{facteurs humains} : main-d'œuvre, traditions, encadrement de l'État, marchés et voies de commercialisation.
\item \textbf{Conclure} en montrant comment la combinaison des deux explique la localisation (exemple : le coton au Nord = climat soudanien à saison sèche marquée + encadrement de la SODECOTON + main-d'œuvre rurale dense).
\end{enumerate}
\end{methode}

\courssub{La diversité des espaces de production}

L'observation d'une carte des productions montre que le Cameroun se partage en grandes \cle{zones agro-écologiques}, c'est-à-dire en ensembles où un type de milieu naturel correspond à un type de production dominant.

\begin{popfigure}
\begin{center}
\begin{tikzpicture}[scale=1.05]
% Contour simplifié du Cameroun
\draw[thick,popdark,fill=popcream] (0,0) -- (1.4,-0.2) -- (2.2,0.6) -- (2.6,1.8) -- (3.4,2.6) -- (4.6,3.4) -- (4.4,4.6) -- (3.2,5.4) -- (2.6,6.6) -- (2.0,7.6) -- (1.2,8.2) -- (0.6,7.2) -- (0.2,6.0) -- (-0.4,4.6) -- (-0.8,3.2) -- (-0.6,1.6) -- cycle;
% Zones agro-écologiques
\fill[popgreen!55] (-0.6,1.6) -- (-0.8,3.2) -- (-0.4,4.6) -- (0.2,6.0) -- (0.8,5.6) -- (1.0,4.0) -- (0.8,2.4) -- (0.2,1.2) -- cycle;
\fill[popblue!45] (0.8,2.4) -- (1.0,4.0) -- (0.8,5.6) -- (1.6,6.2) -- (2.4,5.6) -- (2.6,4.2) -- (2.2,2.6) -- (1.6,1.6) -- cycle;
\fill[poppurple!40] (2.2,2.6) -- (2.6,4.2) -- (2.4,5.6) -- (3.0,6.2) -- (3.4,5.0) -- (3.6,3.6) -- (3.0,2.4) -- cycle;
\fill[popgold!55] (3.0,2.4) -- (3.6,3.6) -- (3.4,5.0) -- (3.0,6.2) -- (2.6,6.6) -- (2.0,7.6) -- (1.2,8.2) -- (0.6,7.2) -- (0.2,6.0) -- (0.8,5.6) -- (1.6,6.2) -- (2.4,5.6) -- cycle;
% Côte atlantique
\draw[line width=2.5pt,popblue] (-0.6,1.6) -- (0.2,1.2) -- (0.8,2.4);
% Villes
\fill[popink] (0.4,2.2) circle (0.07) node[right,font=\scriptsize]{Yaoundé};
\fill[popink] (-0.1,1.7) circle (0.07) node[left,font=\scriptsize]{Douala};
\fill[popink] (0.9,5.0) circle (0.06) node[right,font=\scriptsize]{Bafoussam};
\fill[popink] (2.9,4.3) circle (0.06) node[right,font=\scriptsize]{Ngaoundéré};
\fill[popink] (1.6,7.2) circle (0.06) node[right,font=\scriptsize]{Maroua};
% Nord
\draw[->,thick,popdark] (4.2,7.6) -- (4.2,8.4) node[above,font=\scriptsize]{N};
% Légendes des zones
\node[font=\scriptsize,align=center] at (0.1,3.6) {\textbf{Sud}\\forestier};
\node[font=\scriptsize,align=center] at (1.7,4.6) {\textbf{Ouest}\\montagneux};
\node[font=\scriptsize,align=center] at (3.0,3.4) {\textbf{Adamaoua}};
\node[font=\scriptsize,align=center] at (2.0,6.9) {\textbf{Nord}\\soudanien};
\node[font=\scriptsize,popblue,align=center] at (-1.7,1.4) {Côte\\atlantique};
\end{tikzpicture}
\end{center}
\legende{Croquis schématique du Cameroun : les grandes zones agro-écologiques (contour très simplifié, sans précision cartographique). 1 : Sud forestier (vert) ; 2 : Ouest montagneux (bleu) ; 3 : Adamaoua (violet) ; 4 : Nord soudanien (jaune) ; 5 : côte atlantique (trait bleu épais). À compléter en TP avec les productions dominantes de chaque zone.}
\end{popfigure}

\begin{center}
\begin{tabularx}{\linewidth}{|l|X|X|}
\hline
\rowcolor{popblueL} \textbf{Zone} & \textbf{Milieu naturel dominant} & \textbf{Productions principales} \\
\hline
Sud forestier & Climat équatorial, forêt dense & Cacao, plantain, manioc, palmier à huile, hévéa \\
\hline
Ouest montagneux & Climat tropical de montagne, sols volcaniques & Café arabica, banane, maraîchages, pommes de terre, maïs \\
\hline
Adamaoua & Hauts plateaux, savane arborée & Élevage bovin (zébus), maïs, igname \\
\hline
Nord soudanien & Climat soudanien, savane, saison sèche longue & Coton, mil, sorgho, arachide, élevage \\
\hline
Côte atlantique & Climat équatorial côtier, mer & Banane d'exportation, pêche maritime \\
\hline
\end{tabularx}
\end{center}

\begin{exoresolu}[Localiser et expliquer une production]
\textbf{Énoncé.} Un élève affirme : « On cultive le cacao dans l'Extrême-Nord parce que le climat y est sec. » Corriger cette affirmation en justifiant la localisation réelle du cacao au Cameroun.
\tcblower
\textbf{Solution détaillée.} L'affirmation est fausse. Le cacao est une culture pérenne qui exige chaleur et humidité toute l'année : il est donc localisé dans le \textbf{Sud forestier} (régions du Centre, du Sud, de l'Est et du Littoral), où règne le climat équatorial. Le raisonnement suit la méthode en trois étapes : (1) facteurs naturels : pluies abondantes et régulières, températures élevées, couvert forestier qui protège les cacaoyers jeunes ; (2) facteurs humains : main-d'œuvre rurale disponible, encadrement par la SODECAO, accès aux ports de Douala et Kribi pour l'exportation ; (3) conclusion : le cacao ne peut pas prospérer dans l'Extrême-Nord, où la longue saison sèche (plus de six mois sans pluie) le ferait mourir ; là-bas, ce sont le coton, le mil et le sorgho, plantes résistant à la sécheresse, qui dominent.
\end{exoresolu}

\begin{exercice}[Pour vérifier vos acquis]
\begin{enumerate}
\item Définir en une phrase chacun des termes suivants : agriculture, élevage, pêche, secteur agropastoral et piscicole.
\item Expliquer, à l'aide des chiffres du cours, le paradoxe entre la part de l'agriculture dans l'emploi et sa part dans le PIB.
\item Recopier et compléter : « Le Cameroun est surnommé \ldots\ car il réunit \ldots »
\item Sur le croquis des zones agro-écologiques, associer chaque zone à deux productions dominantes et justifier l'une de ces associations par le climat.
\end{enumerate}
\end{exercice}

\begin{corrige}[Exercice « Pour vérifier vos acquis »]
\begin{enumerate}
\item Agriculture : ensemble des activités de production végétale destinées à l'alimentation humaine et animale ou à l'industrie. Élevage : ensemble des activités de production animale (viande, lait, œufs, peaux). Pêche : capture de poissons et de produits aquatiques en milieu naturel ou aménagé. Secteur agropastoral et piscicole : ensemble regroupant l'agriculture, l'élevage et la pêche.
\item L'agriculture occupe environ 60 \% des actifs mais ne fournit qu'environ 15 à 17 \% du PIB : beaucoup de travailleurs produisent une part faible de la richesse, ce qui traduit une faible productivité (peu de mécanisation, faibles rendements) et explique la pauvreté rurale.
\item Le Cameroun est surnommé « l'Afrique en miniature » car il réunit presque tous les climats et paysages du continent africain.
\item Sud forestier : cacao et plantain ; Ouest montagneux : café arabica et maraîchages ; Adamaoua : élevage bovin et maïs ; Nord soudanien : coton et mil ; côte atlantique : banane et pêche maritime. Justification possible : le coton au Nord s'explique par le climat soudanien à saison sèche marquée, qui convient à cette plante et facilite sa récolte à sec.
\end{enumerate}
\end{corrige}

\begin{aretenir}[Bilan de la section]
Les activités agropastorales et piscicoles regroupent l'agriculture (cultures), l'élevage (animaux) et la pêche. Elles occupent environ 60 \% des actifs camerounais mais ne produisent qu'environ 15 à 17 \% du PIB, tout en fournissant d'importants produits d'exportation (cacao, café, coton). Ce secteur repose sur des conditions naturelles exceptionnelles (diversité des climats, des sols, hydrographie, façade maritime) et humaines (main-d'œuvre rurale, traditions, politiques de l'État), réparties en cinq grandes zones agro-écologiques : Sud forestier, Ouest montagneux, Adamaoua, Nord soudanien et côte atlantique. Reste à comprendre pourquoi, malgré ce potentiel, Amina, son oncle et les pêcheurs de Kribi rencontrent tant de difficultés : ce sera l'objet des sections suivantes.
\end{aretenir}

\coursec{L'agriculture au Cameroun}

Dans la section précédente, nous avons vu que les activités agropastorales et piscicoles constituent la base de l'économie rurale camerounaise : elles occupent la majorité de la population active et fournissent l'essentiel de l'alimentation du pays. Parmi ces activités, l'\cle{agriculture} occupe la première place. Mais toutes les agricultures ne se ressemblent pas : le paysan qui cultive son champ de manioc autour de son village ne travaille ni avec les mêmes outils, ni avec les mêmes objectifs que la grande plantation de palmiers à huile du Littoral. Cette section étudie les deux grands types d'agriculture pratiqués au Cameroun — l'agriculture vivrière et l'agriculture de rente — puis leurs apports à l'économie nationale.

\courssub{L'agriculture vivrière}

\textbf{Intuition.} Dans la plupart des villages camerounais, chaque famille possède un ou plusieurs champs. Elle y cultive le maïs, le manioc ou le plantain qu'elle consomme chaque jour. Ce n'est que lorsque la récolte dépasse les besoins du foyer que le surplus est vendu au marché. C'est là l'idée centrale de l'agriculture vivrière : \emph{produire d'abord pour se nourrir}.

\begin{definition}[L'agriculture vivrière]
L'\cle{agriculture vivrière} est une agriculture dont la production est essentiellement destinée à la \textbf{consommation de la famille du producteur}. Elle se pratique sur de petites surfaces (souvent moins de 2 hectares), avec des outils simples et une main-d'œuvre familiale.
\end{definition}

\textbf{Les techniques de l'agriculture vivrière.} L'agriculture vivrière camerounaise repose sur des techniques traditionnelles héritées des générations passées :

\begin{itemize}
\item la \cle{culture sur brûlis} : on défriche une parcelle de forêt ou de savane en coupant la végétation, puis on la brûle ; les cendres fertilisent temporairement le sol ;
\item la \cle{jachère} : après quelques années de culture, le sol s'épuise ; on abandonne alors la parcelle pendant plusieurs années pour qu'elle retrouve sa fertilité naturelle, et on défriche ailleurs ;
\item la \cle{houe} : c'est l'outil de base du paysan ; le travail du sol reste manuel, sans tracteur ni engrais chimique dans la plupart des cas.
\end{itemize}

Ces techniques donnent des \textbf{rendements faibles}, mais elles permettent de nourrir la majorité de la population rurale.

\textbf{Les principales cultures vivrières.} Les cultures vivrières varient selon les régions climatiques du pays :

\begin{itemize}
\item le \cle{maïs} : céréale dominante dans les hautes terres de l'Ouest et du Nord-Ouest ;
\item le \cle{manioc} : tubercule de la zone forestière du Sud, du Centre et du Littoral, consommé sous forme de bâtons, de gari ou de feuilles (feuilles de manioc) ;
\item le \cle{plantain} : banane à cuire, aliment de base du Sud forestier ;
\item le \cle{mil} et le \cle{sorgho} : céréales résistant à la sécheresse, cultivées dans les savanes du Nord (Nord, Extrême-Nord, Adamaoua) ;
\item l'\cle{igname} et l'\cle{arachide} : cultivés un peu partout, l'arachide surtout dans les savanes.
\end{itemize}

\begin{aretenir}[Répartition spatiale des cultures vivrières]
\begin{itemize}
\item \textbf{Ouest et Nord-Ouest} : maïs (hautes terres fraîches) ;
\item \textbf{Nord et Extrême-Nord} : mil et sorgho (savanes sèches) ;
\item \textbf{Sud forestier} (Centre, Sud, Littoral, Est) : plantain et manioc.
\end{itemize}
Cette répartition s'explique par le climat : les céréales exigeantes en eau (maïs) occupent les hautes terres humides, les céréales résistantes (mil, sorgho) les savanes sèches, et les tubercules et bananiers la forêt équatoriale.
\end{aretenir}

\begin{center}
\begin{tabularx}{\linewidth}{|l|X|X|}\hline
\rowcolor{popblueL} \textbf{Culture} & \textbf{Zone principale} & \textbf{Produit consommé} \\ \hline
Maïs & Ouest, Nord-Ouest & grains bouillis, grillés, farine \\ \hline
Manioc & Centre, Sud, Littoral, Est & bâtons, gari, feuilles \\ \hline
Plantain & Sud forestier (Centre, Sud, Littoral, Est, Sud-Ouest) & plantain bouilli, frit, braisé \\ \hline
Mil / Sorgho & Nord, Extrême-Nord, Adamaoua & boule, couscous, bière de mil \\ \hline
Igname & Centre, Ouest, Adamaoua & tubercule bouilli, pilé \\ \hline
Arachide & Savanes (Nord, Adamaoua, Nord-Ouest) & graines grillées, huile, sauce \\ \hline
\end{tabularx}
\end{center}

\begin{popfigure}
\begin{center}
\begin{tikzpicture}
\begin{axis}[
    ybar, bar width=14pt,
    width=0.85\linewidth, height=6cm,
    symbolic x coords={Manioc, Plantain, Maïs, Mil/Sorgho},
    xtick=data,
    ylabel={Production (millions de tonnes)},
    ymin=0, ymax=6,
    ytick={0,1,2,3,4,5,6},
    nodes near coords,
    every node near coord/.append style={font=\small},
    axis lines=left,
    ymajorgrids=true, grid style={popline!40},
]
\addplot[fill=popblue!70, draw=popblue] coordinates {(Manioc,5.4) (Plantain,4.5) (Maïs,2.2) (Mil/Sorgho,1.2)};
\end{axis}
\end{tikzpicture}
\end{center}
\legende{Production approximative de quelques cultures vivrières au Cameroun (ordres de grandeur récents, en millions de tonnes). Le manioc et le plantain dominent la production vivrière nationale.}
\end{popfigure}

\begin{attention}[Erreur fréquente]
Ne crois pas que l'agriculture vivrière ne rapporte \emph{aucun} revenu. Une partie de la production — les \cle{excédents} — est vendue sur les marchés locaux : le plantain de la région du Centre alimente les marchés de Yaoundé, le maïs de l'Ouest est vendu à Douala. L'agriculture vivrière est donc \emph{essentiellement} (et non \emph{exclusivement}) destinée à l'autoconsommation.
\end{attention}

\begin{savaistu}[Le plantain, un trésor camerounais]
La banane plantain est un aliment de base au Cameroun, qui figure parmi les \textbf{premiers producteurs mondiaux} de plantain. Consommé bouilli, frit (alloco) ou braisé (poulet braisé-plantain), il nourrit des millions de citadins et fait vivre des centaines de milliers de producteurs et de commerçantes (les « bayam-sellam »).
\end{savaistu}

\courssub{L'agriculture de rente (ou de plantation)}

\textbf{Intuition.} À l'opposé du petit champ familial, certaines exploitations couvrent des milliers d'hectares, emploient des centaines de salariés et vendent toute leur production à l'étranger. C'est une agriculture d'entreprise, tournée vers le marché mondial.

\begin{definition}[L'agriculture de rente]
L'\cle{agriculture de rente} (ou agriculture de plantation) est une agriculture \textbf{commerciale} pratiquée sur de \textbf{grandes surfaces}, destinée à la \textbf{vente} et surtout à l'\textbf{exportation}. Elle se caractérise par la monoculture, l'emploi de capitaux importants, une main-d'œuvre salariée et des techniques modernes.
\end{definition}

\textbf{Les caractéristiques de l'agriculture de rente :}
\begin{enumerate}
\item de \textbf{grandes surfaces} : des plantations de plusieurs centaines ou milliers d'hectares ;
\item la \textbf{monoculture} : une seule culture par plantation (que du palmier à huile, que de la banane...) ;
\item des \textbf{capitaux importants} : machines, engrais, pesticides, usines de transformation ;
\item une \textbf{orientation vers l'exportation} : la production est vendue sur le marché mondial et rapporte des devises.
\end{enumerate}

\textbf{Les principales cultures d'exportation et leur localisation :}

\begin{itemize}
\item le \cle{cacao} : première culture d'exportation, cultivé dans la zone forestière du \textbf{Centre-Sud} (Centre, Sud, Littoral, Est) ;
\item le \cle{café} : l'\textbf{arabica} sur les hautes terres de l'\textbf{Ouest} (et Nord-Ouest), le \textbf{robusta} dans la zone forestière du \textbf{Centre-Sud} ;
\item le \cle{coton} : la culture de rente des savanes du \textbf{Nord} (Nord, Extrême-Nord, Adamaoua) ;
\item la \cle{banane} (banane douce) : dans la plaine côtière du \textbf{Littoral} (Penja) et du \textbf{Sud-Ouest} (Tiko, Muyuka) ;
\item le \cle{palmier à huile} : dans la zone côtière et forestière humide (Littoral, Sud-Ouest, Centre, Sud) ;
\item l'\cle{hévéa} (caoutchouc) : dans le Sud forestier (Sud, Centre, Littoral) ;
\item le \cle{thé} : sur les hauts plateaux (Nord-Ouest, Adamaoua Ouest).
\end{itemize}

\begin{exemplebox}[Le cacao et le coton, deux cultures emblématiques]
Le Cameroun est le \textbf{1\textsuperscript{er} producteur de cacao de la CEMAC} et figure parmi les premiers producteurs africains (derrière la Côte d'Ivoire et le Ghana). Le cacao est cultivé par de petits planteurs du Centre-Sud. De son côté, la \textbf{SODECOTON} encadre des \textbf{centaines de milliers de cotonculteurs} dans le Nord : elle fournit semences, engrais et intrants, puis achète et transforme le coton-graine. Le coton fait vivre une grande partie de la population rurale du Grand Nord.
\end{exemplebox}

\begin{popfigure}
\begin{center}
\begin{tikzpicture}[scale=0.95]
% Contour très simplifié du Cameroun
\draw[thick, popdark, fill=popcream] 
  (0,0) -- (2.2,-0.2) -- (3.4,0.6) -- (3.8,1.8) -- (3.2,3.0) -- (2.6,3.4) 
  -- (2.2,4.6) -- (1.7,5.6) -- (1.9,6.6) -- (1.4,7.4) -- (0.8,6.4) -- (0.6,5.2) 
  -- (0.2,4.2) -- (0.5,3.0) -- (-0.3,2.0) -- (-0.5,0.8) -- cycle;
% Nord
\draw[->, thick, popdark] (4.6,6.8) -- (4.6,7.6);
\node[above] at (4.6,7.6) {\small \textbf{N}};
% Zones cultures
\node[popdark, font=\small, align=center] at (1.3,6.0) {Coton};
\fill[poppurple] (1.3,5.6) circle (3pt);
\node[popdark, font=\small, align=center] at (0.9,3.4) {Café arabica};
\fill[popgold!80!black] (0.9,3.0) circle (3pt);
\node[popdark, font=\small, align=center] at (2.4,1.4) {Cacao, café robusta};
\fill[popgold!80!black] (2.0,1.2) circle (3pt);
\node[popdark, font=\small, align=center] at (0.4,0.4) {Banane};
\fill[popgreen] (0.0,0.2) circle (3pt);
\node[popdark, font=\small, align=center] at (2.9,0.2) {Palmier à huile\\ Hévéa};
\fill[poporange] (2.4,0.5) circle (3pt);
% Villes repères
\fill[popink] (0.3,1.0) circle (2pt); \node[left, font=\footnotesize] at (0.3,1.0) {Douala};
\fill[popink] (1.6,1.6) circle (2pt); \node[right, font=\footnotesize] at (1.6,1.6) {Yaoundé};
\fill[popink] (1.5,6.8) circle (2pt); \node[right, font=\footnotesize] at (1.5,6.8) {Maroua};
% Légende
\node[font=\footnotesize, anchor=west] at (4.2,3.6) {\textcolor{popgold!80!black}{$\bullet$} Cacao / Café robusta};
\node[font=\footnotesize, anchor=west] at (4.2,3.1) {\textcolor{poppurple}{$\bullet$} Coton};
\node[font=\footnotesize, anchor=west] at (4.2,2.6) {\textcolor{popgreen}{$\bullet$} Banane};
\node[font=\footnotesize, anchor=west] at (4.2,2.1) {\textcolor{poporange}{$\bullet$} Palmier / Hévéa};
\node[font=\footnotesize, anchor=west] at (4.2,1.6) {\textcolor{popgold!80!black}{$\bullet$} Café arabica};
\end{tikzpicture}
\end{center}
\legende{Croquis schématique de localisation des principales cultures de rente du Cameroun : coton dans le Nord, café arabica à l'Ouest, cacao et café robusta au Centre-Sud, banane et palmier à huile/hévéa sur la côte. Échelle indicative : le pays mesure environ \SI{1200}{km} du sud au nord.}
\end{popfigure}

\begin{methode}[Construire un croquis de localisation des cultures]
\begin{enumerate}
\item \textbf{Tracer le fond de carte} : contour simplifié du pays, limites des grandes zones (Nord, Ouest, Centre-Sud, côte), principales villes repères.
\item \textbf{Placer les symboles} : un symbole ou une couleur par culture, positionné dans sa zone de production réelle.
\item \textbf{Rédiger la légende} : chaque symbole doit être expliqué ; ajouter l'orientation (flèche du nord) et une échelle indicative.
\item \textbf{Donner un titre} précis au croquis (exemple : « Les principales cultures de rente du Cameroun »).
\end{enumerate}
Au Probatoire et au Baccalauréat technique, un croquis sans titre ou sans légende perd des points.
\end{methode}

\textbf{Les grandes sociétés agro-industrielles.} L'agriculture de rente est dominée par de grandes sociétés, souvent créées à l'époque coloniale et aujourd'hui partiellement privatisées :

\begin{center}
\begin{tabularx}{\linewidth}{|l|X|X|}\hline
\rowcolor{popblueL} \textbf{Société} & \textbf{Activité principale} & \textbf{Localisation principale} \\ \hline
CDC (Cameroon Development Corporation) & Banane, palmier à huile, hévéa & Sud-Ouest (Tiko, Muyuka, Ekona) \\ \hline
PHP (Plantations du Haut Penja) & Banane (exportation) & Littoral (Penja) \\ \hline
SOCAPALM & Huile de palme & Littoral (Dibombari, Mbongo), Centre (Eseka), Sud (Kienke) \\ \hline
SODECOTON & Coton (encadrement, transformation) & Nord, Extrême-Nord, Adamaoua \\ \hline
HEVECAM & Hévéa (caoutchouc naturel) & Sud (Bipindi, Kribi, Niete) \\ \hline
\end{tabularx}
\end{center}

\begin{exoresolu}[Localiser les cultures de rente]
\textbf{Énoncé.} Un élève affirme : « Le coton est la principale culture de rente du Sud forestier et le cacao celle du Nord. » Corrige cette affirmation et justifie ta réponse.
\tcblower
\textbf{Solution.} L'affirmation inverse les localisations. Le \textbf{coton} est la culture de rente des \textbf{savanes du Nord} (Nord, Extrême-Nord, Adamaoua) : c'est une plante qui supporte la saison sèche marquée et qui est encadrée par la SODECOTON. Le \textbf{cacao}, lui, est la principale culture de rente de la \textbf{zone forestière du Centre-Sud} (Centre, Sud, Littoral, Est), car le cacaoyer exige chaleur et humidité toute l'année. La répartition des cultures de rente suit donc les grandes zones climatiques : savane au Nord, forêt au Sud.
\end{exoresolu}

\courssub{Agriculture traditionnelle et agriculture moderne}

On peut opposer deux grands modèles agricoles au Cameroun, qui coexistent dans le paysage rural.

\begin{popfigure}
\begin{center}
\begin{tikzpicture}[scale=1]
% Colonne gauche : traditionnelle
\draw[fill=poporangeL, draw=poporange, thick, rounded corners] (-4.4,-2.6) rectangle (-0.2,2.6);
\node[font=\small\bfseries, popdark] at (-2.3,2.2) {Agriculture traditionnelle (vivrière)};
\node[font=\footnotesize, align=left, anchor=north west] at (-4.3,1.8) {
$\bullet$ Petites surfaces ($<$ 2 ha)\\[2pt]
$\bullet$ Outils simples : houe, machette\\[2pt]
$\bullet$ Main-d'œuvre familiale\\[2pt]
$\bullet$ Brûlis, jachère, pas d'engrais\\[2pt]
$\bullet$ Rendements faibles\\[2pt]
$\bullet$ But : nourrir la famille\\ \phantom{$\bullet$ } (excédents vendus localement)};
% Colonne droite : moderne
\draw[fill=popgreenL, draw=popgreen, thick, rounded corners] (0.2,-2.6) rectangle (4.4,2.6);
\node[font=\small\bfseries, popdark] at (2.3,2.2) {Agriculture moderne (de rente / encadrée)};
\node[font=\footnotesize, align=left, anchor=north west] at (0.3,1.8) {
$\bullet$ Grandes surfaces (plantations)\\[2pt]
$\bullet$ Machines, tracteurs, engrais\\[2pt]
$\bullet$ Main-d'œuvre salariée\\[2pt]
$\bullet$ Semences sélectionnées, irrigation\\[2pt]
$\bullet$ Rendements élevés\\[2pt]
$\bullet$ But : vendre et exporter\\ \phantom{$\bullet$ } (recherche du profit)};
% Flèche centrale
\draw[<->, very thick, poppurple] (-0.15,0) -- (0.15,0);
\node[font=\footnotesize, poppurple, rotate=90] at (0, -0.5) {Transition possible};
\end{tikzpicture}
\end{center}
\legende{Schéma comparatif : agriculture traditionnelle et agriculture moderne au Cameroun. Les deux modèles se distinguent par les techniques, les rendements et les objectifs. De nombreux paysans sont en transition (semences améliorées, engrais).}
\end{popfigure}

L'\textbf{agriculture traditionnelle} correspond à l'agriculture vivrière : petites parcelles, outils manuels, main-d'œuvre familiale, rendements faibles, objectif d'autoconsommation. L'\textbf{agriculture moderne} correspond surtout aux plantations (CDC, PHP, SOCAPALM, HEVECAM) et aux exploitations paysannes encadrées (coton avec SODECOTON, café/cacao avec ONCC/structures d'appui) : mécanisation partielle ou totale, engrais et pesticides, semences améliorées, rendements élevés, production destinée à la vente. Entre les deux, de nombreux paysans camerounais sont en transition : ils commencent à utiliser des engrais et des semences améliorées sur leurs petits champs.

\courssub{Les apports de l'agriculture à l'économie camerounaise}

L'agriculture joue un rôle majeur dans la vie nationale :

\begin{enumerate}
\item une \textbf{autosuffisance alimentaire relative} : le Cameroun produit l'essentiel de ce qu'il consomme (manioc, plantain, maïs, mil), ce qui est rare en Afrique centrale ; le pays est parfois appelé « le grenier de la CEMAC » ;
\item des \textbf{devises} : le cacao, le café, le coton, la banane et l'huile de palme sont exportés et rapportent des devises étrangères indispensables au budget de l'État ;
\item des \textbf{emplois} : l'agriculture occupe la majorité de la population active, surtout en milieu rural ;
\item des \textbf{matières premières pour l'industrie} : le coton alimente la filière textile (CICAM), le cacao les chocolateries (Chococam), l'huile de palme les savonneries (Socapalm, Savonnerie du Cameroun), le maïs et le sorgho les brasseries (SABC, UCAM).
\end{enumerate}

\begin{aretenir}[L'essentiel de la section]
\begin{itemize}
\item L'\textbf{agriculture vivrière} (maïs, manioc, plantain, mil, igname, arachide) nourrit les familles ; ses techniques sont le brûlis, la jachère et la houe.
\item L'\textbf{agriculture de rente} (cacao, café, coton, banane, palmier à huile, hévéa, thé) produit pour l'exportation sur de grandes surfaces.
\item La répartition des cultures suit le climat : coton et mil dans les savanes du Nord, maïs et café arabica sur les hautes terres de l'Ouest, cacao, plantain et palmier dans le Sud forestier.
\item L'agriculture fournit nourriture, devises, emplois et matières premières : c'est le pilier de l'économie camerounaise.
\end{itemize}
\end{aretenir}

\begin{exercice}[Pour s'entraîner]
\begin{enumerate}
\item Définis l'agriculture vivrière et l'agriculture de rente en donnant pour chacune deux exemples de cultures.
\item Construis un tableau présentant trois différences entre l'agriculture traditionnelle et l'agriculture moderne.
\item Sur un croquis du Cameroun, localise le coton, le cacao, le café arabica et la banane, avec titre et légende.
\end{enumerate}
\end{exercice}

\begin{corrige}[Exercice]
\begin{enumerate}
\item L'agriculture vivrière est une agriculture dont la production est essentiellement destinée à la consommation de la famille du producteur (exemples : manioc, plantain). L'agriculture de rente est une agriculture commerciale pratiquée sur de grandes surfaces, destinée à la vente et surtout à l'exportation (exemples : cacao, coton).
\item Tableau attendu : surfaces (petites / grandes), outils (houe / machines et engrais), main-d'œuvre (familiale / salariée), rendements (faibles / élevés), objectif (autoconsommation / vente et exportation). Trois de ces différences suffisent.
\item Croquis : coton dans le Nord, cacao au Centre-Sud, café arabica à l'Ouest, banane sur la côte (Littoral et Sud-Ouest). Ne pas oublier le titre, la légende et la flèche du nord.
\end{enumerate}
\end{corrige}

\coursec{L'élevage au Cameroun}

Après avoir analysé les systèmes agraires et les productions végétales qui structurent l'espace rural camerounais, nous nous tournons maintenant vers la seconde composante majeure du secteur primaire : l'élevage. Au Cameroun, l'agriculture et l'élevage sont intimement liés, tant par les interactions écologiques (fumure, traction animale) que par les conflits d'usage de l'espace (divagation, couloirs de transhumance). Cette section examine la diversité des systèmes d'élevage, leur répartition spatiale, la dynamique de la transhumance, l'émergence de filières modernes et les défis structurels qui limitent la couverture des besoins nationaux en protéines animales.

\courssub{Généralités et types d'élevage}

L'élevage désigne l'ensemble des techniques mises en œuvre par l'homme pour élever, nourrir, faire reproduire et soigner des animaux domestiques en vue d'en tirer des produits (viande, lait, œufs, cuir, force de travail) ou des services (transport, labour, épargne sur pied). Au Cameroun, deux grands modèles coexistent, souvent imbriqués au sein d'une même exploitation familiale.

\begin{definition}[Élevage extensif (traditionnel)]
Système d'élevage où les animaux paissent librement sur de vastes espaces naturels (parcours, jachères, brousse) avec très peu d'intrants extérieurs (aliments achetés, médicaments, bâtiments). La productivité par tête est faible, mais la productivité par hectare ou par unité de travail peut être rationnelle dans les zones sèches. C'est le mode dominant au Nord et dans l'Adamaoua.
\end{definition}

\begin{definition}[Élevage intensif (moderne)]
Système caractérisé par un fort niveau d'intrants par unité de surface ou par animal : alimentation concentrée, sélection génétique, prophylaxie rigoureuse, bâtiments spécialisés. Il vise une productivité maximale (lait, viande, œufs) sur de petites surfaces. Il se développe en zone périurbaine (Yaoundé, Douala, Bafoussam) et dans quelques ranches.
\end{definition}

Le cheptel national comprend cinq grandes espèces, aux aires de répartition contrastées :
\begin{itemize}
    \item \textbf{Bovins :} Zébu (Blanc du Logone, Goudali, Bakosi, Doayo/Kapsiki), Taurins (Namchi). Effectif estimé à \num{6,5} millions de têtes (données MINEPIA récentes).
    \item \textbf{Ovins :} Moutons à queue fine (Sahel, Nord) et à queue grasse (Adamaoua, Ouest). Effectif : \num{3,5} millions.
    \item \textbf{Caprins :} Chèvres naines (Sud) et chèvres saheliennes (Nord). Effectif : \num{4,2} millions.
    \item \textbf{Porcins :} Élevage traditionnel en divagation (Sud) et porcheries modernes (périurbain). Effectif : \num{1,8} million.
    \item \textbf{Volailles :} Poulets locaux (cycle long, rustiques) et poulets de chair/pondeuses (cycle court, 42 jours). Effectif : \num{35} millions (dont \num{80\%} villageois).
\end{itemize}
On note aussi l'importance des \textbf{animaux de basse-cour} (pintades, pigeons, canards) et des \textbf{animaux de trait} (bœufs de labour, ânes au Nord) pour l'agriculture.

\courssub{Les grandes zones d'élevage et la transhumance}

La répartition spatiale du cheptel bovin, qui concentre l'essentiel de la biomasse et des enjeux fonciers, est dictée par l'écologie (disponibilité en eau et pâturages, absence de mouche tsé-tsé) et l'histoire (peuplement Peul).

\begin{exemplebox}[L'Adamaoua : cœur de l'élevage bovin]
L'Adamaoua, avec ses hauts plateaux herbeux (alt. 1000--1300 m) débarrassés de la mouche tsé-tsé grâce à l'altitude et au brûlis séculaire, concentre la plus grande partie du cheptel bovin camerounais (plus de \num{40\%} des effectifs nationaux, soit plusieurs millions de têtes). C'est une zone de \emph{montagne pâturée} par excellence, où l'élevage est l'activité dominante, structurant le paysage (parcs à bœufs, cases peules, points d'eau).
\end{exemplebox}

Deux autres pôles majeurs complètent le dispositif :
\begin{itemize}
    \item \textbf{Le Nord et l'Extrême-Nord (Bénoué, Mayo-Rey, Logone, Diamaré) :} Zones soudanaises et sahéliennes. Élevage extensif transhumant. Forts effectifs ovins/caprins. Contraintes : sécheresses récurrentes, pression foncière (coton, riz), insécurité (Boko Haram dans l'Extrême-Nord).
    \item \textbf{Les hautes terres de l'Ouest (Bamiléké, Bamoun) :} Élevage bovin taurin (race Namchi/Kapsiki) en stabulation partielle ou parcours sur jachères. Forte densité humaine $\rightarrow$ parcelles petites, intégration agriculture-élevage (fumure). Développement récent de l'aviculture et porcherie intensives.
\end{itemize}

\begin{attention}[Piège : l'élevage bovin dans le Sud forestier]
Erreur fréquente : situer l'élevage bovin dans le Sud forestier (Régions du Centre, Sud, Est, Littoral hors zones côtières). La \textbf{mouche tsé-tsé} (vecteur de la \emph{trypanosomiase} ou \emph{nagana}) y est omniprésente en dessous de 800--1000 m d'altitude. L'humidité élevée et la forêt dense limitent drastiquement l'élevage bovin. Seules quelques races trypanotolérantes (Doayo, Bakosi, Muturu) y survivent en petits effectifs. L'élevage y est porcin et avicole.
\end{attention}

\paragraph{La transhumance : définition et mécanismes.}

\begin{definition}[Transhumance]
Déplacement saisonnier, cyclique et organisé des troupeaux (et d'une partie de la famille de l'éleveur) entre des aires géographiques complémentaires à la recherche de pâturages et d'eau. Au Cameroun, c'est une \textbf{transhumance Nord-Sud} (ou montagne-plaine).
\end{definition}

\begin{itemize}
    \item \textbf{Saison sèche (nov.--mars) :} Les troupeaux quittent l'Adamaoua et le Nord (pâturages brûlés, points d'eau taries) vers le Sud : plaine de la Bénoué, bassins du Logone, piémonts de l'Adamaoua, voire bordure de la forêt (zones de \emph{soudure}).
    \item \textbf{Saison des pluies (avr.--oct.) :} Retour vers le Nord/Adamaoua (herbe jeune, points d'eau pleins, éloignement des cultures pluviales).
\end{itemize}

\begin{savaistu}[Conflits agriculteurs-éleveurs]
Les conflits entre agriculteurs (souvent autochtones ou migrants récents) et éleveurs peuls (souvent perçus comme allogènes) sont récurrents dans le Nord et l'Adamaoua. Causes : \textbf{divagation des bœufs dans les champs} (coton, maïs, sorgho) avant la récolte ou juste après ; blocage des \textbf{couloirs de transhumance} par l'extension des cultures (cottonnières SODECOTON, rizicoles SEMRY) ; non-respect du calendrier pastoral ; absence de bornage des pistes. Conséquences : destructions de récoltes, blessures/morts, insécurité, blocage de la commercialisation du bétail.
\end{savaistu}

\paragraph{Illustration : Zones d'élevage et flux de transhumance.}

\begin{popfigure}
\begin{tikzpicture}[scale=0.85]
% Fond carte Cameroun simplifié (polygone approximatif)
\draw[popdark, thick, fill=popcream] 
(0,0) -- (1,0.5) -- (1.5,1.5) -- (1.2,3) -- (0.5,4) -- (-0.5,4.5) -- (-1.5,4) -- (-2,3) -- (-2.5,2) -- (-2.2,1) -- (-1.5,0.2) -- (-0.5,0) -- cycle;

% Villes repères
\node[circle,fill=popblue,inner sep=1.5pt,label=above right:\tiny Yaoundé] (Y) at (-0.2,1.8) {};
\node[circle,fill=popblue,inner sep=1.5pt,label=above right:\tiny Douala] (D) at (-1.0,0.8) {};
\node[circle,fill=popblue,inner sep=1.5pt,label=above:\tiny Ngaoundéré] (N) at (0.2,3.2) {};
\node[circle,fill=popblue,inner sep=1.5pt,label=above:\tiny Garoua] (G) at (0.5,4.0) {};
\node[circle,fill=popblue,inner sep=1.5pt,label=above:\tiny Maroua] (M) at (0.8,4.5) {};
\node[circle,fill=popblue,inner sep=1.5pt,label=left:\tiny Bafoussam] (B) at (-1.2,2.5) {};

% Zones d'élevage (hachures/remplissage)
% Adamaoua
\fill[popgreen!30, pattern=north west lines, pattern color=popgreen] (-0.5,2.8) ellipse (1.2 and 0.6);
\node[popgreen!80!popdark, font=\bfseries\small] at (-0.5,2.8) {ADAMAOUA};

% Nord / Extrême-Nord
\fill[poporange!30, pattern=north east lines, pattern color=poporange] (0.5,3.5) ellipse (1.0 and 0.8);
\node[poporange!80!popdark, font=\bfseries\small] at (0.5,3.5) {NORD / EXTRÊME-NORD};

% Ouest (Hautes terres)
\fill[poppurple!30, pattern=dots, pattern color=poppurple] (-1.5,2.2) ellipse (0.6 and 0.5);
\node[poppurple!80!popdark, font=\bfseries\small] at (-1.5,2.2) {OUEST};

% Sud forestier (peu d'élevage bovin)
\fill[popmuted!20] (-1.5,1) ellipse (1.0 and 1.2);
\node[popmuted, font=\small\itshape, align=center] at (-1.5,1){Zone forestière\\(tsé-tsé)};

% Flèches Transhumance (Saison sèche : Nord -> Sud)
\draw[popred, very thick, -{Stealth[length=3mm]}, bend left=20] (0.5,3.8) to node[midway, above, sloped, font=\tiny, popred] {Saison sèche} (-0.5,1.5);
\draw[popred, very thick, -{Stealth[length=3mm]}, bend right=15] (0.2,3.0) to node[midway, below, sloped, font=\tiny, popred] {Saison sèche} (-1.0,1.0);

% Flèches Retour (Saison des pluies : Sud -> Nord)
\draw[popblue, very thick, -{Stealth[length=3mm]}, bend right=20] (-0.5,1.5) to node[midway, below, sloped, font=\tiny, popblue] {Saison pluies} (0.5,3.8);
\draw[popblue, very thick, -{Stealth[length=3mm]}, bend left=15] (-1.0,1.0) to node[midway, above, sloped, font=\tiny, popblue] {Saison pluies} (0.2,3.0);

% Légende manuelle
\node[anchor=south west, font=\tiny, fill=white, fill opacity=0.8, text opacity=1, rounded corners, inner sep=2pt, align=center] at (-2.8, -0.5){
\textbf{Légende :}\\
\hspace{2pt}\tikz\draw[popgreen, pattern=north west lines] (0,0) rectangle (4mm,2mm); Adamaoua (Bovins ++) \\
\hspace{2pt}\tikz\draw[poporange, pattern=north east lines] (0,0) rectangle (4mm,2mm); Nord/Extrême-Nord (Bovins, Ovins, Caprins) \\
\hspace{2pt}\tikz\draw[poppurple, pattern=dots] (0,0) rectangle (4mm,2mm); Ouest (Bovins taurins, Aviculture) \\
\hspace{2pt}\tikz\draw[popred, very thick, -{Stealth[length=2mm]}] (0,0) -- (4mm,0); Transhumance saison sèche \\
\hspace{2pt}\tikz\draw[popblue, very thick, -{Stealth[length=2mm]}] (0,0) -- (4mm,0); Retour saison pluies \\
\hspace{2pt}\tikz\fill[popblue] (0,0) circle (1.5pt); Villes principales
};
% Flèche Nord
\node[anchor=south] at (0,4.8) {\textbf{N}};
\draw[->, thick] (0,4.6) -- (0,4.9);
\end{tikzpicture}
\legende{Carte schématique du Cameroun : principales zones d'élevage (hachurées) et flux de transhumance bovins (flèches rouges : saison sèche vers le sud ; flèches bleues : retour saison des pluies vers le nord). La zone forestière au sud est défavorable à l'élevage bovin (mouche tsé-tsé).}
\end{popfigure}

\courssub{L'élevage moderne : ranches et filières periurbaines}

Face à l'insuffisance de l'offre extensive et à l'urbanisation galopante, un élevage moderne se structure, soutenu par l'État (MINEPIA, projets PNDP, PADFA) et le privé.

\begin{itemize}
    \item \textbf{Ranches (élevage bovin amélioré en parcours clôturé) :} Exemples : Ranch de \emph{Ngaoundaba} (Adamaoua, station de sélection Goudali), Ranch de \emph{Ladji} (Nord), Ranch de \emph{Mbanga} (Littoral, race Bakosi/Doayo), Ranch de \emph{Sangmelima} (Sud, race trypanotolérante). Objectifs : sélection génétique, production de semence, démonstration, engraissement. Surfaces limitées, coûts élevés (clôtures, eau, feu).
    \item \textbf{Aviculture industrielle (poulets de chair et pondeuses) :} Concentrée autour de \textbf{Yaoundé} (Mfou, Obala, Monatélé), \textbf{Douala} (Nkongsamba, Manjo, Penja) et \textbf{Bafoussam} (Bamougoum, Bandjoun). Cycle court (42 j pour chair, 18 mois pour ponte). Aliments composés (maïs, tourteau soja, prémix) souvent importés ou fabriqués localement (SMA, ALZO, PROVIMI). Problème : coût aliment (70\% des charges), concurrence importations illégales (poulets surgelés), grippe aviaire.
    \item \textbf{Porcherie industrielle :} Même bassins periurbains. Races : Large White, Landrace, Piétrain. Gestion des déjections (biogaz, compost) en devenir.
\end{itemize}

\begin{aretenir}[Dualisme structurel]
L'élevage camerounais reste marqué par un fort \textbf{dualisme} : un secteur \emph{traditionnel extensif} (90\% des effectifs, 80\% de la viande, faible productivité, vulnérabilité climatique) et un secteur \emph{moderne intensif} (10\% effectifs, croissance rapide, forte productivité, dépendance aux intrants importés, concentration periurbaine). La politique vise l'\textbf{intensification durable} du secteur traditionnel (amélioration fourragère, santé animale, sédentarisation partielle) et le soutien au moderne.
\end{aretenir}

\courssub{Apports, limites et perspectives du secteur}

\paragraph{Les apports multiples de l'élevage.}

L'élevage n'est pas seulement une « usine à viande ». Ses fonctions sont systémiques :
\begin{enumerate}
    \item \textbf{Sécurité alimentaire :} Viande (bœuf, mouton, chèvre, porc, volaille), lait (local et importé en poudre), œufs. Couverture \num{\approx 60\%} des besoins en protéines animales (déficit structurel comblé par importations : poisson surgelé, lait poudre, poulets).
    \item \textbf{Revenus et épargne :} « Banque sur pied » pour les éleveurs peuls ; revenus réguliers pour aviculteurs/porculteurs periurbains. Contribution PIB primaire \num{\approx 15\%}.
    \item \textbf{Fumure organique :} Bouse, fientes $\rightarrow$ amendement des sols (coton, maïs, riz, maraîchage). Crucial pour l'agriculture de conservation.
    \item \textbf{Force de traction :} Bœufs de labour (Nord, Adamaoua, Ouest), ânes (Extrême-Nord). Permet l'extension des surfaces cultivées.
    \item \textbf{Produits dérivés :} Cuir/peaux (artisanat, export), cornes, os, biogaz.
    \item \textbf{Lien social/culturel :} Dot, sacrifices, fêtes, prestige (taille du troupeau).
\end{enumerate}

\paragraph{Les limites structurelles.}

\begin{itemize}
    \item \textbf{Faible productivité du Zébu local :} Âge à la première vêlage tardif (4--5 ans), intervalle vêlage-vêlage long (\num{> 18} mois), taux de réforme élevé, rendement carcasse faible (\num{150--180} kg vs \num{300+} kg races améliorées), production laitière faible (\num{1--2} L/jour vs \num{15--20} L race laitière).
    \item \textbf{Santé animale :} \textbf{Trypanosomiase} (mouche tsé-tsé) limite l'expansion vers le Sud. \textbf{Peste bovine} (éradiquée 2011, surveillance active), \textbf{Peste des petits ruminants} (PPR), \textbf{Fièvre aphteuse}, \textbf{Dermatose nodulaire}, \textbf{Grippe aviaire}. Couverture vaccinale insuffisante. Résistance aux acaricides (tiques).
    \item \textbf{Alimentation :} Déficit fourrager en saison sèche (qualité paille, quantité). Coût élevé des aliments concentrés (maïs, soja) pour le moderne. Conflits pâturages/cultures.
    \item \textbf{Commercialisation :} Filières longues, informelles, manque d'abattoirs aux normes (hygiène), pesée à l'œil, évasion fiscale, export informel vers Nigeria/Gabon/Guinée Équatoriale (prix plus attractifs).
    \item \textbf{Foncier :} Absence de titres fonciers pour les parcours, appropriation par élites/agro-industries, blocage couloirs transhumance.
\end{itemize}

\paragraph{Illustration : Structure du cheptel national (effectifs approximatifs).}

\begin{popfigure}
\begin{tikzpicture}
\begin{axis}[
    ybar,
    bar width=12pt,
    width=\linewidth,
    height=8cm,
    enlarge x limits=0.25,
    ylabel={Effectifs (millions de têtes)},
    ylabel style={font=\small},
    xlabel={Espèces},
    xlabel style={font=\small},
    xtick=data,
    xticklabels={Bovins, Ovins, Caprins, Porcins, Volailles ($\times$10)},
    xticklabel style={rotate=30, anchor=east, font=\tiny},
    ymin=0,
    ymax=40,
    ymajorgrids=true,
    grid style=dashed,
    legend style={at={(0.5,-0.3)}, anchor=north, legend columns=-1, font=\tiny, fill=white, draw=popline},
    nodes near coords,
    nodes near coords align={vertical},
    every node near coord/.append style={font=\tiny, rotate=90, anchor=west, yshift=2pt},
]
\addplot[popblue, fill=popblue!60] coordinates {(1,6.5) (2,3.5) (3,4.2) (4,1.8) (5,3.5)}; 
\legend{Effectifs (M têtes) -- Volailles divisées par 10}
\end{axis}
% Note manuelle pour volailles
\node[anchor=north, font=\tiny, poporange] at (current bounding box.south) {Note : L'effectif volailles réel est de \num{35} millions (hors échelle) ; divisé par 10 sur le graphique pour lisibilité comparative. Sources : MINEPIA / FAOSTAT estimations récentes.};
\end{tikzpicture}
\legende{Diagramme en bâtons des effectifs approximatifs des principales espèces d'élevage au Cameroun (en millions de têtes). L'effectif volailles (\num{35} millions) est représenté divisé par 10 pour lisibilité. Le bovin domine la biomasse, la volaille domine les effectifs.}
\end{popfigure}

\begin{methode}[Analyser une photographie d'élevage (type Bac/Probatoire)]
\textbf{Étape 1 : Description factuelle (ce que je vois).} Identifier : les animaux (espèce, race, âge, état corporel), le milieu (pâture, bâtiment, case, point d'eau, culture), les hommes (présence, activité, tenue), les infrastructures (clôture, abreuvoir, mangeoire, case de vaccination).
\textbf{Étape 2 : Identification du système.} Croiser indices : \emph{Extensif/Traditionnel} = animaux en liberté, brousse/parcours, absence de bâtiment, éleveur à pied/cheval, race locale. \emph{Intensif/Moderne} = bâtiment fermé, mangeoires/abreuvoirs artificiels, race améliorée, aliments en sacs, technicien/ouvrier.
\textbf{Étape 3 : Dégager les problèmes/atouts visibles.} Ex. : pâture surpâturée (désherbée), animaux maigres (saison sèche), divagation près de champ (conflit), ou au contraire : animaux engraissés, fourrage vert distribué, bâtiment ventilé, biosécurité (pediluve).
\textbf{Étape 4 : Conclusion géographique.} Situer dans l'espace (zone écologique, proximité ville), relier aux politiques publiques (projet, ranch, zone de transhumance).
\end{methode}

\begin{exoresolu}[Analyse d'un document : Extrait de rapport MINEPIA 2023]
\textbf{Énoncé :} « La production nationale de viande est estimée à 320 000 tonnes équivalent carcasse (TEC) pour une demande de 450 000 TEC. Le déficit est comblé par les importations (poisson, poulets, abats). Le cheptel bovin croît de 2,5\%/an, la population de 2,6\%/an. »
\begin{enumerate}
    \item Calculer le taux d'autosuffisance en viande.
    \item Expliquer pourquoi la croissance du cheptel bovin ne résout pas le déficit.
    \item Proposer deux leviers techniques pour augmenter l'offre sans augmenter les effectifs.
\end{enumerate}
\tcblower
\textbf{Solution :}
\begin{enumerate}
    \item Taux d'autosuffisance = (Production / Demande) $\times$ 100 = (320 000 / 450 000) $\times$ 100 = \textbf{71,1\%}. Le Cameroun importe près de 30\% de sa consommation de viande.
    \item La croissance du cheptel (\num{2,5\%}) est \textbf{inférieure} à la croissance de la demande (population \num{2,6\%} + hausse revenu/urbanisation $\rightarrow$ élasticité-revenu positive). De plus, la productivité par tête reste faible (âge abattage tardif, poids carcasse bas). L'augmentation extensive (plus de têtes) heurte les limites foncières et environnementales.
    \item \textbf{Levier 1 : Amélioration génétique et zootechnique.} Insémination artificielle (IA) avec semence de races laitières/viandes (Girolando, Montbéliarde, Charolaise) sur zébu local $\rightarrow$ croisés F1 précoces, plus lourds, meilleurs laitiers. \textbf{Levier 2 : Intensification fourragère et alimentation.} Cultures fourragères (Brachiaria, Panicum, Stylosanthes), ensilage, blocs multi-nutritionnels (urée-mélasse) en saison sèche $\rightarrow$ réduction intervalle vêlage-vêlage, engraissement rapide (90--120 j) des jeunes mâles réformés.
\end{enumerate}
\end{exoresolu}

\begin{exercice}[Application : Diagnostic d'une exploitation]
Une exploitation familiale dans le département du Mayo-Rey (Nord) possède 50 bovins (race Goudali), 30 ovins, 20 caprins. L'éleveur pratique la transhumance annuelle vers la plaine de la Bénoué (saison sèche) et revient au village (saison pluies). Il ne vaccinne pas systématiquement. Il souhaite moderniser.
\begin{enumerate}
    \item Caractérisez le système d'exploitation actuel (type, mobilité, races).
    \item Identifiez trois risques majeurs pour cette exploitation.
    \item Proposez un parcours de modernisation progressif en 3 étapes (sanitaire, alimentation, génétique) compatible avec la transhumance.
\end{enumerate}
\end{exercice}

\begin{corrige}[Exercice n°1]
\begin{enumerate}
    \item \textbf{Système :} Élevage traditionnel extensif transhumant. Mobilité saisonnière Nord-Sud (piémont Bénoué en saison sèche, village terroir en saison pluies). Races locales rustiques (Goudali, Sahel pour ovins/caprins). Capital social/épargne.
    \item \textbf{Risques :} 1) \textbf{Santé :} Absence vaccination $\rightarrow$ PPR, Peste bovine (si réintroduction), Fièvre aphteuse, Trypanosomiase (zone de transition). Mortalité jeune élevée. 2) \textbf{Alimentation :} Déficit fourrager saison sèche (paille de riz/sorgho pauvre), perte de poids, avortements. 3) \textbf{Foncier/Conflits :} Blocage couloirs/pâturages par cultures coton/riz, vols, conflits agriculteurs-éleveurs.
    \item \textbf{Modernisation progressive :}
    \begin{itemize}
        \item \textbf{Étape 1 (Sanitaire - Court terme) :} Adhésion au calendrier vaccinal obligatoire (PPR, Fièvre aphteuse, Charbon symptomatique). Traitement antiparasitaire stratégique (vermifuge, trypanocide préventif en zone à risque). Tenue registre troupeau (identification boucles).
        \item \textbf{Étape 2 (Alimentation - Moyen terme) :} Aménagement de \emph{champs fourragers} (Brachiaria ruziziensis, Stylosanthes) sur jachères près du village de saison des pluies. Stockage foin/ensilage pour supplémentation saison sèche (point d'eau transhumance). Blocs urée-mélasse transportables.
        \item \textbf{Étape 3 (Génétique/Valorisation - Long terme) :} Introduction de taureaux améliorateurs (Goudali sélectionnés station Ngaoundaba ou croisés) via IA ou monte naturelle contrôlée. Engraissement ciblé des mâles réformés (120 j) avant vente (poids + valeur). Organisation en groupement (GIC) pour négociation prix, accès crédit, abattoir mobile.
    \end{itemize}
\end{enumerate}
\end{corrige}

\coursec{La pêche au Cameroun}

Après l'agriculture et l'élevage, les activités de production animale au Cameroun se complètent par un troisième pilier : la \cle{pêche}. Comme l'élevage fournit la viande, la pêche fournit le poisson, qui est la principale source de \cle{protéines animales} bon marché pour les populations camerounaises. Mais contrairement à l'élevage, qui produit un animal domestiqué, la pêche exploite une ressource naturelle sauvage, qu'il faut prélever sans l'épuiser. Cette section présente les types de pêche, les grandes zones de production, la pisciculture, les produits et les apports de ce secteur, ainsi que ses limites.

\courssub{Définition et types de pêche}

\begin{definition}[La pêche]
La \cle{pêche} est l'activité qui consiste à capturer les animaux aquatiques (poissons, crustacés comme les crevettes, mollusques) dans leur milieu naturel, en mer ou dans les eaux intérieures, pour l'alimentation ou le commerce.
\end{definition}

On distingue deux grands types de pêche selon le milieu où elle se pratique.

\begin{definition}[Pêche maritime et pêche continentale]
La \cle{pêche maritime} est pratiquée en mer, dans les eaux salées de l'océan. La \cle{pêche continentale} est pratiquée dans les \cle{eaux intérieures} : fleuves, lacs, lagunes et retenues de barrages, par opposition à la pêche maritime.
\end{definition}

La pêche continentale se subdivise elle-même en deux catégories :
\begin{itemize}
\item la \cle{pêche fluviale}, pratiquée dans les fleuves et les rivières (Sanaga, Wouri, Logone, Chari, Nyong, Ntem) ;
\item la \cle{pêche lacustre}, pratiquée dans les lacs (lac Tchad, lac Ossa, lacs de Barombi, retenues des barrages de la Sanaga comme Mbakaou).
\end{itemize}

\begin{attention}[Erreur fréquente]
Beaucoup d'élèves réduisent la pêche camerounaise à la pêche en mer. C'est une erreur : la \cle{pêche continentale} est tout aussi importante. Le \cle{lac Tchad} et le bassin \cle{Logone-Chari} (avec ses plaines d'inondation, les Yaérés) sont des zones de pêche majeures qui nourrissent tout l'Extrême-Nord. Ne jamais oublier ces zones dans un croquis ou une dissertation.
\end{attention}

\courssub{Pêche artisanale et pêche industrielle}

Selon les moyens techniques utilisés, on distingue deux modes d'exploitation.

\textbf{La pêche artisanale} est pratiquée par des pêcheurs individuels ou regroupés en petites équipes, avec des moyens simples :
\begin{itemize}
\item embarcations : \cle{pirogues} en bois, parfois motorisées (moteur hors-bord) ;
\item engins : filets maillants, lignes et hameçons, nasses (pièges en vannerie), sennes tournantes tirées depuis la plage ;
\item acteurs : pêcheurs camerounais, mais aussi de nombreux pêcheurs étrangers installés sur les côtes (Nigérians, Ghanéens notamment) et sur les fleuves du Nord.
\end{itemize}
Elle fournit l'essentiel du poisson consommé localement, vendu frais, fumé ou séché par les femmes (les « mareyeuses »).

\textbf{La pêche industrielle} est pratiquée par des sociétés disposant de navires puissants :
\begin{itemize}
\item embarcations : \cle{chalutiers} équipés de chambres froides ;
\item engins : \cle{chaluts} (grands filets traînés au fond de l'eau), sennes coulissantes ;
\item acteurs : sociétés camerounaises et étrangères opérant au large de Douala, Limbé et Kribi, souvent ciblant les \cle{crevettes} destinées à l'exportation.
\end{itemize}

\begin{popfigure}
\begin{center}
\begin{tikzpicture}[scale=0.9]
% Pêche artisanale
\draw[fill=popblueL] (-6.2,-0.4) rectangle (-0.4,-1.6);
\node[popblue] at (-3.3,-1.9) {\small Mer ou fleuve};
\draw[fill=popgold!60,draw=popdark] (-4.6,0) .. controls (-4.2,0.35) and (-2.4,0.35) .. (-2,0) -- (-2.2,-0.25) -- (-4.4,-0.25) -- cycle;
\node at (-3.3,0.75) {\textbf{Pêche artisanale}};
\draw[popdark] (-2.1,0) -- (-1.3,-0.9);
\draw[popdark,dashed] (-1.3,-0.9) -- (-0.7,-0.9);
\node[right,align=left,font=\small] at (-0.4,-0.9) {filet simple,\\ligne, nasse};
\draw[fill=popgreen!70] (-3.6,-1.0) ellipse (0.28 and 0.1);
\draw[fill=popgreen!70] (-2.9,-1.25) ellipse (0.28 and 0.1);
\node[font=\small] at (-3.3,0.35) {pirogue};
% Pêche industrielle
\draw[fill=popblueL] (1.0,-0.4) rectangle (7.4,-1.6);
\draw[fill=popmuted!50,draw=popdark] (2.4,0) -- (5.6,0) -- (5.2,-0.35) -- (2.8,-0.35) -- cycle;
\draw[fill=popmuted!70,draw=popdark] (3.4,0) rectangle (4.4,0.5);
\draw[popdark] (4.4,0.5) -- (4.4,0.9);
\node at (4.2,1.25) {\textbf{Pêche industrielle}};
\node[font=\small] at (4.0,0.72) {chalutier};
\draw[popdark,thick] (2.8,-0.35) -- (1.6,-1.3);
\draw[popdark,thick] (2.4,-0.35) -- (1.6,-1.3);
\draw[popdark,pattern=north east lines,pattern color=popdark] (1.6,-1.3) -- (2.6,-1.5) -- (2.4,-1.05) -- cycle;
\node[left,font=\small,align=right] at (1.5,-1.35) {chalut traîné\\au fond};
\draw[fill=popgreen!70] (3.4,-1.1) ellipse (0.28 and 0.1);
\draw[fill=popgreen!70] (4.3,-1.35) ellipse (0.28 and 0.1);
\draw[fill=popgreen!70] (5.2,-1.05) ellipse (0.28 and 0.1);
\end{tikzpicture}
\end{center}
\legende{Schéma comparatif : à gauche, la pêche artisanale (pirogue, filets simples, petite équipe) ; à droite, la pêche industrielle (chalutier tirant un grand filet, capture massive).}
\end{popfigure}

\courssub{Les grandes zones de pêche au Cameroun}

Le Cameroun dispose d'un potentiel halieutique (relatif au poisson) varié, réparti entre la mer et les eaux intérieures.

\begin{itemize}
\item \textbf{La côte atlantique} (environ 400 km de littoral) : les principaux ports et centres de pêche sont \cle{Douala} (estuaire du Wouri), \cle{Limbé}, \cle{Kribi} et \cle{Campo}. On y pêche le capitaine, le carpe rouge, le thon, le maquereau, les soles et surtout les \cle{crevettes}.
\item \textbf{Le lac Tchad} et le bassin \cle{Logone-Chari} : zone majeure de pêche continentale dans l'Extrême-Nord. Les plaines d'inondation (Yaérés) sont de véritables nurseries à poissons.
\item \textbf{Les grands fleuves} : la \cle{Sanaga} (et ses barrages), le \cle{Wouri}, le Nyong, le Ntem, le Dibamba.
\item \textbf{Les lagunes et les petits lacs} : lac Ossa, lacs de Barombi (Ouest), lagunes côtières du Sud.
\end{itemize}

\begin{popfigure}
\begin{center}
\begin{tikzpicture}[scale=1.05]
% Contour simplifié du Cameroun
\draw[thick,popdark,fill=popcream] (1.6,0) -- (3.4,0.2) -- (3.8,1.4) -- (4.6,2.2) -- (4.4,3.4) -- (5.0,4.4) -- (4.2,5.6) -- (4.6,6.6) -- (3.4,7.6) -- (2.2,7.2) -- (1.4,6.2) -- (0.6,5.0) -- (0.2,3.6) -- (0.8,2.4) -- (0.4,1.2) -- cycle;
% Mer
\draw[fill=popblueL] (-1.4,-0.6) rectangle (1.6,1.6);
\node[popblue,font=\small] at (-0.4,-0.2) {Océan Atlantique};
% Lac Tchad
\draw[fill=popblue!60,draw=popblue] (3.1,7.9) ellipse (0.75 and 0.45);
\node[font=\small] at (3.1,8.65) {Lac Tchad};
% Fleuves : Logone-Chari (coule vers le NORD vers le lac Tchad)
\draw[popblue,thick] (2.2,5.4) .. controls (2.6,6.2) and (3.2,6.8) .. (3.6,7.6);
\node[font=\scriptsize,popblue,rotate=55] at (3.35,6.5) {Logone-Chari};
% Sanaga (coule vers l'OUEST)
\draw[popblue,thick] (2.6,4.6) .. controls (2.2,3.6) and (1.6,2.6) .. (1.3,1.5);
\node[font=\scriptsize,popblue,rotate=70] at (2.35,3.2) {Sanaga};
% Wouri
\draw[popblue,thick] (1.5,2.2) -- (1.15,0.9);
\node[font=\scriptsize,popblue] at (1.85,1.05) {Wouri};
% Villes côtières
\fill[popink] (1.15,0.75) circle (1.6pt); \node[right,font=\scriptsize] at (1.25,0.75) {Douala};
\fill[popink] (0.55,1.15) circle (1.6pt); \node[left,font=\scriptsize] at (0.45,1.15) {Limbé};
\fill[popink] (1.5,0.25) circle (1.6pt); \node[right,font=\scriptsize] at (1.6,0.25) {Kribi};
\fill[popink] (1.7,-0.05) circle (1.6pt); \node[right,font=\scriptsize] at (1.8,-0.05) {Campo};
% Zone maritime
\draw[poporange,very thick,->] (0.2,0.5) -- (-0.9,0.5);
\node[font=\scriptsize,poporange,align=center] at (-0.5,1.15) {zone de pêche\\maritime};
% Nord
\draw[popdark,->] (5.6,7.4) -- (5.6,8.2); \node[font=\small] at (5.6,8.45) {N};
\end{tikzpicture}
\end{center}
\legende{Croquis schématique des zones de pêche du Cameroun : pêche maritime sur la côte atlantique (Douala, Limbé, Kribi, Campo) ; pêche continentale sur le lac Tchad, le Logone-Chari (coulant vers le nord), la Sanaga et le Wouri. Contour très simplifié.}
\end{popfigure}

\courssub{La pisciculture}

Face à la stagnation des captures naturelles, une solution se développe : produire le poisson au lieu de seulement le capturer.

\begin{definition}[La pisciculture]
La \cle{pisciculture} (ou aquaculture) est l'\cle{élevage de poissons} dans des étangs ou des cages contrôlés par l'homme : on y nourrit, on y protège et on y récolte les poissons comme on élève des poulets ou des porcs.
\end{definition}

Au Cameroun, la pisciculture connaît un développement récent mais prometteur :
\begin{itemize}
\item espèces élevées : le \cle{tilapia}, la \cle{carpe} et le \cle{clarias} (poisson-chat africain, très rustique) ;
\item exemples : fermes piscicoles de la région de l'Ouest (filière de la carpe autour de Foumban et Bafoussam), stations piscicoles publiques, étangs privés dans le Centre et le Littoral, élevage de clarias en bassins dans les zones urbaines ;
\item avantages : production régulière toute l'année, proximité des villes, création d'emplois pour les jeunes.
\end{itemize}

\courssub{Les produits et les apports de la pêche}

\textbf{Les produits.} Le poisson est commercialisé sous trois formes principales : \cle{frais} (près des zones de capture), \cle{fumé} et \cle{séché} (pour être conservé et transporté vers l'intérieur du pays). Les \cle{crevettes}, pêchées par les chalutiers, sont congelées et exportées.

\textbf{Les apports de la pêche} sont considérables :
\begin{itemize}
\item \textbf{alimentation} : le poisson est la source de \cle{protéines animales} la moins chère ; il représente une part essentielle de la ration alimentaire des Camerounais ;
\item \textbf{emplois} : pêcheurs, mareyeuses, fumeurs de poisson, transporteurs, commerçants : des centaines de milliers de personnes vivent de la filière ;
\item \textbf{devises} : l'exportation des crevettes rapporte des devises à l'État.
\end{itemize}

\courssub{Les limites du secteur de la pêche}

Malgré ce potentiel, le secteur rencontre de sérieuses difficultés :
\begin{itemize}
\item la \cle{surexploitation} des ressources : la pêche excessive épuise les stocks, surtout près des côtes ;
\item la \cle{pêche illégale} : des navires étrangers pillent les eaux camerounaises sans autorisation ;
\item la \cle{vétusté du matériel} : les pirogues et les engins artisanaux sont anciens et peu productifs ;
\item les \cle{importations massives} de poisson congelé.
\end{itemize}

\begin{exemplebox}[Un paradoxe chiffré]
Malgré son potentiel (côte atlantique, fleuves, lacs), le Cameroun ne couvre pas ses besoins : il importe chaque année des \cle{centaines de milliers de tonnes} de poisson congelé (sardinelle, maquereau) pour nourrir sa population. La production nationale ne couvre qu'une partie de la demande, ce qui coûte cher en devises.
\end{exemplebox}

\begin{savaistu}[Le rétrécissement du lac Tchad]
En quelques décennies, le lac Tchad est passé d'environ \SI{25000}{km^2} dans les années 1960 à \cle{moins de \SI{2500}{km^2}} aujourd'hui, à cause de la sécheresse, de la baisse des pluies et des prélèvements pour l'irrigation. Ce rétrécissement menace la pêche, l'agriculture et la vie de millions de personnes dans le bassin (Cameroun, Tchad, Niger, Nigeria). C'est l'une des plus graves crises écologiques d'Afrique.
\end{savaistu}

\begin{methode}[Interpréter un tableau de données halieutiques]
\begin{enumerate}
\item \textbf{Lire le titre et les unités} : de quoi parle le tableau (production, importations) ? En tonnes, en milliards de FCFA ?
\item \textbf{Repérer les grandeurs} : identifier les lignes (produits, zones) et les colonnes (années).
\item \textbf{Comparer les valeurs} : calculer les écarts et les taux de variation (en \%).
\item \textbf{Dégager une tendance} : hausse, baisse, stagnation ?
\item \textbf{Commenter} : expliquer la tendance par les connaissances du cours (surexploitation, pisciculture, importations).
\end{enumerate}
\end{methode}

\begin{exoresolu}[Exploiter un tableau de données sur la pêche]
Énoncé : le tableau ci-dessous donne la production nationale de poisson et les importations du Cameroun (données simplifiées, en milliers de tonnes).

\begin{center}
\begin{tabularx}{\linewidth}{|l|X|X|X|}\hline
\rowcolor{popblueL} \textbf{Année} & \textbf{Production nationale} & \textbf{Importations} & \textbf{Besoins totaux} \\ \hline
2000 & 100 & 80 & 180 \\ \hline
2010 & 120 & 160 & 280 \\ \hline
2020 & 130 & 250 & 380 \\ \hline
\end{tabularx}
\end{center}

1. Calcule le taux de variation des importations entre 2000 et 2020. 2. Dégage la tendance générale et commente-la.
\tcblower
Solution détaillée :

1. Taux de variation des importations entre 2000 et 2020 :
\[ \text{Taux} = \frac{250 - 80}{80} \times 100 = \frac{170}{80} \times 100 = 212{,}5\ \%. \]
Les importations ont plus que triplé en vingt ans.

2. Tendance : la production nationale stagne (elle passe de 100 à 130 milliers de tonnes, soit $+30\ \%$ seulement), tandis que les besoins totaux explosent (de 180 à 380 milliers de tonnes) sous l'effet de la croissance démographique. Le déficit est donc comblé par des importations massives de poisson congelé.

Commentaire : ce tableau illustre le paradoxe du secteur de la pêche camerounaise : un potentiel important (mer, fleuves, lac Tchad) mais une production qui stagne à cause de la surexploitation des côtes, de la pêche illégale et de la vétusté du matériel artisanal. Le développement de la pisciculture (tilapia, carpe, clarias) apparaît comme une voie d'avenir pour réduire la dépendance aux importations.
\end{exoresolu}

\begin{exercice}[Pour s'entraîner]
1. Définis : pêche continentale, pisciculture. 2. Cite trois zones de pêche continentale et deux ports de pêche maritime du Cameroun. 3. Explique deux causes de la stagnation de la production nationale de poisson et propose une solution.
\end{exercice}

\begin{corrige}[Exercice]
1. La pêche continentale est la pêche pratiquée dans les eaux intérieures (fleuves, lacs, lagunes), par opposition à la pêche maritime en mer. La pisciculture est l'élevage de poissons en étangs ou en cages contrôlés par l'homme. 2. Zones continentales : le lac Tchad, le bassin Logone-Chari, la Sanaga (ou le Wouri). Ports maritimes : Douala, Limbé (ou Kribi, Campo). 3. Causes : la surexploitation des ressources côtières, la pêche illégale par des navires étrangers, la vétusté du matériel artisanal, le rétrécissement du lac Tchad. Solution : développer la pisciculture (tilapia, carpe, clarias), moderniser les engins des pêcheurs artisans et surveiller les eaux territoriales.
\end{corrige}

\begin{aretenir}[L'essentiel]
La pêche au Cameroun est \cle{maritime} (côte atlantique : Douala, Limbé, Kribi, Campo) et \cle{continentale} (lac Tchad, Logone-Chari, Sanaga, Wouri). Elle est \cle{artisanale} (pirogues, filets simples) ou \cle{industrielle} (chalutiers, crevettes d'exportation). Elle fournit des protéines bon marché, des emplois et des devises, mais elle est limitée par la surexploitation, la pêche illégale et la vétusté du matériel, ce qui oblige le pays à importer massivement du poisson congelé. La \cle{pisciculture} est la voie d'avenir.
\end{aretenir}

\coursec{TD : Problèmes et perspectives du secteur agropastoral et piscicole}

Nous venons d'étudier la pêche au Cameroun, activité qui complète l'agriculture et l'élevage pour nourrir le pays. Ces trois secteurs, regroupés sous le nom de secteur \cle{agropastoral et piscicole}, occupent la majorité de la population active camerounaise et fournissent l'essentiel de notre alimentation. Pourtant, malgré un potentiel exceptionnel (sols fertiles, climats variés, eaux poissonneuses), ce secteur reste fragile : il produit moins qu'il ne pourrait, et le Cameroun importe encore du riz, du poisson congelé et de la farine de blé. Pourquoi un pays si riche en ressources peine-t-il à nourrir pleinement sa population ? C'est toute la problématique de ce travail dirigé : \textbf{identifier les problèmes du secteur agropastoral et piscicole, les classer, puis proposer des solutions réalistes}.

\courssub{Les problèmes du secteur : un diagnostic}

\textbf{1. Les problèmes naturels}\par

L'agriculture, l'élevage et la pêche dépendent directement de la nature. Or celle-ci impose au Cameroun de lourdes contraintes :

\begin{itemize}
\item Les \cle{sécheresses} : dans l'Extrême-Nord et le Nord, la saison des pluies est courte et irrégulière (3 à 5 mois). Les années de déficit pluviométrique provoquent des pertes de récoltes de mil et de sorgho et la mort de nombreux bovins, comme lors des grandes sécheresses sahéliennes des années 1970 et 1980.
\item L'\cle{érosion des sols} : sur les hautes terres de l'Ouest et du Nord-Ouest, les pluies violentes sur les pentes déboisées emportent la couche fertile. La culture sur brûlis et le raccourcissement des jachères épuisent aussi les sols.
\item Les \cle{parasites et maladies} : la mouche \cle{tsé-tsé}, qui transmet la trypanosomiase, empêche l'élevage bovin dans de vastes zones forestières du Sud ; les chenilles et les sauteriaux détruisent les cultures vivrières ; la maladie du \emph{swollen shoot} attaque les cacaoyers.
\item Les \cle{aléas climatiques} : pluies tardives ou diluviennes, inondations des plaines du Logone, vents violents détruisant les bananeraies du littoral.
\end{itemize}

\textbf{2. Les problèmes techniques}\par

Le paysan camerounais travaille encore souvent comme il y a cinquante ans :

\begin{itemize}
\item L'\cle{outillage rudimentaire} : la houe, la machette et la daba dominent ; la \cle{mécanisation} (tracteurs, motoculteurs) concerne moins d'une exploitation sur dix.
\item Le faible usage des \cle{intrants} : engrais, pesticides et semences améliorées coûtent cher et sont peu distribués. Résultat : les rendements restent faibles, souvent autour de 1 tonne de maïs à l'hectare, alors que des variétés améliorées bien conduites en donnent 3 à 4.
\item Les techniques d'élevage restent extensives : le zébu peulh, bien adapté mais peu productif, donne peu de lait et de viande comparé aux races améliorées.
\end{itemize}

\textbf{3. Les problèmes économiques}\par

\begin{itemize}
\item La \cle{faiblesse des prix} payés aux producteurs : le planteur de cacao ou de café reçoit une faible part du prix final, le reste étant capté par les intermédiaires et les coûts d'exportation.
\item Le manque de \cle{crédit agricole} : les banques prêtent peu aux paysans, jugés trop risqués ; sans argent, impossible d'acheter engrais, équipements ou alevins.
\item La \cle{concurrence des importations} : le riz asiatique, le poulet congelé et le poisson importé inondent les marchés urbains à des prix que le producteur local peine à égaler.
\item La baisse et l'instabilité des \cle{cours mondiaux} du cacao et du café : quand les cours s'effondrent, comme dans les années 1990, les revenus de centaines de milliers de planteurs s'effondrent avec eux.
\end{itemize}

\textbf{4. Les problèmes sociaux et fonciers}\par

\begin{itemize}
\item Le \cle{vieillissement} et l'\cle{exode rural} : les jeunes quittent les villages pour Douala et Yaoundé ; la moyenne d'âge des exploitants dépasse souvent 50 ans, ce qui menace le renouvellement de la main-d'œuvre agricole.
\item Les \cle{conflits agriculteurs-éleveurs} : dans l'Adamaoua, le Nord-Ouest et l'Extrême-Nord, les troupeaux qui piétinent les champs provoquent des affrontements parfois sanglants entre cultivateurs et pasteurs.
\item L'\cle{insécurité foncière} : beaucoup de paysans n'ont pas de titre foncier ; leurs terres peuvent être récupérées par l'État ou des investisseurs, ce qui décourage les investissements durables (plantations, arbres fruitiers).
\end{itemize}

\textbf{5. Les problèmes d'infrastructures}\par

\begin{itemize}
\item Le mauvais état des \cle{pistes rurales} : en saison des pluies, nombre de villages du Centre, de l'Est ou du Nord-Ouest sont enclavés ; les plantains et tomates pourrissent au bord des champs faute de pouvoir rejoindre les villes.
\item Le déficit de \cle{stockage et de conservation} : peu de magasins, de silos, de chambres froides ou de séchoirs. Les \cle{pertes post-récolte} sont considérables.
\item La faible \cle{transformation locale} : le Cameroun exporte le cacao brut et importe du chocolat ; il produit du manioc mais peu de gari ou d'amidon industrialisé. La valeur ajoutée échappe au pays.
\end{itemize}

\begin{exemplebox}[Un exemple chiffré : les pertes post-récolte]
On estime qu'environ \textbf{30 \% à 40 \%} des récoltes vivrières du Cameroun (fruits, légumes, tubercules, poisson frais) sont perdues après la récolte, faute de routes praticables, de moyens de conservation (froid, séchage, fumage amélioré) et d'unités de transformation. Concrètement : sur 10 tonnes de tomates produites dans l'Ouest, 3 à 4 tonnes n'arrivent jamais au consommateur. Ces pertes réduisent les revenus des producteurs, renchérissent les prix en ville et aggravent l'\cle{insécurité alimentaire}, alors même que les champs ont bien produit.
\end{exemplebox}

\courssub{Construire l'arbre des problèmes}

Pour organiser ce diagnostic, l'outil le plus efficace est l'\cle{arbre des problèmes} : il relie les causes au problème central, puis le problème central à ses conséquences.

\begin{methode}[Construire un arbre des problèmes]
\begin{enumerate}
\item Formuler le \textbf{problème central} en une phrase claire (le « tronc ») : ici, « le secteur agropastoral et piscicole camerounais est peu productif et ne couvre pas les besoins du pays ».
\item Classer les \textbf{causes} par catégories (les « racines ») : naturelles, techniques, économiques, sociales et foncières, infrastructures. Chaque grande cause peut être décomposée en causes secondaires.
\item Identifier les \textbf{conséquences} (les « branches ») : importations alimentaires coûteuses, insécurité alimentaire, pauvreté rurale, exode rural.
\item Vérifier la logique : chaque racine doit pouvoir se lire « parce que... » vers le tronc, et chaque branche « donc... » à partir du tronc.
\end{enumerate}
\end{methode}

\begin{popfigure}
\begin{tikzpicture}[
  font=\small,
  cause/.style={draw=popblue, fill=popblueL, rounded corners=2pt, text width=3.1cm, align=center, inner sep=2pt},
  centre/.style={draw=popdark, fill=poporangeL, rounded corners=3pt, text width=4.6cm, align=center, font=\small\bfseries, inner sep=3pt},
  cons/.style={draw=poppurple, fill=poppurpleL, rounded corners=2pt, text width=3.1cm, align=center, inner sep=2pt},
  lien/.style={-{Stealth[length=2mm]}, thick, popmuted}
]
\node[centre] (P) at (0,0) {Problème central : faible productivité du secteur agropastoral et piscicole};
\node[cause] (c1) at (-5.6,-3) {Naturelles : sécheresses (Nord), érosion, tsé-tsé, chenilles};
\node[cause] (c2) at (-2.8,-3.6) {Techniques : houe, peu d'intrants, faible mécanisation};
\node[cause] (c3) at (0,-4) {Économiques : prix faibles, pas de crédit, importations, cours mondiaux};
\node[cause] (c4) at (2.8,-3.6) {Sociales et foncières : exode rural, conflits, insécurité foncière};
\node[cause] (c5) at (5.6,-3) {Infrastructures : pistes dégradées, pertes post-récolte, peu de transformation};
\node[cons] (s1) at (-4,3) {Importations alimentaires coûteuses (riz, poisson, blé)};
\node[cons] (s2) at (0,3.6) {Insécurité alimentaire et hausse des prix en ville};
\node[cons] (s3) at (4,3) {Pauvreté rurale et exode des jeunes};
\draw[lien] (c1.north) -- (P.south west);
\draw[lien] (c2.north) -- (P.south);
\draw[lien] (c3.north) -- (P.south);
\draw[lien] (c4.north) -- (P.south);
\draw[lien] (c5.north) -- (P.south east);
\draw[lien] (P.north west) -- (s1.south);
\draw[lien] (P.north) -- (s2.south);
\draw[lien] (P.north east) -- (s3.south);
\end{tikzpicture}
\legende{Arbre des problèmes du secteur agropastoral et piscicole camerounais : en bas, les causes classées par catégories (les racines) ; au centre, le problème central (le tronc) ; en haut, les conséquences (les branches).}
\end{popfigure}

\begin{attention}[Erreur fréquente à l'examen]
Beaucoup de candidats se contentent d'\textbf{énumérer des problèmes} sans les classer ni les illustrer. Or l'examinateur attend : (1) un \textbf{classement} par catégories (naturels, techniques, économiques, sociaux, infrastructures) ; (2) des \textbf{exemples camerounais précis et localisés} : dire « la sécheresse dans l'Extrême-Nord », « la mouche tsé-tsé en zone forestière du Sud », « les conflits agriculteurs-éleveurs sur les hautes terres du Nord-Ouest », et non « il y a des problèmes climatiques ». Un argument non localisé est un argument à moitié perdu.
\end{attention}

\courssub{Les solutions et perspectives}

À chaque catégorie de problèmes correspondent des solutions, que l'État, les collectivités, les organisations de producteurs et les partenaires extérieurs peuvent mettre en œuvre.

\begin{itemize}
\item \textbf{Face aux contraintes naturelles} : aménagements hydro-agricoles (retenues d'eau, périmètres irrigués comme ceux de la SEMRY dans la plaine du Logone), lutte anti-érosive (terrasses, reboisement des pentes), campagnes de traitement contre la tsé-tsé et les parasites, choix de variétés résistantes à la sécheresse.
\item \textbf{Face aux problèmes techniques} : \cle{vulgarisation agricole} (encadrement et conseil aux paysans par des techniciens), \cle{subvention des intrants} (engrais et semences améliorées à prix réduit), location de matériel, modernisation de l'élevage (races améliorées, parcs de vaccination).
\item \textbf{Face aux problèmes économiques} : \cle{crédit rural} adapté (microfinance, banques agricoles), prix planchers et stabilisation des revenus des planteurs, protection raisonnable contre les importations déloyales, soutien à la \cle{pisciculture} continentale pour réduire les importations de poisson.
\item \textbf{Face aux problèmes sociaux et fonciers} : sécurisation du foncier rural (titres fonciers, cadastre), délimitation de couloirs de transhumance et médiation entre agriculteurs et éleveurs, incitation des jeunes à s'installer (installation aidée, formation).
\item \textbf{Face aux problèmes d'infrastructures} : aménagement et entretien des \cle{pistes rurales}, construction de magasins de stockage, de séchoirs et de chambres froides, création d'\cle{agro-industries} de transformation (huileries, minoteries, unités de transformation du manioc, conserveries de poisson).
\end{itemize}

\begin{center}
\begin{tabularx}{\linewidth}{|l|X|X|}
\hline
\rowcolor{popblueL} \textbf{Catégorie} & \textbf{Problèmes} & \textbf{Solutions possibles} \\
\hline
Naturels & Sécheresses (Nord), érosion, tsé-tsé, chenilles & Irrigation, reboisement, traitements, variétés résistantes \\
\hline
Techniques & Outillage rudimentaire, peu d'intrants & Vulgarisation, subvention des intrants, mécanisation \\
\hline
Économiques & Prix faibles, pas de crédit, importations & Crédit rural, prix planchers, soutien à la production locale \\
\hline
Sociaux et fonciers & Exode rural, conflits, insécurité foncière & Titres fonciers, couloirs de transhumance, installation des jeunes \\
\hline
Infrastructures & Pistes dégradées, pertes post-récolte & Aménagement des pistes, stockage, agro-industries \\
\hline
\end{tabularx}
\end{center}

\begin{savaistu}[L'agriculture de deuxième génération]
Pour moderniser le secteur, le Cameroun a lancé des programmes d'« \cle{agriculture de deuxième génération} » et des \cle{agropôles} : des zones aménagées (eau, pistes, électricité) où de jeunes agropreneurs cultivent de façon mécanisée et intensive, avec accès au crédit et à la transformation. L'objectif est double : augmenter la production et rendre le métier agricole attractif pour la jeunesse. C'est un excellent exemple d'actualité à mobiliser dans une conclusion de dissertation au Probatoire ou au Baccalauréat technique.
\end{savaistu}

\courssub{Rédaction guidée d'une conclusion argumentée}

\begin{methode}[Rédiger et justifier une conclusion]
\begin{enumerate}
\item \textbf{Rappeler la problématique} en la reformulant brièvement : ne recopiez pas le sujet mot à mot.
\item \textbf{Synthétiser les idées clés} du développement en deux ou trois phrases : ici, la richesse du potentiel, la diversité des problèmes, l'existence de solutions.
\item \textbf{Ouvrir sur une perspective} : une évolution souhaitable, un défi à venir (souveraineté alimentaire, modernisation).
\item \textbf{Ne jamais introduire d'idée nouvelle} dans la conclusion : tout ce qui s'y trouve doit avoir été démontré avant.
\end{enumerate}
\end{methode}

\begin{exoresolu}[Exercice de synthèse : rédiger une conclusion]
Énoncé. À partir du diagnostic établi dans ce TD, rédigez une conclusion argumentée (6 à 10 lignes) répondant à la problématique : « Le secteur agropastoral et piscicole peut-il faire du Cameroun un pays souverain sur le plan alimentaire ? »
\tcblower
Solution détaillée (modèle de réponse). \par
« Ainsi, le secteur agropastoral et piscicole camerounais dispose d'un potentiel remarquable : des sols fertiles, des climats variés, un cheptel important et des eaux poissonneuses. Pourtant, nous l'avons montré, ce secteur reste freiné par des contraintes naturelles (sécheresses du Nord, érosion, parasites), techniques (outillage rudimentaire, faible usage des intrants), économiques (prix faibles, manque de crédit, concurrence des importations), sociales (exode rural, conflits agriculteurs-éleveurs) et d'infrastructures (pistes dégradées, pertes post-récolte de l'ordre de 30 à 40 \%). Ces obstacles expliquent que le Cameroun importe encore une partie de sa nourriture. Mais les solutions existent : vulgarisation agricole, crédit rural, aménagement des pistes, pisciculture moderne, agropôles et agro-industries de transformation. Si ces politiques sont poursuivies et amplifiées, le secteur agropastoral et piscicole pourra devenir le véritable pilier de la \cle{souveraineté alimentaire} du Cameroun, c'est-à-dire de sa capacité à nourrir durablement sa population par ses propres moyens, tout en créant des emplois pour sa jeunesse. »
\end{exoresolu}

\begin{exercice}[À vous de jouer]
\begin{enumerate}
\item Classez les problèmes suivants dans la bonne catégorie : mouche tsé-tsé ; absence de chambres froides ; exode des jeunes ; houe et machette ; baisse des cours du cacao ; conflits autour des champs piétinés.
\item Complétez le tableau problèmes/solutions avec deux lignes supplémentaires de votre choix, en citant un exemple camerounais localisé pour chaque problème.
\item Rédigez une conclusion de 6 à 10 lignes sur le thème : « Moderniser l'agriculture camerounaise : une nécessité pour la jeunesse ».
\end{enumerate}
\end{exercice}

\begin{corrige}[Exercice « À vous de jouer »]
\begin{enumerate}
\item Naturel : mouche tsé-tsé. Infrastructures : absence de chambres froides. Social : exode des jeunes ; conflits autour des champs piétinés. Technique : houe et machette. Économique : baisse des cours du cacao.
\item Exemple attendu : problème « sécheresse dans l'Extrême-Nord » $\rightarrow$ solution « irrigation et retenues d'eau (SEMRY) » ; problème « importations de poisson » $\rightarrow$ solution « développement de la pisciculture continentale ». Toute ligne cohérente, classée et localisée est acceptée.
\item La conclusion doit rappeler la problématique, synthétiser (potentiel, problèmes, solutions) et ouvrir sur une perspective (emploi des jeunes, souveraineté alimentaire), sans idée nouvelle.
\end{enumerate}
\end{corrige}

\begin{aretenir}[L'essentiel du TD]
Le secteur agropastoral et piscicole camerounais souffre de problèmes \textbf{naturels} (sécheresses, érosion, parasites), \textbf{techniques} (outillage rudimentaire, peu d'intrants), \textbf{économiques} (prix faibles, manque de crédit, importations), \textbf{sociaux et fonciers} (exode rural, conflits, insécurité foncière) et d'\textbf{infrastructures} (pistes, stockage, transformation). Les solutions passent par la vulgarisation, la subvention des intrants, le crédit rural, l'aménagement des pistes, la pisciculture moderne et les agro-industries. Enjeu majeur : faire de ce secteur le pilier de la \cle{souveraineté alimentaire} du Cameroun.
\end{aretenir}

\newpage

\coursec{Activité d'intégration}

\courssub{Objectifs de l'activité}

À l'issue de ce travail pratique, l'élève doit être capable de :
\begin{itemize}
\item \cle{localiser} sur un fond de carte vierge les grandes zones de production agricole, d'élevage et de pêche du Cameroun ;
\item \cle{construire} une légende simple et un diagramme en bâtons à partir de données chiffrées ;
\item \cle{exploiter} un dossier documentaire (cartes, photographies, tableaux statistiques, article de presse) pour établir un diagnostic ;
\item \cle{rédiger et justifier} une conclusion argumentée répondant à une question-problème ;
\item \cle{appliquer} les notions de la leçon (agriculture, élevage, pêche, problèmes du secteur) à une situation concrète de la vie courante.
\end{itemize}

\courssub{Situation}

En mars 2026, le conseil régional de l'Adamaoua organise à Ngaoundéré un forum intitulé \emph{« Nourrir le Cameroun : mythe ou réalité ? »}. Ta classe de Première technique est invitée à présenter un stand. Le proviseur remet à ton groupe un dossier documentaire et te demande de produire une \cle{carte de synthèse} des activités agropastorales et piscicoles du pays, un \cle{diagramme} des productions vivrières et un \cle{diagnostic écrit} du secteur.

\textbf{Document 1 — Les zones agro-écologiques du Cameroun (résumé).} Le Cameroun compte cinq grandes zones agro-écologiques : la zone soudano-sahélienne (Extrême-Nord et Nord : mil, sorgho, coton, élevage), la zone de hautes savanes (Adamaoua : élevage bovin, maïs), la zone de hauts plateaux (Ouest et Nord-Ouest : cultures vivrières, café arabica, maraîchage), la zone forestière monomodale (Littoral et Sud-Ouest : bananes, palmiers à huile, hévéa) et la zone forestière bimodale (Centre, Sud, Est : cacao, café robusta, cultures vivrières).

\textbf{Document 2 — Photographies.} Photo A : une plantation de cacao à Ngoumou (région du Centre). Photo B : un troupeau de zébus conduit par un berger peul sur les pâturages de l'Adamaoua. Photo C : des pêcheurs artisans remontant leurs filets sur la plage de Kribi (région du Sud).

\textbf{Document 3 — Productions vivrières du Cameroun en 2022 (milliers de tonnes, sources : MINADER/FAO, valeurs arrondies).}

\begin{center}
\begin{tabularx}{\linewidth}{|l|>{\raggedleft\arraybackslash}X|}\hline
\rowcolor{popblueL} \textbf{Produit vivrier} & \textbf{Production (milliers de tonnes)} \\ \hline
Maïs & 2 300 \\ \hline
Manioc & 5 600 \\ \hline
Banane plantain & 4 700 \\ \hline
Mil/sorgho & 1 300 \\ \hline
Igname & 700 \\ \hline
Riz paddy & 400 \\ \hline
\end{tabularx}
\end{center}

\textbf{Document 4 — Quelques repères sur l'élevage et la pêche.} Le cheptel camerounais compte environ 7 millions de bovins (concentrés à l'Adamaoua, au Nord et à l'Extrême-Nord), 8 millions de petits ruminants et plus de 50 millions de volailles. La production nationale de poisson est estimée à environ 330 000 tonnes par an (pêche artisanale maritime autour de Kribi, Limbe et Douala ; pêche continentale sur le Logone, la Bénoué et le lac Tchad ; aquaculture encore modeste). Or la demande nationale dépasse 500 000 tonnes : le Cameroun importe chaque année près de 200 000 tonnes de poisson congelé, pour une facture supérieure à 100 milliards de FCFA.

\textbf{Document 5 — Extrait de presse (\emph{Cameroun Tribune}, 2025).} « Malgré un potentiel agricole considérable, le Cameroun consacre chaque année des centaines de milliards de FCFA à l'importation de riz, de blé, de poisson et de produits laitiers. Les pistes rurales dégradées, la faible mécanisation, le vieillissement des planteurs, les conflits agriculteurs-éleveurs et les aléas climatiques freinent la modernisation du secteur. »

\textbf{Question-problème :} le secteur agropastoral et piscicole camerounais est-il « un géant aux pieds d'argile » ?

\courssub{Consignes}

\begin{enumerate}
\item Sur le fond de carte vierge du Cameroun fourni, place et légende : deux cultures de rente (cacao, café ou banane), deux cultures vivrières, la principale zone d'élevage bovin et deux zones de pêche (une maritime, une continentale). N'oublie ni le titre, ni l'orientation, ni la légende.
\item À partir du document 3, construis un diagramme en bâtons des productions vivrières (échelle : 1 cm pour 1 000 milliers de tonnes). Calcule la part du manioc dans le total des six productions.
\item À partir des documents 2, 4 et 5, complète le tableau « problèmes / solutions » en proposant pour chaque problème identifié une solution réaliste.
\item Rédige individuellement une conclusion d'environ quinze lignes répondant à la question-problème : « Le secteur agropastoral et piscicole camerounais : un géant aux pieds d'argile ? ». Ta réponse doit présenter d'abord la puissance du secteur, puis ses faiblesses, et se terminer par une position justifiée.
\end{enumerate}

\courssub{Ressources à mobiliser}

\begin{aretenir}[Notions utiles de la leçon]
\begin{itemize}
\item \cle{Agriculture} : distinction entre cultures vivrières (manioc, maïs, plantain...) et cultures de rente ou d'exportation (cacao, café, coton, banane, palmier à huile) ; systèmes de culture (extensif sur brûlis, plantations modernes).
\item \cle{Élevage} : élevage extensif traditionnel (zébus de l'Adamaoua et du Nord) et élevage moderne ou semi-intensif (volailles, porcs autour des villes).
\item \cle{Pêche} : pêche artisanale maritime (Kribi, Limbe, Douala), pêche continentale (fleuves, lac Tchad) et aquaculture.
\item \cle{Problèmes du secteur} : faible productivité, enclavement, aléas climatiques, conflits agriculteurs-éleveurs, importations alimentaires coûteuses.
\item Savoir-faire statistique : part en pourcentage $= \dfrac{\text{partie}}{\text{total}} \times 100$ ; construction d'un diagramme en bâtons ; règles d'un croquis géographique (titre, légende, orientation).
\end{itemize}
\end{aretenir}

\courssub{Démarche et corrigé commenté}

\begin{exoresolu}[Corrigé commenté du TP]
\textbf{Étape 1 — La carte de synthèse.} Sur le fond de carte vierge, on reporte : le \cle{cacao} dans le Centre, le Sud et le Littoral (zone forestière bimodale) ; le \cle{café} dans l'Ouest (arabica) et le Centre-Sud (robusta) ; le \cle{coton} dans le Nord et l'Extrême-Nord ; les \cle{cultures vivrières} un peu partout mais surtout dans l'Ouest et le Centre ; l'\cle{élevage bovin} sur l'Adamaoua, le Nord et l'Extrême-Nord ; la \cle{pêche maritime} sur la côte (Kribi, Limbe, Douala) et la \cle{pêche continentale} sur le Logone, la Bénoué et le lac Tchad. Le croquis doit comporter un titre (« Les activités agropastorales et piscicoles du Cameroun »), une flèche du nord et une légende numérotée.

\begin{popfigure}
\begin{tikzpicture}[scale=0.9]
% Contour simplifié du Cameroun
\draw[thick,popdark,fill=popcream] (0,0) -- (1.6,0.2) -- (2.4,1.4) -- (2.2,2.8) -- (2.8,4.2) -- (3.6,5.4) -- (4.4,6.6) -- (4.0,7.8) -- (3.0,8.6) -- (2.2,7.6) -- (1.4,6.4) -- (0.6,5.0) -- (-0.6,3.6) -- (-0.8,2.0) -- cycle;
% Zones agro-écologiques (aplats de couleur)
\fill[popgreenL,opacity=0.5] (-0.4,1.2) rectangle (2.2,2.9); % Zone forestière bimodale / monomodale
\fill[popgoldL,opacity=0.5] (0.8,4.4) rectangle (3.4,6.4);   % Adamaoua
\fill[poporangeL,opacity=0.5] (1.6,6.4) rectangle (4.0,8.4);  % Soudano-sahélienne
% Villes repères
\fill[popink] (0.2,0.9) circle (0.07) node[right,font=\scriptsize]{Douala};
\fill[popink] (0.5,1.5) circle (0.07) node[right,font=\scriptsize]{Yaoundé};
\fill[popink] (-0.5,0.4) circle (0.07) node[right,font=\scriptsize]{Kribi};
\fill[popink] (2.1,5.4) circle (0.07) node[right,font=\scriptsize]{Ngaoundéré};
\fill[popink] (3.2,7.6) circle (0.07) node[right,font=\scriptsize]{Maroua};
% Légendes des zones (texte)
\node[font=\scriptsize\bfseries,popdark,align=center] at (0.9,2.2) {Zone forestière\\Cacao, café, vivrier};
\node[font=\scriptsize\bfseries,popdark,align=center] at (2.1,5.0) {Hautes savanes\\Élevage bovin};
\node[font=\scriptsize\bfseries,popdark,align=center] at (2.8,7.3) {Soudano-sahélienne\\Coton, mil, élevage};
% Flèches pêche
\draw[popblue,thick,->] (-0.9,0.6) -- (-1.6,0.2) node[left,font=\scriptsize,align=center]{Pêche\\maritime};
\draw[popblue,thick,->] (3.6,8.2) -- (4.6,8.6) node[right,font=\scriptsize,align=center]{Pêche continentale\\Lac Tchad, Logone};
% Nord
\draw[->,thick,popdark] (4.6,6.0) -- (4.6,7.0) node[above,font=\scriptsize]{N};
\end{tikzpicture}
\legende{Croquis de synthèse attendu : 1. zone forestière (cacao, café, cultures vivrières) ; 2. hautes savanes de l'Adamaoua (élevage bovin) ; 3. zone soudano-sahélienne (coton, mil/sorgho, élevage) ; 4. zones de pêche maritime (côte) et continentale (lac Tchad, fleuves). Contour volontairement simplifié.}
\end{popfigure}

\textbf{Étape 2 — Le diagramme en bâtons.} Total des six productions : $2 300 + 5 600 + 4 700 + 1 300 + 700 + 400 = 15 000$ milliers de tonnes. Part du manioc :
\[ \frac{5 600}{15 000} \times 100 \approx 37,3\,\%. \]
Le manioc est donc la première production vivrière du pays, devant le plantain (environ 31,3\,\%) et le maïs (environ 15,3\,\%). Le diagramme comporte six bâtons de hauteurs proportionnelles (5,6 cm pour le manioc, 4,7 cm pour le plantain, 2,3 cm pour le maïs, 1,3 cm pour le mil/sorgho, 0,7 cm pour l'igname, 0,4 cm pour le riz).

\begin{popfigure}
\begin{tikzpicture}
\begin{axis}[
    ybar,
    bar width=0.6cm,
    width=0.95\linewidth,
    height=6cm,
    ylabel={Production (milliers de tonnes)},
    symbolic x coords={Maïs,Manioc,Plantain,Mil/Sorgho,Igname,Riz},
    xtick=data,
    ymin=0,
    ymax=6200,
    nodes near coords,
    nodes near coords style={font=\scriptsize, rotate=90, anchor=west, yshift=2pt},
    y label style={font=\small},
    xticklabel style={font=\scriptsize, rotate=30, anchor=east},
    enlarge x limits=0.15,
    major x tick style={draw=none}
]
\addplot[fill=popblue, draw=popdark] coordinates {
    (Maïs,2300)
    (Manioc,5600)
    (Plantain,4700)
    (Mil/Sorgho,1300)
    (Igname,700)
    (Riz,400)
};
\end{axis}
\end{tikzpicture}
\legende{Diagramme en bâtons des productions vivrières du Cameroun en 2022 (sources : MINADER/FAO, valeurs arrondies). Le manioc domine largement (37,3\,\% du total).}
\end{popfigure}

\textbf{Étape 3 — Le tableau problèmes / solutions.}

\begin{center}
\begin{tabularx}{\linewidth}{|X|X|}\hline
\rowcolor{popblueL} \textbf{Problèmes identifiés (doc. 4, 5)} & \textbf{Solutions possibles réalistes} \\ \hline
Faible mécanisation et vieillissement des planteurs & Former et installer les jeunes agriculteurs ; subventionner l'acquisition de matériel (tracteurs, décortiqueuses) via des coopératives \\ \hline
Pistes rurales dégradées (enclavement des zones de production) & Programme prioritaire de réhabilitation des routes rurales ; entretien du réseau ferré (Transcamerounais) pour l'évacuation des récoltes \\ \hline
Aléas climatiques (sécheresses au Nord, pluies irrégulières au Sud) & Développer l'irrigation (barrages, pompage) ; diffuser des variétés à cycle court et résistantes à la sécheresse ; mettre en place une assurance indicielle agricole \\ \hline
Conflits agriculteurs-éleveurs (divagation, destruction des champs) & Délimiter et baliser les couloirs de transhumance ; créer des réserves de pâturage et des points d'eau ; commissions locales de concertation \\ \hline
Importations massives de poisson (200 000 t/an) et de riz & Soutenir l'aquaculture (alevins, aliments, formation) ; relancer la riziculture pluviale et irriguée (SEMRY, UNVDA) ; taxer modérément les importations pour financer la filière locale \\ \hline
\end{tabularx}
\end{center}

\textbf{Étape 4 — La conclusion rédigée (modèle).} Le Cameroun possède un secteur agropastoral et piscicole puissant : cinq zones agro-écologiques complémentaires, plus de 15 millions de tonnes de produits vivriers produits en 2022, un cheptel de 7 millions de bovins, des cultures de rente réputées (cacao, café, banane, coton, hévéa) et 330 000 tonnes de poisson capturées ou élevées par an. C'est un « géant » qui nourrit l'essentiel de la population, fait vivre la majorité des ruraux et génère des devises. Mais ce géant a des « pieds d'argile » : faible productivité à l'hectare et par tête, mécanisation quasi inexistante en milieu paysan, pistes rurales dégradées qui enclavent les bassins de production, aléas climatiques récurrents, conflits fonciers entre agriculteurs et éleveurs. Résultat : le pays importe près de 200 000 tonnes de poisson et des quantités croissantes de riz, de blé et de lait, pour une facture annuelle de plusieurs centaines de milliards de FCFA. Le secteur est donc un géant potentiel encore fragile : pour consolider ses pieds, l'État et les producteurs doivent moderniser l'agriculture familiale, développer l'aquaculture et la riziculture, et désenclaver durablement les campagnes.
\tcblower
\textbf{Commentaires pour la correction.} (1) La carte est évaluée sur l'exactitude des localisations (régions, non villes seules) et la présence obligatoire du titre, de la légende organisée et de l'orientation (flèche Nord). (2) Le calcul de la part du manioc doit faire apparaître la formule du pourcentage ; l'erreur fréquente est d'oublier le facteur 100 ou de diviser par 100 au lieu de multiplier. (3) Le tableau exige une solution \emph{spécifique} et \emph{réaliste} pour chaque problème cité, pas une phrase vague. (4) La conclusion est notée sur la structure en trois temps (puissance / faiblesses / position nuancée) et sur l'exploitation \emph{chiffrée} du dossier : une copie sans aucune donnée numérique ne peut obtenir la note maximale.
\end{exoresolu}

\begin{attention}[Pièges fréquents]
Ne confonds pas \cle{cultures vivrières} et \cle{cultures de rente} sur la carte : le cacao et le café sont des cultures de rente (exportation), le manioc et le maïs sont vivriers (autoconsommation/marché local). N'oublie jamais la légende ni le titre du croquis : au Probatoire et au Baccalauréat technique, un croquis sans légende perd des points systématiquement. Dans le diagramme, vérifie que les hauteurs des bâtons sont strictement proportionnelles aux valeurs (échelle respectée) et que les axes sont légendés avec l'unité.
\end{attention}

\courssub{Bilan de l'activité}

Ce TP a permis de mettre en évidence trois idées essentielles de la leçon. D'abord, la \cle{diversité} des activités agropastorales et piscicoles du Cameroun, qui s'explique par la variété des zones agro-écologiques : chaque région a sa spécialité (cacao au Sud, coton au Nord, élevage à l'Adamaoua, pêche sur la côte et les fleuves). Ensuite, le \cle{poids économique et social} du secteur : il assure l'essentiel de la sécurité alimentaire, emploie plus de 60\,\% de la population active et fournit les principales exportations agricoles. Enfin, la \cle{fragilité} de ce secteur, résumée par l'image du « géant aux pieds d'argile » : malgré un potentiel considérable, le Cameroun reste structurellement dépendant des importations alimentaires. En construisant une carte, un diagramme et un diagnostic argumenté, tu as mobilisé à la fois les savoirs de la leçon (agriculture, élevage, pêche, problèmes du secteur) et les savoir-faire exigibles : appliquer les notions à une situation concrète et rédiger une réponse justifiée, deux compétences directement évaluées au Probatoire et au Baccalauréat technique.

\newpage

\coursec{Méthodes et exercices résolus}

\coursec{Les activités agropastorales}

\courssub{Généralités sur les activités agropastorales et piscicoles}

\begin{definition}[Secteur agropastoral et piscicole]
Ensemble des activités économiques liées à la production primaire d'origine biologique : \textbf{agriculture} (cultures vivrières, industrielles, pérennes), \textbf{élevage} (bovin, ovin, caprin, porcin, avicole, cunicole) et \textbf{pêche} (maritime artisanale/industrielle, continentale, aquaculture). Ce secteur constitue le socle de la sécurité alimentaire et un vivier d'emplois majeur.
\end{definition}

\begin{aretenir}[Place dans l'économie camerounaise (Données récentes INS/Minader/Minepia)]
\begin{itemize}
    \item \textbf{Emploi :} $\approx$ \textbf{60\%} de la population active (majorité informelle, exploitation familiale).
    \item \textbf{PIB :} $\approx$ \textbf{15--17\%} (part en baisse relative face aux services/industries, mais volume absolu en hausse).
    \item \textbf{Balance commerciale :} Déficit structurel alimentaire (importations massives : riz, blé, poisson surgelé, huile, lait, viande) $\approx$ \textbf{1\,500--2\,000 milliards FCFA/an}.
    \item \textbf{Exportations :} Cacao, coton, café, banane, hévéa, bois, ananas (produits bruts ou peu transformés).
\end{itemize}
\end{aretenir}

\begin{propriete}[Facteurs de localisation des activités agropastorales]
La répartition spatiale résulte de la combinaison de :
\begin{enumerate}
    \item \textbf{Facteurs naturels (Potentiel) :} Climat (pluviométrie, température, saisonnalité) $\to$ Zones agro-écologiques ; Sols (ferralitiques, vertisols, alluviaux) ; Ressources en eau (fleuves, nappes, lacs, barrages) ; Ressources halieutiques (plateau continental, estuaires, cours d'eau).
    \item \textbf{Facteurs humains (Mise en valeur) :} Densité de population / Main-d'œuvre ; Niveau technique / Vulgarisation ; Politique foncière (droit coutumier vs droit écrit) ; Infrastructures (routes, stockage, transformation, énergie) ; Marchés / Demande urbaine (Yaoundé, Douala) ; Politiques publiques (PNDA, PIDMA, SND30, ZAAP).
\end{enumerate}
\end{propriete}

\begin{exemplebox}[Zonage agro-écologique majeur du Cameroun (Repères pour croquis)]
\begin{center}
\begin{tabularx}{\linewidth}{|l|X|X|}\hline
\rowcolor{popblueL} \textbf{Zone} & \textbf{Régions concernées} & \textbf{Spécifications dominantes} \\ \hline
\textbf{Équatoriale} (Bimodale) & Sud, Littoral, Sud-Ouest, Centre-Sud, Est & Cacao, Hévéa, Palmier à huile, Banane/Plantain, Manioc, Pêche maritime/continentale \\ \hline
\textbf{Guinéenne} (Unimodale longue) & Adamaoua, Est, Centre, Sud (parties nord) & Café robusta, Cacao, Tabac, Ananas, Maïs, Ignames, Élevage porcin/avicole familial \\ \hline
\textbf{Soudanienne} (Unimodale moyenne) & Nord de l'Adamaoua, Nord (Bénoué), Nord-Ouest & Coton, Maïs, Sorgho, Mil, Arachide, Niébé, Élevage bovin (transhumance/ranching) \\ \hline
\textbf{Soudano-sahélienne} (Unimodale courte) & Extrême-Nord, Nord (Diamaré, Mayo-Danay) & Mil, Sorgho, Riz (irrigué SEMRY), Oignon, Arachide, Élevage bovin extensif, Pêche lacustre (Tchad, Lagdo, Maga) \\ \hline
\textbf{Hauts Plateaux} (Altitude) & Ouest, Nord-Ouest, parties Adamaoua & Café arabica, Pomme de terre, Haricot, Maïs, Maraîchage, Élevage laitier amélioré (ranchs) \\ \hline
\end{tabularx}
\end{center}
\end{exemplebox}

\courssub{L'agriculture au Cameroun}

\begin{definition}[Types d'agriculture au Cameroun]
\begin{itemize}
    \item \textbf{Agriculture paysanne / familiale (90\% des exploitants) :} Exploitations $<$ 2 ha, polyculture, itinéraires techniques extensifs (itinérante sur brûlis, jachère), faible mécanisation, autoconsommation majeure, main-d'œuvre familiale.
    \item \textbf{Agriculture de plantation / agro-industrielle :} Grandes surfaces (siècles d'hectares), monoculture, capital intensif, main-d'œuvre salariée, forte mécanisation/intrants, exportation. Acteurs : \textbf{CDC} (Banane, Palmier, Hévéa - Sud-Ouest), \textbf{SOSUCAM} (Canne à sucre - Moungo, Nkoteng), \textbf{HEVECAM} (Hévéa - Sud), \textbf{SOCAPALM} (Palmier - Sanaga Maritime, Kienké), \textbf{SPFS} (Banane - Njombé), \textbf{UNEX PALM}.
    \item \textbf{Agriculture moderne paysanne / émergente :} Exploitations 5--50 ha, semi-mécanisées, intrants, contractualisation (filières coton, maïs, riz, ananas), souvent issues de programmes (PNDA, PIDMA, PEAJ).
\end{itemize}
\end{definition}

\begin{aretenir}[Principales filières végétales et organisation]
\begin{itemize}
    \item \textbf{Coton (Nord, Extrême-Nord, Adamaoua) :} Filière intégrée \textbf{SODECOTON} (encadrement, intrants à crédit, collecte, égrenage, export). 1er produit d'exportation agricole hors bois. Problèmes : cours mondial, parasites (jassides), foncier, changement climatique.
    \item \textbf{Cacao (Sud, Centre, Est, Littoral, SO) :} 4ème/5ème producteur mondial. Filière libéralisée encadrée par \textbf{CICC} (Comité Interprofessionnel Cacao-Café) et \textbf{ONCC}. Productivité faible ($\approx$ 300--500 kg/ha). Enjeu : rénovation vergers, qualité (fermentation), transformation locale (Chococam, Sic Cacaos, Neo Industry, Atlantic Cocoa).
    \item \textbf{Café (Arabica : Ouest/NO ; Robusta : Littoral/Sud/Est) :} Effondrement production (années 90), relance timide (Plan de relance café). Qualité arabica reconnue.
    \item \textbf{Banane/Plantain :} 1er producteur africain plantain. Banane dessert export (CDC, PHP, Boh Plantations) vers UE. Plantain vivrier national (flux Ouest/Centre $\to$ Yaoundé/Douala).
    \item \textbf{Céréales (Maïs, Riz, Mil/Sorgho) :} Maïs : 1ère céréale (toutes zones). Riz : Bassin du Logone/Chari (SEMRY, UNVDA), Ndop, Haut-Nyong, Mbam. Déficit couvert par importations (Asie). Mil/Sorgho : Nord/EN (sécurité alimentaire zone sahélienne).
    \item \textbf{Racines/Tubercules (Manioc, Ignames, Macabo, Patate douce) :} Bassin de production : Sud, Centre, Est, Ouest, Littoral. Manioc : transformation (foufou, gari, amidon, farine) en essor.
\end{itemize}
\end{aretenir}

\begin{experience}[Analyse de l'itinéraire technique du coton (SODECOTON)]
\textbf{Objectif :} Comprendre l'intensification encadrée.
\begin{enumerate}
    \item \textbf{Préparation :} Labour/traction animale ou motorisée (nov--fév).
    \item \textbf{Semis :} Poquet (3--4 graines) $\times$ 0,80 m $\times$ 0,40 m (juin, 1ères pluies utiles). Variétés : IRMA, FK 37.
    \item \textbf{Entretien :} 2--3 sarclages (manuel/chimique), buttage. Traitements phytosanitaires (jassides, pucerons, chenilles) : 3--5 passages avion/porté.
    \item \textbf{Récolte :} Manuelle (égrenage à la main) oct--déc. Rendement paysan moyen 1\,200--1\,500 kg/ha graine coton.
    \item \textbf{Commercialisation :} Prix indicateur fixé par l'État/SODECOTON $\to$ paiement différé (acompte + solde).
\end{enumerate}
\legende{Itinéraire type campagne cotonnière Nord-Cameroun. Source : SODECOTON.}
\end{experience}

\begin{attention}[Pièges fréquents : Agriculture]
\begin{itemize}
    \item Ne pas confondre \textbf{production} (volume total, tonnes) et \textbf{rendement} (productivité surfacique, kg/ha). Une production peut augmenter par extension des surfaces sans gain de rendement.
    \item Situer \textbf{précisément} les bassins : « Nord » ne suffit pas ; distinguer Bénoué (coton/maïs), Mayo-Rey (riz/élevage), Diamaré (coton/mil), Mayo-Danay (riz/SEMRY).
    \item Ne pas oublier l'\textbf{agriculture péri-urbaine} (maraîchage : tomate, chou, piment, amarante autour de Yaoundé, Douala, Bafoussam, Ngaoundéré) qui approvisionne les villes.
\end{itemize}
\end{attention}

\courssub{L'élevage au Cameroun}

\begin{definition}[Systèmes d'élevage (Typologie programme)]
\begin{enumerate}
    \item \textbf{Élevage traditionnel / Pastoral (Transhumance) :} Système extensif, mobilité spatiale (Nord $\leftrightarrow$ Sud selon saisons), parcours libres, races locales rustiques (Goudali, Red Fulani, Kouri, Azawak), cheptel = capital social/épargne. Zones : Adamaoua, Nord, Extrême-Nord, Nord-Ouest (Monts Mandara).
    \item \textbf{Élevage sédentaire amélioré / Semi-intensif :} Exploitation familiale fixe, pâturages aménagés (PNIA : Péncélés, Ndokayo, Waza-Logone), cultures fourragères, alimentation complémentaire (son, tourteaux, résidus), races locales améliorées ou croisées. Zones : Adamaoua (fermes pilotes), Ouest (lait), Péri-urbain.
    \item \textbf{Élevage moderne / Intensif (Ranching, Hors-sol) :} Capital élevé, races exotiques/améliorées (Holstein, Montbéliarde, Charolaise, croisés F1), alimentation maîtrisée (ensilage, foin, concentrés, tourteaux de coton/soja), bâtiments, santé vétérinaire rigoureuse. Zones : \textbf{Ranchs} (Ngaoundéré, Maroua, Kounden, Bamyanga, Waza) ; \textbf{Hors-sol avicole/porcin/laitier} (Ceintures vertes Yaoundé, Douala, Bafoussam, Ebolowa).
\end{enumerate}
\end{definition}

\begin{aretenir}[Cheptel national (Estimations MINEPIA récentes)]
\begin{itemize}
    \item \textbf{Bovins :} $\approx$ \textbf{6--7 millions} de têtes (majorité Nord/Adamaoua/EN). Races : Goudali (lait/viande), Red Fulani (viande), Kouri (lac Tchad, viande).
    \item \textbf{Ovins/Caprins :} $\approx$ \textbf{4--5 millions} (partout, rustiques).
    \item \textbf{Porcins :} $\approx$ \textbf{1,5--2 millions} (Ouest, Centre, Sud, Péri-urbain). Race locale + Piétrain/Large White.
    \item \textbf{Volaille :} $\approx$ \textbf{30--40 millions} (poulets de chair, pondeuses, pintades). Production industrielle (SIPA, CAMEROUN FARMS, fermes moyennes) + villageoise.
    \item \textbf{Production laitière :} $\approx$ \textbf{200--250 millions L/an} (majorité pastoral, faible collecte formelle $\approx$ 10\%). Laiteries : LAITERIE DU BERGER (Ngaoundéré), SOTUC (Douala), fermes laitières Ouest.
\end{itemize}
\end{aretenir}

\begin{exemplebox}[Transhumance : Dynamique spatiale et conflits]
\begin{itemize}
    \item \textbf{Couloirs de transhumance :} Axes Nord-Sud balisés (ex: Mayo Rey $\to$ Mbéré $\to$ Adamaoua). Largeur légale 50--100 m, souvent empiétés par cultures.
    \item \textbf{Calendrier :} Montée (saison sèche nov--avr) vers zones humides (Bénoué, Logone, Mbéré, Sanaga) ; Descente (saison des pluies mai--oct) vers zones sèches (Adamaoua, plateaux).
    \item \textbf{Conflits agro-pastoraux :} Destruction cultures par bétail, blocage couloirs, point d'eau, vol. Mécanismes : Commissions de conciliation (Préfets, Lamibé), Loi 74/1 (foncier), Loi 98/005 (régime pastoral).
\end{itemize}
\end{exemplebox}

\begin{methode}[Comparer les systèmes d'élevage (Tableau de synthèse)]
\textbf{Objectif :} Différencier les 3 modèles selon critères géographiques.
\begin{enumerate}
    \item \textbf{Critères (colonnes) :} Mobilité / Occupation espace ; Alimentation ; Races / Productivité ; Capital / Main-d'œuvre ; Commercialisation ; Contraintes majeures.
    \item \textbf{Remplissage :} Mots-clés, données chiffrées (rendements lait/viande).
    \item \textbf{Synthèse :} Gradient extensif (pastoral) $\to$ intensif (hors-sol) ; Localisation : Nord/Adamaoua vs Péri-urbain/Ouest.
\end{enumerate}
\textbf{Attention :} Distinguer \emph{Ranching} (extensif sur grands espaces clôturés, races améliorées) et \emph{Hors-sol} (intensif, bâtiment, zéro pâturage).
\end{methode}

\courssub{La pêche au Cameroun}

\begin{definition}[Types de pêche]
\begin{enumerate}
    \item \textbf{Pêche maritime :}
    \begin{itemize}
        \item \textbf{Artisanale :} Pirogues (4\,500--5\,000), moteurs hors-bord (15--40 CV), filets maillants, senne, palangre. Équipages 3--8 pêcheurs. Zones : Côte (Limbé, Idenau, Kribi, Campo, Douala/Wouri). Cibles : Poissons pélagiques (sardinelles, chinchards), démersaux (bars, capitaines), crevettes (pêche côtière).
        \item \textbf{Industrielle :} Chalutiers (crevettiers, poissonniers), senneurs. $\approx$ 80--100 navires. Basés Douala/Kribi. Export crevettes (UE), poisson local/export. Puissance $\approx$ 600 kW.
    \end{itemize}
    \item \textbf{Pêche continentale :} Lac Tchad (partage), Lac Lagdo, Lac Mbakaou, Barrage Lom Pangar, Fleuves (Sanaga, Bénoué, Logone, Chari, Nyong, Wouri). Pirogues non motorisées/motorisées. Filets, nasses, épuisettes. Cibles : Tilapia, Clarias (poisson-chat), Heterotis, Carpes. Activité majeure pour protéines Nord/Est.
    \item \textbf{Aquaculture :} Écloseries (Melen, Foumban, Kotto, Maga), cages flottantes (Lagdo, Lom Pangar), étangs villageois. Espèces : Tilapia nilotica, Clarias gariepinus, Heterotis niloticus. Production encore faible ($<$ 5\,000 t) mais en croissance (PDA, PEAJ).
\end{enumerate}
\end{definition}

\begin{aretenir}[Indicateurs clés de la pêche (2022 env.)]
\begin{itemize}
    \item \textbf{Captures totales :} $\approx$ \textbf{250\,000--270\,000 tonnes} (Artisanale maritime $\approx$ 70--75\%, Continentale $\approx$ 20--25\%, Industrielle $\approx$ 5--10\%).
    \item \textbf{CPUE (Capture Par Unité d'Effort) :} Indicateur d'abondance/pression.
    \begin{itemize}
        \item Artisanale : kg / (pirogue $\times$ jour de mer). Tendance \textbf{baissière} (surpêche croissance, juvéniles, destruction mangroves).
        \item Industrielle : kg / (kW $\times$ jour) ou t/jour. Variable selon stocks crevettes/poissons.
    \end{itemize}
    \item \textbf{Emploi :} $\approx$ 200\,000--300\,000 pêcheurs directs (artisanale), 10\,000 industriels. Femmes : transformation (fumage, séchage, étuvage) et commercialisation.
\end{itemize}
\end{aretenir}

\begin{exoresolu}[Diagnostic de la pêche maritime au Cameroun (Données 2020-2022)]
\textbf{Énoncé :} Le tableau présente les captures et l'effort de pêche pour les deux segments de la flottille camerounaise.
\begin{center}
\begin{tabular}{|l|c|c|c|}\hline
\rowcolor{popblueL} \textbf{Indicateur} & \textbf{Artisanale (2022)} & \textbf{Industrielle (2022)} & \textbf{Total 2020} \\ \hline
Captures (tonnes) & 185\,000 & 65\,000 & 230\,000 \\ \hline
Nombre d'unités (pirogues / navires) & 4\,500 & 85 & - \\ \hline
Jours de mer moyens / an / unité & 180 & 220 & - \\ \hline
Puissance moyenne (kW) & 15 (moteur hors-bord) & 600 & - \\ \hline
\end{tabular}
\end{center}
\textbf{Travail :}
\begin{enumerate}
    \item Calculer le CPUE (en kg/jour de mer) pour chaque segment en 2022.
    \item Calculer la part de la pêche artisanale dans les captures totales 2022.
    \item En 2020, le CPUE artisanale était de 240 kg/jour. Commenter l'évolution de la productivité artisanale entre 2020 et 2022.
    \item Proposer deux mesures de gestion durable adaptées au contexte camerounais.
\end{enumerate}
\tcblower
\textbf{Solution détaillée :}
\begin{enumerate}
    \item \textbf{CPUE 2022 :}
    \begin{itemize}
        \item \textbf{Artisanale :} Effort total = $4\,500 \times 180 = 810\,000$ pirogue-jour.
        $CPUE_{art} = \frac{185\,000\,000}{810\,000} \approx \textbf{228 kg / pirogue-jour}$.
        \item \textbf{Industrielle :} Effort total = $85 \times 220 = 18\,700$ navire-jour.
        $CPUE_{ind} = \frac{65\,000\,000}{18\,700} \approx \textbf{3\,476 kg / navire-jour}$ ($\approx 5,8$ t/jour).
    \end{itemize}
    \item \textbf{Part artisanale 2022 :} Total 2022 = $250\,000$ t. $\%_{art} = \frac{185\,000}{250\,000} \times 100 = \textbf{74\%}$.
    \item \textbf{Évolution CPUE Artisanale :} $240 \to 228$ kg/j. Variation $= \frac{228-240}{240} \times 100 = \textbf{-5\%}$.
    \textbf{Interprétation :} Baisse de la productivité unitaire. Signe d'\textbf{augmentation de l'effort} (motorisation, nombre pirogues) non compensée par le stock. Risque de \cle{surpêche} (croissance/recrutement), aggravé par destruction mangroves (nurseries) et pollution estuaire Wouri.
    \item \textbf{Mesures de gestion durable :}
    \begin{enumerate}
        \item \textbf{Respect repos biologique \& mailles :} Contrôle strict mailles ($\ge$ 40 mm crevettes, $\ge$ 60 mm poissons), interdiction pêche zones frayères (embouchures, mangroves) pics reproduction (juil--sept).
        \item \textbf{Gestion concertée capacités :} Plafonnement pirogues motorisées / licences industrielles ; programme \textbf{reconversion/réduction flotte} artisanale excédentaire (rachat, formation) + \textbf{valorisation captures} (chaîne de froid, transformation) pour revenu/kg $\uparrow$ au lieu volume $\uparrow$.
    \end{enumerate}
\end{enumerate}
\textbf{Conclusion :} Pêche artisanale domine volume (74\%) mais productivité décline (-5\%/2 ans), révélant pression excessive stock démersal côtier. Durabilité impose gestion effort (licences) et protection habitats (mangroves).
\end{exoresolu}

\courssub{TD : Problèmes et perspectives du secteur agropastoral et piscicole}

\begin{methode}[Analyser un tableau statistique de production agricole (rendements, taux de variation, parts relatives)]
\textbf{Objectif :} Exploiter des données chiffrées (superficies, productions, rendements) pour caractériser une spéculation et comparer son évolution.
\begin{enumerate}
    \item \textbf{Identifier les variables :} Années (lignes), Spéculations (colonnes), Unités (tonnes, hectares, kg/ha, \%).
    \item \textbf{Calculer le rendement :} $R = \frac{\text{Production (t)}}{\text{Superficie (ha)}} \times 1000$ (en kg/ha).
    \item \textbf{Calculer le taux de variation (TV) :} $TV = \frac{V_{\text{final}} - V_{\text{initial}}}{V_{\text{initial}}} \times 100$ (en \%).
    \item \textbf{Calculer la part relative :} $Part = \frac{\text{Valeur spéculation}}{\text{Total}} \times 100$.
    \item \textbf{Interpréter :} Comparer rendements aux moyennes nationales/mondiales ; relier évolution aux facteurs (climat, intrants, prix, politique agricole).
    \item \textbf{Rédiger la conclusion :} Phrase synthétique chiffrée (ex : « La production de maïs a augmenté de 15\% entre 2018 et 2022, portée par l'extension des superficies (+10\%) et l'amélioration du rendement (+5\%). »).
\end{enumerate}
\textbf{Réflexe Bac :} Vérifier l'homogénéité des unités avant tout calcul. Arrondir au dixième pour les taux, à l'entier pour les rendements.
\end{methode}

\begin{exoresolu}[Évolution de la filière coton au Cameroun (Données inspirées SODECOTON)]
\textbf{Énoncé :} Le tableau ci-dessous présente l'évolution de la campagne cotonnière dans la région du Nord sur 3 ans.
\begin{center}
\begin{tabularx}{\linewidth}{|c|X|X|X|}\hline
\rowcolor{popblueL} \textbf{Campagne} & \textbf{Superficie emblavée (ha)} & \textbf{Production graine (tonnes)} & \textbf{Rendement moyen (kg/ha)} \\ \hline
2020-2021 & 220\,000 & 286\,000 & ? \\ \hline
2021-2022 & 245\,000 & 318\,500 & ? \\ \hline
2022-2023 & 230\,000 & 276\,000 & ? \\ \hline
\end{tabularx}
\end{center}
\textbf{Travail demandé :}
\begin{enumerate}
    \item Compléter le tableau en calculant les rendements moyens.
    \item Calculer le taux de variation de la production entre 2020-2021 et 2022-2023.
    \item Analyser l'évolution de la productivité (rendement) et proposer deux explications géographiques à la baisse observée en 2022-2023.
\end{enumerate}
\tcblower
\textbf{Solution détaillée :}
\begin{enumerate}
    \item \textbf{Calcul des rendements (kg/ha) :}
    \begin{itemize}
        \item 2020-2021 : $R_1 = \frac{286\,000}{220\,000} \times 1000 = 1\,300$ kg/ha.
        \item 2021-2022 : $R_2 = \frac{318\,500}{245\,000} \times 1000 = 1\,300$ kg/ha.
        \item 2022-2023 : $R_3 = \frac{276\,000}{230\,000} \times 1000 = 1\,200$ kg/ha.
    \end{itemize}
    Tableau complété :
    \begin{center}
    \begin{tabularx}{\linewidth}{|c|X|X|X|}\hline
    \rowcolor{popblueL} \textbf{Campagne} & \textbf{Superficie (ha)} & \textbf{Production (t)} & \textbf{Rendement (kg/ha)} \\ \hline
    2020-2021 & 220\,000 & 286\,000 & 1\,300 \\ \hline
    2021-2022 & 245\,000 & 318\,500 & 1\,300 \\ \hline
    2022-2023 & 230\,000 & 276\,000 & 1\,200 \\ \hline
    \end{tabularx}
    \end{center}
    \item \textbf{Taux de variation de la production (2020-2021 $\to$ 2022-2023) :}
    $$TV = \frac{276\,000 - 286\,000}{286\,000} \times 100 = \frac{-10\,000}{286\,000} \times 100 \approx -3,5\%$$
    La production a baissé de \textbf{3,5\%} sur la période.
    \item \textbf{Analyse de la productivité :}
    Le rendement est stable à 1\,300 kg/ha les deux premières années, puis chute à 1\,200 kg/ha en 2022-2023 (-7,7\%).
    \textbf{Explications géographiques :}
    \begin{itemize}
        \item \textbf{Facteur climatique :} La campagne 2022-2023 a connu un début de saison pluvieuse tardif et des poches de sécheresse en juillet-août (phénomène fréquent en zone soudano-sahélienne), perturbant la floraison et la formation des pommes.
        \item \textbf{Facteur technique / économique :} Difficultés d'accès aux intrants (engrais, herbicides) à temps suite à la hausse des prix internationaux (contexte post-Covid / guerre en Ukraine) et retards de paiement des cotonculteurs par la SODECOTON, réduisant l'intensification culturale.
    \end{itemize}
\end{enumerate}
\textbf{Conclusion :} Malgré une relative stabilité des superficies, la filière cotonnière subit une érosion de sa productivité en 2022-2023, liée à la conjonction d'aléas climatiques et de contraintes socio-économiques limitant l'itinéraire technique.
\end{exoresolu}

\begin{methode}[Construire et interpréter un calendrier cultural (agro-climatique)]
\textbf{Objectif :} Mettre en relation le régime de pluviométrie (hauteurs, nombre de mois humides) et le cycle cultural d'une spéculation pour justifier sa localisation.
\begin{enumerate}
    \item \textbf{Identifier la zone agro-écologique :} Équatoriale (bimodale), Guinéenne (unimodale longue), Soudanienne (unimodale moyenne), Soudano-sahélienne (unimodale courte).
    \item \textbf{Repérer la saison des pluies :} Mois où $P > \text{ETP}$ (ou $P > 100$ mm). Noter début, fin, durée.
    \item \textbf{Connaître le cycle de la plante :} Durée végétative (semis $\to$ récolte), besoins en eau par stade (levée, floraison, remplissage).
    \item \textbf{Positionner les opérations :} Préparation sol (avant pluies), Semis (début pluies), Sarclage/Buttage (1-2 mois après), Récolte (fin pluies / début saison sèche).
    \item \textbf{Représenter graphiquement :} Diagramme ombrothermique simplifié (barres pluie, courbe T°) + bande colorée sous l'axe des mois pour le cycle cultural.
    \item \textbf{Justifier l'adaptation :} « Le cycle de 90-120 jours du mil correspond à la durée de la saison des pluies (juin-septembre) dans le Nord. »
\end{enumerate}
\textbf{Formule clé :} Durée saison des pluies $\approx$ Nombre de mois $P > 100$ mm. Cycle cultural $\le$ Durée saison des pluies (sauf irrigation/résiduelle).
\end{methode}

\begin{exoresolu}[Calendrier cultural du maïs en zone bimodale (Région du Centre)]
\textbf{Énoncé :} À partir des données climatiques de Yaoundé (région du Centre) ci-dessous, construire le calendrier cultural du maïs (cycle court : 90-100 jours) pour les deux saisons.
\begin{center}
\begin{tabular}{|c|cccccccccccc|}\hline
\rowcolor{popblueL} \textbf{Mois} & J & F & M & A & M & J & J & A & S & O & N & D \\ \hline
\textbf{Pluie (mm)} & 20 & 40 & 130 & 180 & 200 & 130 & 60 & 100 & 220 & 280 & 150 & 30 \\ \hline
\textbf{Temp (°C)} & 25 & 26 & 26 & 25 & 25 & 24 & 23 & 23 & 24 & 24 & 24 & 25 \\ \hline
\end{tabular}
\end{center}
\textbf{Travail :}
\begin{enumerate}
    \item Identifier les deux saisons des pluies (seuil $P \ge 100$ mm).
    \item Déterminer les fenêtres de semis optimales pour le maïs.
    \item Représenter le calendrier sous forme de frise chronologique (texte structuré).
    \item Expliquer pourquoi le maïs de 2ème saison est souvent plus productif.
\end{enumerate}
\tcblower
\textbf{Solution détaillée :}
\begin{enumerate}
    \item \textbf{Saisons des pluies (P $\ge$ 100 mm) :}
    \begin{itemize}
        \item \textbf{1ère saison (Grande saison) :} Mars (130) $\to$ Juin (130). \emph{Durée : 4 mois (Mars, Avril, Mai, Juin).} Note : Juillet (60) et Août (100) marquent la « petite saison sèche » (aoûtien).
        \item \textbf{2ème saison (Petite saison) :} Septembre (220) $\to$ Novembre (150). \emph{Durée : 3 mois (Sept, Oct, Nov).}
    \end{itemize}
    \item \textbf{Fenêtres de semis (Maïs cycle 90-100 j $\approx$ 3-3,5 mois) :}
    \begin{itemize}
        \item \textbf{Saison 1 :} Semis \textbf{mi-mars} $\to$ Récolte \textbf{mi-juin/début juillet}. (Risque : stress hydrique fin mai/juin si pause pluviométrique).
        \item \textbf{Saison 2 :} Semis \textbf{mi-août / début septembre} $\to$ Récolte \textbf{mi-novembre / début décembre}. (Avant l'arrêt des pluies en décembre).
    \end{itemize}
    \item \textbf{Frise chronologique (Calendrier cultural) :}
    \begin{center}
    \begin{tabular}{|l|l|l|}\hline
    \rowcolor{popblueL} \textbf{Période} & \textbf{Opérations culturales} & \textbf{Stade plante / Climat} \\ \hline
    \textbf{Fév - Mi-Mars} & Préparation terrain, apport fumure & Sol nu, T° élevées \\ \hline
    \textbf{Mi-Mars - Mi-Juin} & \textbf{SAISON 1} : Semis, 2 sarclages, récolte & Levée $\to$ Floraison (Mai) $\to$ Maturation \\ \hline
    \textbf{Juillet - Août} & Entretien jachère / Culture légumineuse & Petite saison sèche (Aoûtien) \\ \hline
    \textbf{Mi-Août - Mi-Nov} & \textbf{SAISON 2} : Semis, 2 sarclages, récolte & Levée $\to$ Floraison (Oct) $\to$ Maturation \\ \hline
    \textbf{Déc - Janv} & Récolte, stockage, commercialisation & Saison sèche \\ \hline
    \end{tabular}
    \end{center}
    \item \textbf{Productivité 2ème saison > 1ère saison :}
    \begin{itemize}
        \item \textbf{Régularité pluviométrique :} La 2ème saison (sept-nov) est continue et intense (220-280 mm/mois), sans pause sèche critique pendant la floraison (octobre).
        \item \textbf{Ensoleillement :} Meilleure radiation solaire en octobre-novembre (moins de nébulosité qu'en avril-mai) favorisant la photosynthèse et le remplissage des grains.
        \item \textbf{Pression parasitaire :} Moins forte en début de 2ème saison qu'en fin de 1ère saison (accumulation des populations de foreurs/striga).
    \end{itemize}
\end{enumerate}
\textbf{Conclusion :} Le climat bimodal du Centre permet deux cycles de maïs/an. La 2ème saison (semis août-sept) est plus sûre et plus productive grâce à la fiabilité des pluies d'automne.
\end{exoresolu}

\begin{methode}[Comparer les systèmes d'élevage (Tableau de synthèse / SWOT simplifié)]
\textbf{Objectif :} Différencier l'élevage traditionnel (transhumance), l'élevage amélioré (parcs de vaccination, pâturages aménagés) et l'élevage moderne (ranching, hors-sol) selon des critères géographiques.
\begin{enumerate}
    \item \textbf{Définir les 3 modèles types du programme :}
    \begin{itemize}
        \item \cle{Transhumance} (nomadisme/pastoralisme) : Mobilité spatiale, parcours libres, cheptel important, faible productivité/tête.
        \item \cle{Élevage sédentaire amélioré} : Exploitation familiale, pâturages aménagés (PNIA), alimentation complémentaire, races locales améliorées.
        \item \cle{Élevage moderne / Ranch / Hors-sol} : Capital intensif, races améliorées/exotiques, alimentation maîtrisée (ensilage, tourteaux), forte productivité, proximité marchés urbains.
    \end{itemize}
    \item \textbf{Choisir les critères de comparaison (colonnes) :} Mobilité / Sédentarisation ; Mode d'alimentation ; Race / Productivité (lait/viande/tête/an) ; Capital / Main d'œuvre ; Impact environnemental ; Commercialisation.
    \item \textbf{Remplir le tableau par cases courtes (mots-clés).}
    \item \textbf{Synthétiser :} Identifier le gradient d'intensification (extensif $\to$ intensif) et la localisation spatiale (Nord vs Péri-urbain).
\end{enumerate}
\textbf{Attention :} Ne pas confondre « élevage moderne » et « ranching » (le ranching est un type d'élevage moderne extensif sur grands espaces, le hors-sol est intensif).
\end{methode}

\begin{exoresolu}[Analyse comparative : Transhumance vs Ranch de Ngaoundéré]
\textbf{Énoncé :} Complétez le tableau comparatif ci-dessous pour opposer l'élevage transhumant peul du Nord-Cameroun (Adamaoua/Nord/Extrême-Nord) et l'élevage de type ranch (ex: Ranch de Ngaoundéré ou fermes péri-urbaines de Douala/Yaoundé).
\begin{center}
\begin{tabularx}{\linewidth}{|l|X|X|}\hline
\rowcolor{popblueL} \textbf{Critères} & \textbf{Élevage Transhumant (Pastoral)} & \textbf{Élevage Moderne (Ranch / Hors-sol)} \\ \hline
Mobilité / Occupation espace & \textit{À compléter} & \textit{À compléter} \\ \hline
Alimentation du cheptel & \textit{À compléter} & \textit{À compléter} \\ \hline
Races / Productivité (Lait/Viande) & \textit{À compléter} & \textit{À compléter} \\ \hline
Capital / Main d'œuvre & \textit{À compléter} & \textit{À compléter} \\ \hline
Commercialisation / Circuits & \textit{À compléter} & \textit{À compléter} \\ \hline
Contraintes majeures & \textit{À compléter} & \textit{À compléter} \\ \hline
\end{tabularx}
\end{center}
\tcblower
\textbf{Solution détaillée (Tableau complété) :}
\begin{center}
\begin{tabularx}{\linewidth}{|l|X|X|}\hline
\rowcolor{popblueL} \textbf{Critères} & \textbf{Élevage Transhumant (Pastoral)} & \textbf{Élevage Moderne (Ranch / Hors-sol)} \\ \hline
Mobilité / Occupation espace & \textbf{Forte mobilité} (transhumance Nord-Sud saisonnière). Parcours libres, conflits usage sol (agriculteurs). & \textbf{Sédentaire}. Emprise foncière fixe (titres fonciers), clôtures, paddocks. Proximité axes/ville. \\ \hline
Alimentation du cheptel & \textbf{Parcours naturel} (herbacées sahéliennes, résidus de récolte). Complément rare (sel, son). Déficit fourrager saison sèche. & \textbf{Maîtrisée} : Pâturages améliorés (Brachiaria, Panicum), ensilage, foin, tourteaux (coton, soja), concentrés. Distribution d'eau permanente. \\ \hline
Races / Productivité & \textbf{Races locales} (Goudali, Red Fulani, Kouri). Rustiques, faible rendement (Lait : 1-2 L/j/vache ; Carcasse : 150-200 kg). & \textbf{Races améliorées/exotiques} (Holstein, Montbéliarde, Charolaise, croisés F1). Haute productivité (Lait : 15-25 L/j ; Carcasse : 350-450 kg). \\ \hline
Capital / Main d'œuvre & \textbf{Faible capital} (cheptel = capital), \textbf{Main d'œuvre familiale} (bergers), savoir traditionnel. & \textbf{Capital élevé} (bâtiments, matériel, génétique, intrants), \textbf{Main d'œuvre salariée} qualifiée (techniciens, vétérinaires). \\ \hline
Commercialisation / Circuits & \textbf{Filère longue/informelle} : Marchés à bétail (Garoua, Maroua, Ngaoundéré) $\to$ Grossistes $\to$ Abattoirs urbains (Douala/Yaoundé). Perte de poids au trek. & \textbf{Filière courte/organisée} : Vente directe aux abattoirs modernes, supermarchés, industries laitières (SOTUC, LAITERIE DU BERGER). Traçabilité. \\ \hline
Contraintes majeures & \textbf{1. Foncier} (empiétement cultures, blocage couloirs). \textbf{2. Sanitaire} (péripneumonie, fièvre aphteuse, absence vétérinaire). \textbf{3. Climatique} (sécheresse, incendies). & \textbf{1. Coût intrants} (aliments importés, énergie). \textbf{2. Foncier urbain} (coût, pression). \textbf{3. Marché} (concurrence importations, pouvoir d'achat). \\ \hline
\end{tabularx}
\end{center}
\textbf{Analyse de synthèse :} L'opposition structurelle repose sur le \textbf{gradient d'intensification}. Le système pastoral valorise des espaces marginaux (basse pression foncière) par la mobilité (extensif), mais bute sur la démographie et le changement climatique. Le système moderne substitue le capital à l'espace (intensif), assure l'autosécurité alimentaire urbaine, mais reste dépendant des importations d'intrants (tourteaux, races).
\end{exoresolu}

\begin{methode}[Calculer et interpréter des indicateurs de la pêche (CPUE, Part des captures, Taux d'exploitation)]
\textbf{Objectif :} Quantifier la pression de pêche et la performance des flottilles (artisanale vs industrielle) pour diagnostiquer l'état du stock.
\begin{enumerate}
    \item \textbf{Identifier les données :} Captures totales (tonnes), Effectif navires/pirogues, Puissance moteurs (kW/J), Jours de mer / Marées, Effectif pêcheurs.
    \item \textbf{Calculer le CPUE (Capture Par Unité d'Effort) :} $CPUE = \frac{\text{Captures (kg ou t)}}{\text{Effort de pêche}}$.
    \begin{itemize}
        \item Effort artisanale : Nb pirogues $\times$ Jours de mer (ou Nb pêcheurs $\times$ Jours).
        \item Effort industrielle : Puissance (kW) $\times$ Jours de mer (kW.jour).
    \end{itemize}
    \item \textbf{Calculer la part relative :} $\% \text{Artisanale} = \frac{\text{Captures Art.}}{\text{Captures Totales}} \times 100$.
    \item \textbf{Interpréter l'évolution du CPUE :}
    \begin{itemize}
        \item CPUE \textbf{stable/hausse} : Stock sain, effort maîtrisé.
        \item CPUE \textbf{baisse continue} : \cle{Surpêche} (effort excessif), dégradation habitats, changement climatique.
    \end{itemize}
    \item \textbf{Conclure sur la durabilité :} Croiser CPUE, tailles moyennes capturées (juvéniles ?), et réglementation (repos biologique, mailles).
\end{enumerate}
\textbf{Unités :} CPUE artisanale en kg/pirogue/jour ou kg/pêcheur/jour. CPUE industrielle en kg/kW/jour (ou t/jour de mer).
\end{methode}

\begin{methode}[Rédiger une réponse argumentée « Problèmes et Perspectives » (Structure Problème $\to$ Causes $\to$ Conséquences $\to$ Solutions)]
\textbf{Objectif :} Structurer une dissertation ou une réponse longue (type Bac) sur les contraintes du secteur agropastoral et les politiques de relance.
\begin{enumerate}
    \item \textbf{Analyser le sujet :} Identifier le secteur visé (Agriculture ? Élevage ? Pêche ? Global ?), l'échelle (Cameroun, Région Nord, Exploitation familiale) et l'injonction (Expliquer, Proposer, Évaluer).
    \item \textbf{Structurer en 3 parties logiques (ou 2 parties + synthèse) :}
    \begin{itemize}
        \item \textbf{I. Diagnostic des contraintes (Problèmes) :} Grouper en facteurs \textbf{Naturels} (Climat, Sol, Maladies), \textbf{Humains/Structurels} (Foncier, Formation, Financement, Organisation), \textbf{Économiques/Conjoncturels} (Prix, Intrants, Infrastructures, Concurrence importations).
        \item \textbf{II. Conséquences socio-économiques :} Insécurité alimentaire, déficit balance commerciale, pauvreté rurale, exode, conflits agro-pastoraux.
        \item \textbf{III. Perspectives / Politiques de relance (Solutions) :} Relier \emph{chaque solution à un problème cité en I}. Ex: Problème intrants chers $\to$ Solution : Subvention engrais / Usine engrais locale (Kribi). Problème foncier $\to$ Solution : Loi foncière / Délimitation couloirs transhumance.
    \end{itemize}
    \item \textbf{Illustrer par des exemples précis (Cameroun) :} Citer les programmes (PNDA, PIDMA, PARIC, PEAJ, SODECOTON, CDC, SEMRY, UNVDA, ONUDI), les infrastructures (Barrage de Lom Pangar pour irrigation, Port de Kribi), les zones (Bassin de la Mbéré, Logone, Haut-Nyong).
    \item \textbf{Conclusion ouverte :} Condition de réussite (Gouvernance, Paix sociale, Adaptation climatique).
\end{enumerate}
\textbf{Mots-clés de liaison :} \emph{En effet, Par ailleurs, Toutefois, Dès lors, C'est pourquoi, Corollaire.}
\end{methode}

\begin{exoresolu}[Sujet type Bac : « Les contraintes du développement agricole au Cameroun et les pistes de relance »]
\textbf{Énoncé :} En vous appuyant sur des exemples précis, présentez les principaux obstacles au développement de l'agriculture camerounaise et les stratégies mises en œuvre pour y remédier.
\tcblower
\textbf{Solution détaillée (Plan rédigé type) :}
\textbf{Introduction :}
L'agriculture emploie \textbf{60\% de la population active} et contribue à \textbf{$\approx$ 15-17\% du PIB} (chiffres récents). Malgré des potentialités immenses (diversité climatique, sols, eau), le secteur reste marqué par une \textbf{faible productivité} et une \textbf{insécurité alimentaire structurelle} (importations massives : riz, blé, poisson, huile). Quels en sont les freins et les réponses apportées ?

\textbf{I. Des contraintes multiformes freinent la performance agricole}
\textbf{A. Contraintes naturelles et sanitaires}
L'irrégularité pluviométrique (sécheresses Nord, excès Sud), l'appauvrissement des sols (acidité, érosion, minage minier) et la pression parasitaire (chute armée d'automne sur maïs, striga, cochenille du manioc) limitent les rendements. L'absence de maîtrise de l'eau (seulement \textbf{$\approx$ 1\% des terres irrigables} équipées) rend l'agriculture \textbf{pluviale et aléatoire}.
\textbf{B. Contraintes structurelles et économiques}
Le \textbf{foncier coutumier} freine l'investissement long terme (absence titre foncier). L'\textbf{accès au financement} est quasi inexistant pour les petits exploitants (taux usuraires, absence garanties). Le \textbf{coût des intrants} (engrais, semences améliorées) explose (dépendance importation). Les \textbf{infrastructures rurales} (pistes, électricité, stockage, transformation) sont défaillantes, causant \textbf{30-40\% de pertes post-récolte}.
\textbf{C. Contraintes organisationnelles}
Atomisation des exploitations (moyenne $<$ 1,5 ha), vieillissement paysan, faible encadrement (ratio vulgarisateur/producteur $\approx$ 1/3000), désorganisation des filières (multiplicité intermédiaires, prix producteur bas).

\textbf{II. Des politiques de relance ciblées mais inégalement efficaces}
\textbf{A. Modernisation de l'outil de production (Offre)}
\textit{Programmes phares :} \textbf{PNDA} (Programme National de Développement Agricole) : intrants subventionnés, mécanisation (CUMA), semences. \textbf{PIDMA} (Projet Intégré Développement Maïs) : semences hybrides, irrigation. \textbf{SEMRY/UNVDA} : Aménagement hydro-agricole (riz irrigué Logone/Chari, Ndop). \textbf{Usine d'engrais de Kribi} (projet) pour réduire coût intrants.
\textbf{B. Organisation des acteurs et marchés (Demande/Institutions)}
Création des \textbf{Coopératives / OP} (Organisations de Producteurs) pour mutualiser achats/ventes. \textbf{Contrats de filière} (Coton, Cacao, Banane, Ananas) sécurisant débouchés et prix plancher. \textbf{PEAJ} (Programme Entreprenariat Agricole Jeunes) pour rajeunir et professionnaliser.
\textbf{C. Infrastructure et Gouvernance}
Programme \textbf{« Routes agricoles »} (désenclavement bassins). \textbf{Transformation locale} (Usines cacao, coton, tomate, manioc) pour capter valeur ajoutée. \textbf{Politique de substitution aux importations} (Stratégie Nationale de Développement - SND30) : Objectif \textbf{autosuffisance riz/maïs/poisson}.

\textbf{Conclusion :}
Le passage d'une agriculture de subsistance à une agriculture \textbf{compétitive et durable} exige de lever le \textbf{verrou foncier et financier}, de généraliser la \textbf{maîtrise de l'eau} et de structurer les \textbf{filières amont-aval}. La réussite conditionne la \textbf{souveraineté alimentaire} et la \textbf{création d'emplois décents} pour la jeunesse, dans un contexte d'adaptation au changement climatique (Agriculture Climatiquement Intelligente - ACI).
\end{exoresolu}

\begin{methode}[Réaliser un croquis schématique de localisation des activités agropastorales (Carte mentale / Croquis de synthèse)]
\textbf{Objectif :} Produire un croquis lisible, légendé, localisant les grands domaines (cultures, élevage, pêche) et les flux, à l'échelle du Cameroun.
\begin{enumerate}
    \item \textbf{Contour de base :} Tracer le triangle Cameroun (Nord large, Sud pointe). Repères : Lac Tchad (N), Golfe Guinée (SW), Fleuves (Bénoué, Sanaga, Nyong, Logone/Chari).
    \item \textbf{Zonage agro-écologique (Fond de carte - crayon de couleur clair) :}
    \begin{itemize}
        \item \textbf{Zone Soudano-sahélienne (Nord/Extrême-Nord) :} Coton, Mil/Sorgho, Arachide, Oignon, \textbf{Élevage bovin extensif} (Transhumance), Pêche lacustre (Tchad, Lagdo, Maga).
        \item \textbf{Zone des Hauts Plateaux (Ouest/Nord-Ouest/Adamaoua) :} Café arabica, Pomme de terre, Haricot, Maïs, \textbf{Élevage laitier/bovin amélioré} (Ranchs), Maraîchage.
        \item \textbf{Zone Guinéenne (Adamaoua/Est/Centre-Sud) :} Café robusta, Cacao, Tabac, Ananas, Maïs, Ignames, \textbf{Élevage porcin/avicole} familial.
        \item \textbf{Zone Équatoriale (Sud/Littoral/Sud-Ouest) :} \textbf{Cacao, Hévéa, Palmier à huile, Banane/Plantain}, Forêt/Chasse, \textbf{Pêche maritime} (Kribi, Douala, Limbé, Idenau).
    \end{itemize}
    \item \textbf{Symboles (Figurés) :}
    \begin{itemize}
        \item \textbf{Points proportionnels} : Villes marchés / Ports / Usines (Douala, Yaoundé, Garoua, Ngaoundéré, Kribi, Bafoussam).
        \item \textbf{Flèches d'épaisseur variable} : Flux commerciaux (Coton Nord $\to$ Douala ; Cacao Sud $\to$ Douala/Kribi ; Bétail Nord $\to$ Yaoundé/Douala).
        \item \textbf{Hachures / Points /Traits} : Cultures industrielles (hachures régulières), Vivrières (points), Pâturages (traits horizontaux), Forêt (petits triangles).
    \end{itemize}
    \item \textbf{Légende organisée (Impératif) :} Titre, Auteur, Date, Échelle approximative, Flèche Nord. Ordre : 1. Milieu (Fleuves, Relief), 2. Productions (Cultures, Élevage, Pêche), 3. Infrastructures/Flux, 4. Villes.
    \item \textbf{Soin :} Propreté, écriture fine, coloriage dans le sens du trait, pas de ratures.
\end{enumerate}
\textbf{Barème Bac (indicatif) :} Contour/Repères (2 pts), Zonage/Contenu (6 pts), Figurés/Légende (6 pts), Soin/Esthétique (2 pts) = 16 pts + 4 pts analyse commentée.
\end{methode}

\begin{exoresolu}[Croquis de synthèse : « Le Cameroun : contrastes et complémentarités agropastorales »]
\textbf{Énoncé :} Réalisez un croquis schématique du Cameroun mettant en évidence la répartition spatiale des principales activités agropastorales et piscicoles ainsi que les principaux flux de commercialisation. La légende doit être organisée et hiérarchisée.
\tcblower
\textbf{Solution détaillée (Description du croquis attendu + Code TikZ indicatif) :}
\textbf{1. Analyse du sujet :} Montrer la \textbf{spécialisation régionale} (Nord = Coton/Elevage ; Centre/Sud/Est = Perennes/Cacao ; Ouest = Vivrières/Lait ; Côte = Pêche/Plantations) et l'\textbf{orientation des flux vers les ports/métropoles}.
\textbf{2. Éléments à faire figurer (Check-list pour l'élève) :}
\begin{itemize}
    \item \textbf{Hydrographie :} Lac Tchad, Bénoué, Sanaga, Nyong, Logone, Chari, Mbam.
    \item \textbf{Villes/Ports :} Douala (Port principal), Kribi (Port en eau profonde), Yaoundé, Garoua, Maroua, Ngaoundéré, Bafoussam, Bertoua, Ebolowa, Limbé, Idenau.
    \item \textbf{Productions (Zones colorées/hachurées) :}
    \begin{itemize}
        \item \textcolor{popgold}{\textbf{Zone Cotonnière}} : Nord, Extrême-Nord (hachures jaune/or).
        \item \textcolor{popgreen}{\textbf{Zone Cacaoyère}} : Sud, Centre, Est, Littoral, Sud-Ouest (hachures vert foncé / points marron).
        \item \textcolor{poporange}{\textbf{Zone Palmier/Hévéa/Banane}} : Littoral, Sud-Ouest, Sud (hachures orange).
        \item \textcolor{popblue}{\textbf{Zone Élevage Bovin}} : Adamaoua, Nord, Extrême-Nord (trait horizontal bleu / points bleus). \textit{Symboles Ranch : Carrés rouges (Ngaoundéré, Maroua, Kounden)}.
        \item \textcolor{poppurple}{\textbf{Zone Pêche Maritime}} : Côte (Limbé $\to$ Campo). \textit{Pêche Continentale : Lac Tchad, Lagdo, Mbakaou, Sanaga, Logone}.
    \end{itemize}
    \item \textbf{Flux (Flèches épaisses) :}
    \begin{itemize}
        \item Coton : Nord $\to$ Garoua $\to$ Douala (Export) / Usines locales (Garoua, Kaele).
        \item Cacao/Bois/Hévéa : Sud/Est $\to$ Yaoundé/Douala/Kribi (Export).
        \item Bétail sur pied : Adamaoua/Nord $\to$ Yaoundé/Douala (Abattoirs).
        \item Produits vivriers (Ignames, Plantains, Manioc) : Ouest/Centre/Sud $\to$ Yaoundé/Douala.
    \end{itemize}
\end{itemize}
\textbf{3. Légende type (Hiérarchisée) :}
\begin{center}
\begin{tabular}{|l|l|}\hline
\rowcolor{popblueL} \textbf{TITRE : LE CAMEROUN : CONTRASTES ET COMPLÉMENTARITÉS AGROPASTORALES} & \\ \hline
\textbf{REPÈRES} & \\ \hline
\tikz\draw[popblue, thick] (0,0) -- (1,0); & Limites internationales \\ \hline
\tikz\draw[popblue, thick] (0,0) -- (1,0); & Fleuves principaux (Bénoué, Sanaga, Logone, Nyong) \\ \hline
\tikz\fill[popblue!30] (0.2,0.1) rectangle (0.8,0.3); & Lac Tchad / Barrages (Lagdo, Mbakaou, Lom Pangar) \\ \hline
\tikz\draw[popdark, fill=popdark] (0.5,0.2) circle (0.1); & Villes / Ports majeurs (Douala, Kribi, Yaoundé, Garoua) \\ \hline
\textbf{PRODUCTIONS VÉGÉTALES} & \\ \hline
\tikz\fill[popgold!50] (0,0) rectangle (1,0.4); & Zone cotonnière (Nord, Extrême-Nord) \\ \hline
\tikz\fill[popgreen!50] (0,0) rectangle (1,0.4); & Zone cacaoyère (Sud, Centre, Est, Littoral, SO) \\ \hline
\tikz\fill[poporange!50] (0,0) rectangle (1,0.4); & Plantations industrielles (Palmier, Hévéa, Banane - Côte) \\ \hline
\tikz\fill[popgreen!20!popyellow] (0,0) rectangle (1,0.4); & Vivrières intenses (Ouest, Hauts-Plateaux, Bamiléké) \\ \hline
\textbf{ÉLEVAGE ET PÊCHE} & \\ \hline
\tikz\draw[popblue, thick] (0,0.2) -- (1,0.2); \tikz\draw[popblue, thick] (0,0.3) -- (1,0.3); & Zone d'élevage bovin extensif / Transhumance (Adamaoua, Nord, EN) \\ \hline
\tikz\draw[popred, fill=popred] (0.5,0.2) rectangle (0.6,0.3); & Ranchs / Élevage moderne (Ngaoundéré, Maroua, Kounden) \\ \hline
\tikz\draw[poppurple, ->, thick] (0,0.2) -- (1,0.2); & Pêche maritime (Zone côtière) \\ \hline
\tikz\fill[poppurple!30] (0.2,0.1) rectangle (0.8,0.3); & Pêche continentale (Lacs, Barrages, Fleuves) \\ \hline
\textbf{FLUX ET INFRASTRUCTURES} & \\ \hline
\tikz\draw[popdark, ->, thick, line width=1.5pt] (0,0.2) -- (1,0.2); & Flux d'exportation (Coton, Cacao, Bois, Café, Pétrole) \\ \hline
\tikz\draw[popdark, ->, thick, line width=1pt] (0,0.2) -- (1,0.2); & Flux intérieur (Bétail, Vivrières, Poisson) vers métropoles \\ \hline
\tikz\draw[popgray, dashed] (0,0.2) -- (1,0.2); & Axes routiers/ferroviaires structurants (Transcam, Nord-Sud) \\ \hline
\end{tabular}
\end{center}
\textbf{4. Commentaire guidé (Oral/Écrit) :} « Ce croquis met en évidence la \textbf{division naturelle du travail} : le Nord « grenier » (céréales, coton, viande) et le Sud « plantation » (cacao, bois, huile, pêche). La complémentarité s'opère via les flux Nord-Sud (bétail, coton, oignons) et Sud-Nord (produits manufacturés, poisson séché, huile). Le défi est l'intégration nationale par les infrastructures (route, rail) et la transformation locale pour réduire l'exportation brute. »
\end{exoresolu}

\begin{methode}[Analyser une filière agricole (Approche « Filière » : Acteurs, Flux, Valeur ajoutée, Gouvernance)]
\textbf{Objectif :} Décrire le parcours du produit du champ à l'assiette (amont $\to$ aval) pour identifier les goulets d'étranglement et les leviers de compétitivité.
\begin{enumerate}
    \item \textbf{Définir le périmètre :} Quelle spéculation ? (Cacao, Coton, Banane, Maïs, Lait, Poisson). Quelle échelle ? (Nationale, Bassin de production).
    \item \textbf{Identifier les maillons (Acteurs + Opérations) :}
    \begin{enumerate}
        \item \textbf{Amont (Fournisseurs) :} Semenciers, Engrais, Matériel, Crédit, Recherche (IRAD), Vulgarisation.
        \item \textbf{Production (Exploitants) :} Types d'exploitations (familiales, agro-industrielles), Itinéraires techniques, Rendements.
        \item \textbf{Collecte / Première transformation :} Coopératives, Négociants/Collecteurs, Égrenage (coton), Fermentation/Séchage (cacao), Abattage (bétail/poisson).
        \item \textbf{Transformation industrielle / Conditionnement :} Usines (SODECOTON, CDC, SOSUCAM, Chococam, Laiteries, Conserveries).
        \item \textbf{Distribution / Commercialisation :} Grossistes, Détaillants, Supermarchés, Export (Port Douala/Kribi), Marchés informels.
        \item \textbf{Consommation finale :} Ménages, Industries agro-alimentaires, Restauration.
    \end{enumerate}
    \item \textbf{Analyser les flux :} Physiques (tonnes), Financiers (FCFA), Information (prix, normes).
    \item \textbf{Évaluer la répartition de la valeur ajoutée :} Part Producteur / Part Transformateur / Part Distributeur / État (Taxes). \emph{Idéal : Part producteur > 40-50\% prix final.}
    \item \textbf{Diagnostiquer la gouvernance :} Interprofession (ex: CICC pour cacao/café, AIC pour coton), Contrats, Normes qualité (HACCP, Bio, Equitable), Régulation prix.
\end{enumerate}
\textbf{Sortie type Bac :} Schéma de la filière (Flèches) + Tableau SWOT de la filière + Proposition de maillon à renforcer.
\end{methode}

\begin{exoresolu}[Analyse de la filière Cacao au Cameroun]
\textbf{Énoncé :} À l'aide d'un schéma commenté, décrivez l'organisation de la filière cacao au Cameroun. Identifiez le maillon faible et proposez une mesure de renforcement.
\tcblower
\textbf{Solution détaillée :}
\textbf{1. Schéma de la filière (Description textuelle structurée) :}
\begin{center}
\begin{tikzpicture}[node distance=1.2cm and 0.5cm, every node/.style={font=\small, align=center}]
% Nodes
\node[draw, rectangle, fill=popgreenL, rounded corners, align=center] (amont){AMONT\\Semences (IRAD/SODECAO)\\Engrais/Phyto\\Crédit/Vulgarisation};
\node[draw, rectangle, fill=popgold!18, rounded corners, right=of amont, align=center] (prod){PRODUCTION\\400-600k exploitants familiaux\\~0.5-1 ha/moyen\\Rendement ~300-500 kg/ha\\Récolte, Écabossage,\\Fermentation, Séchage};
\node[draw, rectangle, fill=poporangeL, rounded corners, right=of prod, align=center] (collecte){COLLECTE / PREMIÈRE TRANSF.\\Négociants / Pisteurs\\Coopératives (COOPEC, etc.)\\Contrôle qualité (humidité, moisissures)\\Achats fèves séchées};
\node[draw, rectangle, fill=popblueL, rounded corners, right=of collecte, align=center] (transfo){TRANSFORMATION INDUSTRIELLE\\Torréfaction / Broyage\\Usines : Chococam, Sic Cacaos,\\Neo Industry, Atlantic Cocoa\\Produits : Liqueur, Beurre, Poudre};
\node[draw, rectangle, fill=poppinkL, rounded corners, right=of transfo, align=center] (distrib){DISTRIBUTION / EXPORT\\Exportateurs (Olam, Cargill, Telcar,\\Sociétés locales)\\Port Douala / Kribi\\Marché local (Chocolateries,\\Boulangeries)};
\node[draw, rectangle, fill=poppurpleL, rounded corners, right=of distrib, align=center] (conso){CONSOMMATION\\Industrie chocolatière mondiale\\Ménages camerounais\\Restauration};
% Gouvernance
\node[draw, ellipse, fill=popgrayL, above=of prod, xshift=1cm, align=center] (gov){GOUVERNANCE : CICC (Comité Interprof. Cacao-Café)\\ONCC (Office National) - Prix indicateur, Qualité, Promotion};
% Arrows
\draw[->, thick, popdark] (amont) -- (prod);
\draw[->, thick, popdark] (prod) -- (collecte);
\draw[->, thick, popdark] (collecte) -- (transfo);
\draw[->, thick, popdark] (transfo) -- (distrib);
\draw[->, thick, popdark] (distrib) -- (conso);
\draw[->, dashed, poppurple] (gov) -- (prod);
\draw[->, dashed, poppurple] (gov) -- (collecte);
\draw[->, dashed, poppurple] (gov) -- (transfo);
\end{tikzpicture}
\end{center}
\textbf{2. Analyse des maillons (Forces/Faiblesses) :}
\begin{itemize}
    \item \textbf{Production (Maillon FAIBLE) :} Vieillissement vergers ($>$ 40 ans), variétés non améliorées, bas rendement (300 kg/ha vs 1,5 t/ha potentiel), mauvaises pratiques post-récolte (fermentation insuffisante $\to$ fèves acides/moisies), faible encadrement, vieillissement paysans, difficulté d'accès au crédit/foncier.
    \item \textbf{Collecte :} Multiplicité d'intermédiaires (pisteurs) $\to$ perte de traçabilité, pression sur prix producteur, qualité hétérogène.
    \item \textbf{Transformation :} Capacité installée sous-utilisée (~40-50\%), dépendance export fèves brutes (70-80\% de la prod), faible valeur ajoutée locale.
    \item \textbf{Gouvernance (CICC) :} Rôle clé (prix plancher, kit intrants, rénovation vergers), mais moyens limités face à l'informel.
\end{itemize}
\textbf{3. Maillon faible identifié :} \textbf{La Production (Amont)}. C'est le goulot d'étranglement quantitatif (volume) et qualitatif (qualité fèves).
\textbf{4. Mesure de renforcement prioritaire :}
\textbf{Programme massif de rénovation/réhabilitation des vergers} couplé à la \textbf{diffusion de plants hybrides précoces et productifs (IRAD : « IRAD 1, 2, 3 » ou « Big 5 »)} via les coopératives, avec \textbf{subvention ciblée intrants (engrais, phyto)} et \textbf{formation obligatoire aux bonnes pratiques post-récolte} (fermentation 6j, séchage sur claies solaires). Condition : Sécurisation foncière pour inciter l'investissement long terme (25-30 ans).
\textbf{Conclusion :} La compétitivité de la filière cacao camerounaise (4ème/5ème mondial) repose sur la levée du verrou de la productivité à l'hectare et de la qualité marchande à la ferme.
\end{exoresolu}

\begin{attention}[Les 5 erreurs qui coûtent des points au Bac (Géographie - Activités Agropastorales)]
\begin{enumerate}
    \item \textbf{Confondre « Rendement » et « Production » :} Écrire « La production a augmenté » alors que seul le rendement a augmenté (superficie stable) ou inversement. \textbf{Réflexe :} Toujours préciser : Production = Superficie $\times$ Rendement. Citer les deux chiffres distincts.
    \item \textbf{Négliger la localisation spatiale (Le « Où » absent) :} Répondre « L'élevage est extensif » sans dire \textbf{OÙ} (Adamaoua, Nord, Extrême-Nord, couloirs de transhumance). \textbf{Réflexe :} Toute affirmation géographique doit être ancrée dans l'espace (Région, Bassin, Zone agro-écologique).
    \item \textbf{Lister les problèmes sans les hiérarchiser ni les relier aux causes structurelles :} Faire une « liste de courses » (manque d'argent, mauvaises routes, sécheresse) sans distinguer \textbf{contraintes naturelles} (climat, sol) et \textbf{contraintes humaines/économiques} (foncier, crédit, organisation, gouvernance). \textbf{Réflexe :} Grouper en catégories (Naturel / Structurel / Conjoncturel) et utiliser des connecteurs logiques (\emph{Par conséquent, Ce qui s'explique par...}).
    \item \textbf{Ignorer la dimension « Pêche » ou la traiter à part sans lien avec l'agropastoral :} Oublier la pêche continentale (Logone, Lac Tchad, Barrages) et maritime (crevettes, poissons pélagiques) ou ne pas mentionner les conflits d'usage (Eau : irrigation vs pêche ; Espace : agriculture vs pâturage). \textbf{Réflexe :} Intégrer la pêche dans le système agropastoral global (protéines, revenu, engrais organique).
    \item \textbf{Croquis / Carte : Légende incomplète, figurés inadaptés, absence d'échelle/orientation :} Utiliser des points pour des surfaces (cultures), des traits pour des points (villes), pas de flèche Nord, pas de titre. \textbf{Réflexe :} Respecter la sémiologie graphique (Surface = Hachures/Points ; Linéaire = Traits ; Ponctuel = Signes/Points proportionnels). Toujours : Titre, Légende hiérarchisée, Échelle, Nord.
\end{enumerate}
\end{attention}

\begin{exercice}[Entraînement Bac - Sujet 1]
À partir de la carte de répartition des activités agropastorales (non fournie ici, se reporter au croquis de synthèse), répondre aux questions :
\begin{enumerate}
    \item Localiser et nommer les trois principaux bassins de production cotonnière.
    \item Expliquer pourquoi l'élevage bovin moderne (ranch) se localise préférentiellement dans l'Adamaoua.
    \item Citer deux flux de commercialisation majeurs reliant le Nord au Sud et justifier leur sens.
    \item Calculer le taux d'autosuffisance en riz si la production locale est de 300\,000 t et la consommation nationale de 650\,000 t.
\end{enumerate}
\end{exercice}

\begin{corrige}[Exercice 1]
\begin{enumerate}
    \item \textbf{Bassins cotonniers :} Bassin de la Bénoué (Garoua, Mayo-Rey), Bassin du Diamaré (Maroua), Bassin du Mayo-Danay (Yagoua, Kaélé) -- Zone SODECOTON.
    \item \textbf{Adamaoua :} Altitude (climat tempéré, absence trypanosomose), herbacées palatables (Andropogon, Hyparrhenia), disponibilité foncière (peu dense), présence ranchs historiques (Ngaoundéré, Bamyanga), proximité marché urbain (Ngaoundéré, Garoua) et axes routiers/ferroviaires.
    \item \textbf{Flux 1 :} Bétail sur pied (Adamaoua/Nord $\to$ Yaoundé/Douala) : Demande urbaine protéines animales, pouvoir d'achat supérieur. \textbf{Flux 2 :} Coton fibre (Nord $\to$ Douala/Kribi) : Exportation vers ports maritimes (usines égrenage Garoua/Kaélé $\to$ rail/route $\to$ port).
    \item \textbf{Taux autosuffisance} = $\frac{300\,000}{650\,000} \times 100 \approx \textbf{46,2\%}$. Déficit $\approx$ 53,8\% comblé par importations.
\end{enumerate}
\end{corrige}

\begin{exercice}[Entraînement Bac - Sujet 2 : Analyse de document]
\textbf{Document :} Extrait rapport MINEPIA 2023. « La production halieutique nationale stagne autour de 250\,000 t/an alors que la demande dépasse 500\,000 t. La pêche artisanale maritime voit son CPUE passer de 250 kg/jour (2015) à 220 kg/jour (2023). Le nombre de pirogues motorisées a doublé en 10 ans. Les mangroves de l'estuaire du Wouri ont régressé de 30\% depuis 1990. »
\begin{enumerate}
    \item Calculer le déficit halieutique national et le taux de couverture des besoins par la production locale.
    \item Interpréter l'évolution du CPUE artisanale. Quel terme technique désigne cette situation ?
    \item Expliquer le lien entre régression des mangroves et baisse des captures.
    \item Proposer deux mesures concrètes pour inverser la tendance.
\end{enumerate}
\end{exercice}

\begin{corrige}[Exercice 2]
\begin{enumerate}
    \item \textbf{Déficit} = $500\,000 - 250\,000 = 250\,000$ t. \textbf{Taux de couverture} = $\frac{250\,000}{500\,000} \times 100 = \textbf{50\%}$.
    \item \textbf{Évolution :} CPUE passe de 250 à 220 kg/j (-12\% en 8 ans). Baisse continue = indicateur de \textbf{surpêche} (surpêche de croissance : capture juvéniles avant reproduction ; ou surpêche de recrutement). Effort de pêche (nb pirogues $\times$ motorisation) $>$ capacité de régénération du stock.
    \item \textbf{Lien mangroves/pêche :} Mangroves = \textbf{nurseries} (zones de reproduction, frayères, nourriceries) pour crevettes, poissons juvéniles (bars, capitaines, ethmaloses). Destruction (bois de chauffe, urbanisation, pollution, hydrocarbures) $\to$ perte habitat juvéniles $\to$ effondrement recrutement $\to$ baisse captures futures.
    \item \textbf{Mesures :} 1) \textbf{Restauration mangroves} (reboisement Rhizophora, protection légale zones sensibles, lutte pollution Wouri). 2) \textbf{Gestion effort pêche} : Gel licences nouvelles, rachat pirogues excédentaires, contrôle strict mailles filets, respect repos biologique (juil-sept), promotion aquaculture (tilapia/clarias) pour soulager pression milieu naturel.
\end{enumerate}
\end{corrige}

\begin{savaistu}[Le barrage de Lom Pangar : Clé de voûte de l'agropastoralisme ?]
Mis en service 2017 (stockage 6 milliards m³, puissance 30 MW). \textbf{Rôles multiples :}
\begin{itemize}
    \item \textbf{Énergie :} Régularisation débit Sanaga $\to$ optimisation barrages aval (Song Loulou, Edea) $\to$ électricité industries/urbain.
    \item \textbf{Agriculture :} Potentiel irrigation \textbf{12\,000 ha} aval (riz, maïs, légumes) dans vallée Sanaga (Dizangué, Edea). Projets \textbf{PDAR} (Projet Développement Agricole Rural).
    \item \textbf{Pêche :} Lac de retenue $\to$ pêche artisanale intense (Tilapia, Clarias), revenus riverains.
    \item \textbf{Élevage :} Abreuvement cheptel transhumant saison sèche, cultures fourragères irriguées.
    \item \textbf{Navigation :} Régime hydraulique stabilisé pour transport fluvial (bois, carburant, personnes).
\end{itemize}
\textbf{Enjeu :} Gestion intégrée ressources (Eau-Énergie-Alimentation-Environnement) = Approche \cle{Nexus}.
\end{savaistu}

\newpage

\coursec{Exercices}

\courssub{Je vérifie mes connaissances}

\begin{exercice}[QCM : les activités agropastorales]
Pour chaque question, entoure la (ou les) bonne(s) réponse(s).
\begin{enumerate}
\item L'agriculture vivrière a pour but principal :
\begin{enumerate}
\item[a)] l'exportation ; \quad b)] la nourriture de la famille du producteur ; \quad c)] la spéculation.
\end{enumerate}
\item Le cacao camerounais est principalement cultivé dans :
\begin{enumerate}
\item[a)] l'Adamaoua ; \quad b)] le Centre et le Sud ; \quad c)] l'Extrême-Nord.
\end{enumerate}
\item La principale zone d'élevage bovin du Cameroun est :
\begin{enumerate}
\item[a)] le Littoral ; \quad b)] l'Adamaoua ; \quad c)] le Sud-Ouest.
\end{enumerate}
\item La pêche continentale se pratique :
\begin{enumerate}
\item[a)] en haute mer ; \quad b)] sur les fleuves et les lacs ; \quad c)] dans les étangs piscicoles.
\end{enumerate}
\end{enumerate}
\end{exercice}

\begin{exercice}[Vrai ou faux ? Justifie ta réponse]
Réponds par VRAI ou FAUX, puis justifie en une ou deux phrases.
\begin{enumerate}
\item Le Cameroun exporte la quasi-totalité de sa production de poisson.
\item La transhumance est le déplacement saisonnier des troupeaux à la recherche de pâturages et d'eau.
\item La pisciculture est une forme de pêche maritime.
\item Les cultures de rente (cacao, café, coton) rapportent des devises au Cameroun.
\end{enumerate}
\end{exercice}

\begin{exercice}[Définitions]
Définis en deux ou trois lignes chacun des termes suivants : \cle{agriculture vivrière}, \cle{culture de rente}, \cle{transhumance}, \cle{pisciculture}.
\end{exercice}

\begin{exercice}[Calculs directs]
Le tableau ci-dessous donne les captures de poisson du Cameroun (en milliers de tonnes) pour une année donnée.
\begin{center}
\begin{tabularx}{\linewidth}{|l|X|X|X|X|}
\hline
\rowcolor{popblueL} Origine & Pêche maritime & Pêche continentale & Pisciculture & Importations \\
\hline
Quantité (milliers de t) & 120 & 60 & 5 & 200 \\
\hline
\end{tabularx}
\end{center}
\begin{enumerate}
\item Calcule la production nationale totale de poisson.
\item Calcule la part (en \%) de la pêche maritime dans la production nationale.
\item Calcule le taux de couverture des besoins par la production nationale (production nationale / consommation totale $\times$ 100).
\end{enumerate}
\end{exercice}

\courssub{Je m'entraîne}

\begin{exercice}[Les formes d'agriculture au Cameroun]
\begin{enumerate}
\item Cite les deux grandes formes d'agriculture pratiquées au Cameroun.
\item Pour chacune, donne deux caractéristiques et deux exemples de cultures.
\item Présente en quelques lignes deux limites de l'agriculture vivrière camerounaise.
\end{enumerate}
\end{exercice}

\begin{exercice}[L'élevage dans l'Adamaoua]
\begin{enumerate}
\item Explique deux raisons (climatiques et humaines) qui font de l'Adamaoua la principale zone d'élevage bovin du Cameroun.
\item Dégage deux problèmes liés à la transhumance.
\item Propose une solution pour moderniser l'élevage au Cameroun.
\end{enumerate}
\end{exercice}

\begin{exercice}[Construction graphique : la pêche au Cameroun]
À partir du tableau de l'exercice « Calculs directs » (pêche maritime : 120 ; pêche continentale : 60 ; pisciculture : 5 ; importations : 200, en milliers de tonnes) :
\begin{enumerate}
\item Construis un histogramme en bâtons représentant ces quatre quantités (échelle conseillée : 1 cm pour 20 milliers de tonnes).
\item Compare la production nationale aux importations.
\item Rédige un commentaire de quatre à six lignes sur la dépendance du Cameroun aux importations de poisson.
\end{enumerate}
\end{exercice}

\begin{exercice}[Les cultures d'exportation]
Le tableau ci-dessous donne les principales cultures de rente du Cameroun et leurs zones de production.
\begin{center}
\begin{tabularx}{\linewidth}{|l|X|}
\hline
\rowcolor{popblueL} Culture & Principales zones de production \\
\hline
Cacao & Centre, Sud, Littoral \\
\hline
Café & Ouest, Nord-Ouest, Adamaoua \\
\hline
Coton & Nord, Extrême-Nord \\
\hline
Banane & Littoral, Sud-Ouest \\
\hline
\end{tabularx}
\end{center}
\begin{enumerate}
\item Décris la répartition spatiale de ces cultures : quelles régions dominent ?
\item Explique en quelques lignes l'importance des cultures de rente pour l'économie camerounaise (devises, emplois, recettes de l'État).
\end{enumerate}
\end{exercice}

\courssub{Je me prépare au Bac}

\begin{exercice}[Dissertation type Probatoire]
Sujet : « Après avoir présenté les principales formes d'agriculture pratiquées au Cameroun, montre leurs limites et propose des solutions pour moderniser le secteur agricole. »

Tu rédigeras un devoir organisé comprenant :
\begin{enumerate}
\item une introduction (amorce, définition des termes, annonce du plan) ;
\item un développement en deux ou trois parties clairement délimitées ;
\item une conclusion avec ouverture.
\end{enumerate}
\end{exercice}

\begin{exercice}[Étude de documents : cultures de rente et exportations]
\textbf{Document 1.} Carte (fournie par l'enseignant) : localisation des zones de cacao (Centre-Sud), café (Ouest), coton (Nord) et banane (Littoral).

\textbf{Document 2.} Exportations agricoles du Cameroun (données simplifiées, en milliards de FCFA) :
\begin{center}
\begin{tabularx}{\linewidth}{|l|X|X|X|X|}
\hline
\rowcolor{popblueL} Produit & Cacao & Café & Coton & Banane \\
\hline
Valeur des exportations & 400 & 120 & 150 & 80 \\
\hline
\end{tabularx}
\end{center}
\begin{enumerate}
\item À partir du document 1, décris la répartition spatiale des cultures d'exportation au Cameroun.
\item Calcule la part du cacao (en \%) dans la valeur totale des quatre exportations du document 2.
\item Explique l'importance de ces cultures pour l'économie camerounaise (devises, emplois, équilibre de la balance commerciale).
\item Dégages deux risques liés à la dépendance du Cameroun aux cultures d'exportation.
\end{enumerate}
\end{exercice}

\begin{exercice}[Sujet type Baccalauréat technique : élevage, pêche et paradoxe des importations]
« Le Cameroun dispose d'atouts considérables pour l'élevage et la pêche, pourtant il importe viande et poisson. »
\begin{enumerate}
\item Présente deux atouts naturels et deux atouts humains du Cameroun pour l'élevage.
\item Présente deux atouts du Cameroun pour la pêche (maritime et continentale).
\item Explique deux causes qui obligent le Cameroun à importer de la viande et du poisson malgré ces atouts.
\item Propose trois pistes concrètes de relance du secteur (infrastructures, formation, pisciculture, transformation locale).
\item Rédige une conclusion de cinq lignes répondant au paradoxe énoncé dans le sujet.
\end{enumerate}
\end{exercice}

\newpage

\coursec{Corrigés des exercices}

\begin{corrige}[Exercice 1 — QCM : les activités agropastorales]
\begin{enumerate}
\item \textbf{Bonne réponse : b)}. L'agriculture vivrière a pour but principal de nourrir la famille du producteur ; seul l'éventuel surplus est vendu sur les marchés locaux.
\item \textbf{Bonne réponse : b)}. Le cacao est une culture de rente qui exige chaleur et humidité : il est cultivé dans la zone forestière humide, notamment dans les régions du Centre, du Sud et du Littoral.
\item \textbf{Bonne réponse : b)}. L'Adamaoua, avec ses hauts plateaux, son climat soudanien et ses vastes pâturages, est la principale zone d'élevage bovin du Cameroun.
\item \textbf{Bonne réponse : b)}. La pêche continentale se pratique sur les eaux intérieures : fleuves (Sanaga, Logone), rivières et lacs (lac Tchad, lagunes). La pêche en haute mer relève de la pêche maritime, et les étangs relèvent de la pisciculture.
\end{enumerate}
\end{corrige}

\begin{corrige}[Exercice 2 — Vrai ou faux ?]
\begin{enumerate}
\item \textbf{FAUX.} Le Cameroun consomme sa production nationale de poisson et doit même importer massivement (environ 200\,000 tonnes par an) pour couvrir ses besoins ; il n'exporte pas sa production.
\item \textbf{VRAI.} La transhumance est bien le déplacement saisonnier des troupeaux, menés par les éleveurs, à la recherche de pâturages et de points d'eau, notamment entre la saison sèche et la saison des pluies.
\item \textbf{FAUX.} La pisciculture est l'élevage de poissons en étangs ou en bassins ; ce n'est pas une pêche maritime, qui se pratique en mer.
\item \textbf{VRAI.} Le cacao, le café et le coton sont vendus à l'étranger et rapportent des devises (monnaies étrangères) qui alimentent le budget de l'État et la balance commerciale.
\end{enumerate}
\end{corrige}

\begin{corrige}[Exercice 3 — Définitions]
\begin{itemize}
\item \textbf{Agriculture vivrière} : forme d'agriculture dont la production est destinée avant tout à la consommation de la famille du producteur (manioc, maïs, igname, plantain), sur de petites exploitations, avec des outils simples ; seul le surplus est vendu.
\item \textbf{Culture de rente} : culture destinée à la vente, surtout à l'exportation (cacao, café, coton, banane), qui procure un revenu monétaire au producteur et des devises à l'État.
\item \textbf{Transhumance} : déplacement saisonnier des troupeaux et de leurs éleveurs à la recherche de pâturages et d'eau, entre les zones de saison sèche et les zones de saison des pluies.
\item \textbf{Pisciculture} : élevage de poissons dans des étangs, des bassins ou des enclos aménagés, afin de produire du poisson de manière contrôlée.
\end{itemize}
\end{corrige}

\begin{corrige}[Exercice 4 — Calculs directs]
\begin{enumerate}
\item La production nationale totale est la somme des trois sources nationales :
\[ 120 + 60 + 5 = 185 \text{ milliers de tonnes.} \]
\item Part de la pêche maritime dans la production nationale :
\[ \frac{120}{185} \times 100 \approx 64,9\,\%. \]
La pêche maritime fournit donc environ \textbf{64,9\,\%} de la production nationale.
\item La consommation totale est la production nationale augmentée des importations :
\[ 185 + 200 = 385 \text{ milliers de tonnes.} \]
Le taux de couverture des besoins est :
\[ \frac{185}{385} \times 100 \approx 48,1\,\%. \]
La production nationale ne couvre donc qu'environ \textbf{48\,\%} des besoins : plus de la moitié du poisson consommé au Cameroun est importée.
\end{enumerate}
\textbf{Point de vigilance :} ne pas inclure les importations dans la production nationale ; elles n'interviennent que dans le calcul de la consommation totale.
\end{corrige}

\begin{corrige}[Exercice 5 — Les formes d'agriculture au Cameroun]
\begin{enumerate}
\item Les deux grandes formes d'agriculture sont l'\cle{agriculture vivrière} (traditionnelle) et l'\cle{agriculture de rente} (ou de plantation, commerciale).
\item \textbf{Agriculture vivrière} : petites exploitations familiales, outils simples (machette, houe), faible rendement, production destinée à l'autoconsommation. Exemples de cultures : manioc, maïs (aussi igname, plantain, mil).
\textbf{Agriculture de rente} : production destinée à la vente et à l'exportation, exploitations plus grandes (plantations), besoin de capitaux et de main-d'œuvre salariée. Exemples : cacao, café (aussi coton, banane, palmier à huile).
\item \textbf{Limites de l'agriculture vivrière} : (a) les rendements sont faibles à cause des outils rudimentaires et du peu d'intrants (engrais, semences améliorées) ; (b) les surfaces sont petites et morcelées, ce qui empêche la mécanisation ; on peut ajouter la faible commercialisation, qui maintient les paysans dans la pauvreté.
\end{enumerate}
\end{corrige}

\begin{corrige}[Exercice 6 — L'élevage dans l'Adamaoua]
\begin{enumerate}
\item \textbf{Raisons climatiques} : le climat soudanien de l'Adamaoua, avec une saison sèche marquée, limite la mouche tsé-tsé (qui transmet la trypanosomiase) ; les hautes savanes offrent de vastes pâturages herbeux. \textbf{Raisons humaines} : la présence d'éleveurs peuls (Foulbé) qui possèdent une longue tradition pastorale et un savoir-faire reconnu ; la faible densité de population laisse de grands espaces disponibles pour les troupeaux.
\item \textbf{Problèmes liés à la transhumance} : les conflits entre agriculteurs et éleveurs (destruction des champs par les troupeaux) ; la dégradation des pâturages et des points d'eau par surpâturage ; on peut citer aussi les vols de bétail et les difficultés de suivi sanitaire des animaux.
\item \textbf{Solution de modernisation} : créer des ranchs modernes avec pâturages améliorés, points d'eau (forages) et suivi vétérinaire ; on peut aussi proposer la sédentarisation des éleveurs, la sélection des races et le développement des unités de transformation (laiteries, abattoirs frigorifiques).
\end{enumerate}
\end{corrige}

\begin{corrige}[Exercice 7 — Construction graphique : la pêche au Cameroun]
\begin{enumerate}
\item Avec l'échelle 1 cm pour 20 milliers de tonnes, les hauteurs des bâtons sont : pêche maritime : $120/20 = 6$ cm ; pêche continentale : $60/20 = 3$ cm ; pisciculture : $5/20 = 0,25$ cm ; importations : $200/20 = 10$ cm. Le graphique attendu :
\begin{popfigure}
\begin{tikzpicture}
\begin{axis}[
  ybar, bar width=0.55cm,
  width=\linewidth, height=6.2cm,
  ymin=0, ymax=220,
  ylabel={Milliers de tonnes},
  symbolic x coords={Maritime,Continentale,Pisciculture,Importations},
  xtick=data,
  x tick label style={font=\small, rotate=12, anchor=east},
  ytick={0,50,100,150,200},
  nodes near coords, nodes near coords style={font=\small},
  axis lines=left,
]
\addplot[fill=popblue!70, draw=popblue] coordinates {(Maritime,120) (Continentale,60) (Pisciculture,5) (Importations,200)};
\end{axis}
\end{tikzpicture}
\legende{Histogramme des quantités de poisson au Cameroun (en milliers de tonnes) : production nationale (pêche maritime, pêche continentale, pisciculture) et importations.}
\end{popfigure}
\item La production nationale totale est de $120 + 60 + 5 = 185$ milliers de tonnes, contre $200$ milliers de tonnes d'importations : les importations dépassent la production nationale ($200 > 185$).
\item \textbf{Commentaire attendu} : le Cameroun produit environ 185\,000 tonnes de poisson, dominées par la pêche maritime (environ 65\,\%), tandis que la pisciculture reste marginale (5\,000 tonnes). Or le pays importe 200\,000 tonnes, soit plus que sa propre production. Le taux de couverture des besoins n'est que d'environ 48\,\% : plus de la moitié du poisson consommé vient de l'étranger. Cette dépendance coûte cher en devises et fragilise la sécurité alimentaire. Le développement de la pisciculture et la modernisation des pêcheries apparaissent donc comme des priorités.
\end{enumerate}
\end{corrige}

\begin{corrige}[Exercice 8 — Les cultures d'exportation]
\begin{enumerate}
\item La répartition spatiale suit les grandes zones agro-écologiques : le \textbf{cacao} et la \textbf{banane} occupent le sud forestier et humide (Centre, Sud, Littoral, Sud-Ouest) ; le \textbf{café} domine sur les hautes terres de l'Ouest et du Nord-Ouest ; le \textbf{coton} est la culture des savanes sèches du Nord et de l'Extrême-Nord. Les régions du Sud forestier concentrent donc l'essentiel des cultures de rente, le Nord étant marqué par le coton.
\item \textbf{Importance économique} : ces cultures rapportent des \textbf{devises} grâce aux exportations (le cacao est le premier produit agricole exporté) ; elles fournissent des \textbf{emplois} à des centaines de milliers de paysans et de salariés des plantations ; elles alimentent les \textbf{recettes de l'État} (taxes, droits d'exportation) et contribuent à l'équilibre de la \textbf{balance commerciale}. Elles font vivre aussi des filières de transformation (chocolaterie, huileries, filature du coton).
\end{enumerate}
\end{corrige}

\begin{corrige}[Exercice 9 — Dissertation type Probatoire]
\textbf{Barème indicatif (sur 20)} : introduction 3 pts ; développement 12 pts (4 pts par partie) ; conclusion 3 pts ; présentation et orthographe 2 pts.

\textbf{Introduction attendue.} Amorce : l'agriculture occupe la majorité de la population active camerounaise et nourrit le pays. Définitions : agriculture vivrière (production pour l'autoconsommation) et agriculture de rente (production pour la vente et l'exportation). Problématique : quelles sont les formes d'agriculture au Cameroun, quelles limites freinent-elles le développement, et comment moderniser le secteur ? Annonce du plan : nous présenterons d'abord les formes d'agriculture, puis leurs limites, enfin les solutions de modernisation.

\textbf{Développement.}
\emph{Première partie : les formes d'agriculture.} L'agriculture vivrière traditionnelle : petites exploitations, outils simples, cultures de manioc, maïs, plantain, igname ; elle occupe la majorité des paysans. L'agriculture de rente : plantations de cacao (Centre-Sud), café (Ouest), coton (Nord), banane (Littoral) ; elle est tournée vers l'exportation et rapporte des devises.
\emph{Deuxième partie : les limites.} Pour l'agriculture vivrière : faibles rendements, outillage rudimentaire, manque d'engrais et de semences améliorées, morcellement des terres, difficultés d'écoulement (mauvaises pistes rurales). Pour l'agriculture de rente : baisse des cours mondiaux, vieillissement des plantations et des planteurs, faible transformation locale, concurrence internationale.
\emph{Troisième partie : les solutions.} Mécanisation et vulgarisation agricole (encadrement des paysans) ; distribution d'intrants et de semences améliorées ; construction de pistes rurales et d'infrastructures de stockage ; développement de la transformation locale (usines) ; crédit agricole et organisation des paysans en coopératives ; formation des jeunes aux métiers agricoles.

\textbf{Conclusion.} Bilan : le Cameroun pratique une agriculture duale, vivrière et de rente, toutes deux freinées par de nombreuses contraintes. Ouverture : la modernisation de l'agriculture est une condition de la sécurité alimentaire et de l'émergence économique du pays.

\textbf{Points de vigilance :} définir les termes du sujet dans l'introduction ; annoncer clairement le plan ; illustrer chaque partie par des exemples camerounais précis (cultures, régions) ; éviter le hors-sujet sur l'élevage ou la pêche.
\end{corrige}

\begin{corrige}[Exercice 10 — Étude de documents : cultures de rente et exportations]
\begin{enumerate}
\item \textbf{Répartition spatiale (document 1)} : le cacao occupe le Centre et le Sud (zone forestière humide) ; le café est cultivé sur les hautes terres de l'Ouest ; le coton domine dans les savanes du Nord ; la banane se concentre sur le Littoral. Chaque culture correspond donc à une zone agro-climatique précise, du sud humide au nord sec.
\item \textbf{Part du cacao.} Valeur totale des quatre exportations :
\[ 400 + 120 + 150 + 80 = 750 \text{ milliards de FCFA.} \]
Part du cacao :
\[ \frac{400}{750} \times 100 \approx 53,3\,\%. \]
Le cacao représente donc environ \textbf{53,3\,\%} de la valeur de ces quatre exportations : c'est la première culture d'exportation.
\item \textbf{Importance économique} : ces cultures rapportent des devises indispensables au pays ; elles créent des emplois ruraux (planteurs, ouvriers de plantation, transporteurs, exportateurs) ; elles améliorent la balance commerciale en augmentant la valeur des exportations ; elles alimentent les recettes de l'État par les taxes.
\item \textbf{Deux risques de la dépendance} : (a) la \textbf{baisse des cours mondiaux} (les prix du cacao ou du café sont fixés à l'étranger et peuvent s'effondrer, réduisant brutalement les revenus) ; (b) la \textbf{faible transformation locale} : le Cameroun exporte des produits bruts peu valorisés et réimporte des produits finis chers. On peut citer aussi la vulnérabilité aux aléas climatiques et la concurrence d'autres pays producteurs.
\end{enumerate}
\textbf{Barème indicatif (sur 20)} : question 1 : 4 pts ; question 2 : 4 pts (calcul détaillé exigé) ; question 3 : 6 pts ; question 4 : 4 pts ; présentation : 2 pts.
\textbf{Points de vigilance :} le calcul de la part doit montrer le total (750) avant le pourcentage ; les risques doivent être expliqués, pas seulement cités.
\end{corrige}

\begin{corrige}[Exercice 11 — Sujet type Baccalauréat technique : élevage, pêche et paradoxe des importations]
\begin{enumerate}
\item \textbf{Atouts pour l'élevage.} Atouts naturels : de vastes pâturages de savane dans l'Adamaoua et le Nord ; un climat soudanien qui limite la mouche tsé-tsé ; des points d'eau (fleuves, mares). Atouts humains : une longue tradition pastorale des éleveurs peuls ; une main-d'œuvre disponible et un savoir-faire reconnu ; une forte demande de viande d'une population en croissance.
\item \textbf{Atouts pour la pêche.} Pêche maritime : une façade atlantique d'environ 400 km riche en poissons (côtes du Littoral et du Sud-Ouest). Pêche continentale : de nombreux fleuves et lacs (Sanaga, Logone, lac Tchad, lagunes) ; on peut ajouter un climat favorable à la pisciculture.
\item \textbf{Causes des importations.} (a) Pour l'élevage : techniques traditionnelles peu productives (transhumance, races à faible rendement), manque d'infrastructures (abattoirs, chaîne du froid) et de suivi vétérinaire. (b) Pour la pêche : pêche artisanale aux pirogues et filets rudimentaires, pisciculture embryonnaire (environ 5\,000 tonnes seulement), absence de congélation et de conservation, si bien que la production nationale (environ 185\,000 tonnes) ne couvre qu'environ 48\,\% des besoins, d'où environ 200\,000 tonnes de poisson importé.
\item \textbf{Pistes de relance} : construire des infrastructures (abattoirs frigorifiques, chambres froides, ports de pêche, forages pastoraux) ; former les éleveurs et les pêcheurs (vulgarisation, centres de formation) ; développer massivement la \textbf{pisciculture} (étangs, alevinage) ; encourager la \textbf{transformation locale} (fumage moderne, congélation, laiteries) ; faciliter le crédit et l'organisation en coopératives.
\item \textbf{Conclusion attendue (environ cinq lignes)} : le paradoxe s'explique par l'écart entre un potentiel naturel et humain remarquable et des techniques restées traditionnelles. Le manque d'infrastructures, de formation et de transformation empêche la production nationale de répondre à une demande croissante. Le Cameroun importe donc ce qu'il pourrait produire. En investissant dans la modernisation de l'élevage, des pêcheries et de la pisciculture, le pays peut réduire sa dépendance, économiser des devises et renforcer sa sécurité alimentaire.
\end{enumerate}
\textbf{Barème indicatif (sur 20)} : question 1 : 4 pts ; question 2 : 3 pts ; question 3 : 5 pts ; question 4 : 5 pts ; question 5 : 3 pts.
\textbf{Points de vigilance :} distinguer clairement atouts naturels et atouts humains ; appuyer l'argumentation sur les chiffres de la leçon (185\,000 t produites, 200\,000 t importées, taux de couverture d'environ 48\,\%) ; la conclusion doit répondre explicitement au paradoxe du sujet.
\end{corrige}

\newpage

\coursec{Fiche bilan}

\begin{popfigure}
\begin{tikzpicture}[mindmap, grow cyclic, every node/.style={concept, font=\small},
root concept/.append style={concept color=popblue, font=\small\bfseries},
level 1/.append style={level distance=3.2cm, sibling angle=60},
level 2/.append style={level distance=1.9cm, sibling angle=45, font=\scriptsize}]
\node[root concept]{Activités agropastorales et piscicoles}
child[concept color=popgreen]{node{Généralités}
  child{node{Agriculture, élevage, pêche}}
  child{node{Secteur primaire}}
}
child[concept color=poporange]{node{Agriculture}
  child{node{Vivrière : mil, manioc, maïs}}
  child{node{D'exportation : cacao, café, coton}}
}
child[concept color=poppurple]{node{Élevage}
  child{node{Extensif (Nord)}}
  child{node{Intensif (Ouest)}}
}
child[concept color=popblue!60]{node{Pêche}
  child{node{Maritime (Kribi, Douala)}}
  child{node{Continentale (lacs, fleuves)}}
  child{node{Pisciculture}}
}
child[concept color=poppink]{node{Atouts}
  child{node{Climats variés}}
  child{node{Main-d'œuvre}}
}
child[concept color=popgold]{node{Problèmes et perspectives}
  child{node{Faible mécanisation}}
  child{node{Modernisation, irrigation}}
};
\end{tikzpicture}
\legende{Carte mentale de la leçon : les six branches essentielles des activités agropastorales et piscicoles au Cameroun.}
\end{popfigure}

\begin{bilan}
\textbf{Définitions clés}
\begin{itemize}
\item \cle{Agriculture} : ensemble des activités de culture des végétaux pour nourrir les hommes et fournir des matières premières.
\item \cle{Agriculture vivrière} : agriculture destinée à l'autoconsommation et au marché local (manioc, maïs, mil, igname, plantain).
\item \cle{Agriculture de rente (d'exportation)} : cultures vendues à l'étranger (cacao, café, coton, banane, hévéa, palmier à huile).
\item \cle{Élevage extensif} : élevage traditionnel, troupeaux déplacés sur de grandes surfaces (transhumance), faible rendement.
\item \cle{Élevage intensif} : élevage moderne, animaux parqués et bien nourris, fort rendement (poulets, porcs).
\item \cle{Pêche continentale} : pêche pratiquée dans les eaux intérieures (fleuves, lacs) ; \cle{pêche maritime} : en mer ; \cle{pisciculture} : élevage de poissons en étangs.
\end{itemize}

\textbf{Repères chiffrés à connaître}
\begin{center}
\begin{tabularx}{\linewidth}{|l|X|X|}\hline
\rowcolor{popblueL} \textbf{Activité} & \textbf{Zones principales} & \textbf{Produits phares} \\ \hline
Agriculture vivrière & Partout, surtout Centre, Sud, Nord & Manioc, maïs, mil, plantain \\ \hline
Agriculture d'exportation & Littoral, Sud-Ouest, Centre & Cacao, café, coton, banane \\ \hline
Élevage & Adamaoua, Nord, Extrême-Nord & Bovins, ovins, caprins \\ \hline
Pêche & Côte atlantique, lac Tchad, Logone & Capitaine, carpe, tilapia \\ \hline
\end{tabularx}
\end{center}

\textbf{Méthodes express}
\begin{itemize}
\item Croquis : contour simplifié du Cameroun + zones colorées + villes (points nommés) + flèche du nord + légende numérotée.
\item Analyse d'un problème : constat chiffré $\rightarrow$ causes $\rightarrow$ conséquences $\rightarrow$ solutions.
\item Justifier une réponse : toujours appuyer l'argument par un exemple localisé et un chiffre.
\end{itemize}
\end{bilan}

\begin{aretenir}[Checklist avant le Bac]
\begin{itemize}
\item[$\square$] Je sais définir agriculture vivrière, agriculture de rente, élevage extensif et intensif.
\item[$\square$] Je distingue pêche maritime, pêche continentale et pisciculture.
\item[$\square$] Je cite au moins trois cultures d'exportation du Cameroun et leurs zones.
\item[$\square$] Je situe les grandes zones d'élevage (Adamaoua, Nord, Extrême-Nord).
\item[$\square$] Je connais les principales zones de pêche (côte atlantique, lac Tchad, Logone).
\item[$\square$] Je sais expliquer les atouts naturels du Cameroun (climats, sols, hydrographie).
\item[$\square$] Je cite au moins quatre problèmes du secteur : faible mécanisation, enclavement, aléas climatiques, faible transformation.
\item[$\square$] Je propose des solutions : irrigation, crédit agricole, routes rurales, formation.
\item[$\square$] Je sais construire un croquis avec titre, légende, orientation et échelle indicative.
\item[$\square$] Je rédige une réponse justifiée avec un exemple et un chiffre.
\end{itemize}
\end{aretenir}

\findecours

\end{document}
